% DER_XVIII-c: Grothendieck's text with the commentary set in, then the overview, inside a group that applies the chunk's own definitions
\dgrpart{1}{XVIII-c}{Part 3 of 5, manuscript pages 220--328}{XVIII-c\enspace\textperiodcentered\enspace Chapter XVIII, part 3 of 5\enspace\textperiodcentered\enspace manuscript pages 220--328}{Chapter XVIII \textperiodcentered\ \textup{XVIII-c}}
\dgrchunkbegin
% DER_XVIII-c: the definitions of the chunk's preamble, applied inside its group
\providecommand{\IfDocumentMetadataT}[1]{}
\let\dctauAB\relax\newsavebox{\dctauAB}
\let\dctauABone\relax\newsavebox{\dctauABone}
\let\dctauAAzero\relax\newsavebox{\dctauAAzero}
\let\printedpage\relax\newcommand{\printedpage}[1]{\marginpar{\footnotesize p.~#1}}
\let\msp\relax\newcommand{\msp}[1]{\marginpar{\footnotesize [#1]}}
\let\uHom\relax\newcommand{\uHom}{\underline{\mathrm{Hom}}}
\let\Ind\relax\newcommand{\Ind}{\mathrm{Ind}}
\let\Ob\relax\newcommand{\Ob}{\operatorname{Ob}}
\let\Fl\relax\newcommand{\Fl}{\operatorname{Fl}}
\let\card\relax\newcommand{\card}{\operatorname{card}}
\let\cardred\relax\newcommand{\cardred}{\mathop{\textup{card\,réd}}\nolimits}
\let\cM\relax\newcommand{\cM}{\mathcal{M}}
\let\rbox\relax\newcommand{\rbox}[1]{\tikz[baseline=(X.base)]\node[draw,rounded corners=4pt,inner sep=2pt](X){\textup{#1}};}
\let\bs\relax\newcommand{\bs}{\backslash}
\let\encadre\relax\newcommand{\encadre}[1]{\tikz[baseline=(X.base)]\node[draw,rounded corners,inner sep=2pt](X){#1};}
\let\uGamma\relax\newcommand{\uGamma}{\underline{\Gamma}}
\let\trois\relax\newcommand{\trois}[1]{\mathrel{\overset{#1}{\vcenter{\offinterlineskip\hbox{$\longrightarrow$}\vskip0.5pt\hbox{$\longrightarrow$}\vskip0.5pt\hbox{$\longrightarrow$}}}}}
\let\Hom\relax\newcommand{\Hom}{\operatorname{Hom}}
\let\cercle\relax\newcommand{\cercle}[1]{\tikz[baseline=(X.base)]\node[draw,circle,inner sep=1pt](X){\upshape #1};}
\let\defeq\relax\newcommand{\defeq}{\overset{\text{\upshape déf}}{=}}
\let\into\relax\newcommand{\into}{\lhook\joinrel\longrightarrow}
\let\M\relax\newcommand{\M}{\mathcal{M}}
\let\N\relax\newcommand{\N}{\mathcal{N}}
\let\cH\relax\newcommand{\cH}{\mathcal{H}}
\let\red\relax\newcommand{\red}{\mathop{\textup{réd}}\nolimits}
\let\uS\relax\newcommand{\uS}{\underline{S}}
\let\incl\relax\newcommand{\incl}{\lhook\joinrel\longrightarrow}
\let\ilim\relax\newcommand{\ilim}{\varinjlim\limits}
\let\circlabel\relax\newcommand{\circlabel}[1]{\tikz[baseline=(c.base)]\node[draw,circle,inner sep=1pt](c){#1};}
\let\circled\relax\newcommand{\circled}[1]{\tikz[baseline=(c.base)]\node[draw,circle,inner sep=1pt](c){#1};}
\let\ovalbox\relax\newcommand{\ovalbox}[1]{\tikz[baseline=(c.base)]\node[draw,rounded corners=6pt,inner sep=2pt](c){#1};}

\sbox{\dctauAB}{\circled{$\scriptstyle\tau_{\alpha,\beta}$}}
\sbox{\dctauABone}{\circled{$\scriptstyle\tau_{\alpha,\beta+1}$}}
\sbox{\dctauAAzero}{\circled{$\scriptstyle\tau_{\alpha,\alpha_0}$}}
% ---- from der-ch18_p081-085.tex
\printedpage{79}
% head: A. Grothendieck : Dérivateurs, Ch. XVIII, Édition des Universités de Montpellier II et Paris VII

Le résultat principal dans la présente section est le suivant :

\textbf{Théorème 4.9.3.} \textit{Sous les conditions générales de 4.9.1 pour $C$, $\pi_0$, $J$ (notamment} $\cardred J \leqslant \pi_0$\textit{), le foncteur pleinement fidèle (4.9.3) est une équivalence de catégories dans chacun des deux cas suivants :}
\begin{itemize}[label={}]
{\itshape
\item[a)] $C$ est stable par $\varinjlim$ filtrantes dénombrables.
\item[b)] Il existe un cardinal $\pi_0 < \pi$, et un cardinal $\pi'_0 > \pi_0$ \textup{(}\textbf{\upshape N.B.} On peut supposer que $\pi'_0$ est le successeur de $\pi_0$, donc $\pi'_0 \leqslant \pi$), tels que $C$ soit stable par limites inductives grandes devant $\pi_0$, et de cardinal $\pi'_0$.
}
\end{itemize}
(\footnote[94]{Cas commun : il existe un ordinal limite $\alpha$, tel que 1) $\card \alpha \leqslant \pi_0$. 2) Les limites inductives \texttt{[de]} type $I_\alpha = \{\beta \mid \beta < \alpha\}$ sont représentables dans $C$.}).

\textbf{Remarque 4.9.4.}
\begin{itemize}[label={}]
\item[a)] Ce théorème, cas a), contient le théorème 4.4.13, partie b) (p. 95-96), où on suppose que $C$ est stable par limites inductives de cardinal $\leqslant \pi_0$, l'hypothèse a) de 4.9.3 étant considérablement plus faible.
\item[b)] \msp{page 220}Je présume qu'en dehors d'une hypothèse restrictive du type a) ou b) ci-dessus, le foncteur pleinement fidèle (4.9.3) n'est pas toujours essentiellement surjectif, même dans le cas particulier où $J$ est une catégorie finie, p.ex. $J = B_G$, $G = \{\pm 1\}$. Pour un contre-exemple dans le cas très voisin où on prend $\cM = \Ind(C)$ (au lieu de $\Ind_{\pi_0}(C)$), et où on considère le foncteur (pleinement fidèle si $J$ est une catégorie équivalent à une catégorie \textit{finie}) % [sic] équivalent -> équivalente, agreement with « une catégorie »
\[
\Ind(\uHom(J,C)) \longrightarrow \uHom(J,\Ind(C)),
\]
je renvoie à l'exercice SGA 4 I 8.8.8.
\end{itemize}
L'application principale du théorème 4.9.3 consiste en le

\textbf{Corollaire 4.9.5.} \textit{Soient $\cM$ une catégorie} accessible \textit{(4.8.7), $J$ une catégorie équivalente à une petite catégorie. Alors la catégorie $\cM(J^\circ) = \uHom(J,\cM)$ est accessible. Plus précisément, soient $\pi_0$ un cardinal} utile \textit{pour $\cM$ (i.e. tel que $\Ind_{\pi_0}(\cM_{\pi_0}) \xrightarrow{\sim} \cM$, i.e. un cardinal tel que $\cM$ soit $\pi_0$-accessible), soit $\pi$ un cardinal infini tel que $\pi^{\pi_0} = \pi \geqslant \cardred J$} (\footnote[95]{On a donc $\pi = \pi^{\pi_0} \geqslant 2^{\pi_0} > \pi_0$.}), \textit{i.e. $\pi = c^{\pi_0}$, où $c$ est n'importe quel cardinal $c \geqslant \sup(\pi_0, \cardred J)$, alors $\uHom(J,\cM)$ est $\pi$-accessible, plus précisément on a une équivalence de catégories} % doubt resolved against PDF; see corrections.csv
\[
\uHom(J,\cM) \xleftarrow{\ \sim\ } \Ind_\pi(\uHom(J,\cM_\pi)) \ . \tag*{(4.9.5.1)}
\]

\msp{page 221}On sait en effet, avec les hypothèses faites sur $\pi_0$, $\pi$, que $\pi$ est utile pour $\cM$, i.e. que
\[
\cM \xleftarrow{\ \sim\ } \Ind_\pi(\cM_\pi) \ ,
\]
d'autre part, $\cM_\pi$ est stable par $\varinjlim$ grandes devant $\pi_0$ et de cardinal $\leqslant \pi$. On peut donc appliquer 4.9.3 b), en prenant $C = \cM_\pi$, $\pi$ et $\pi_0$ comme ici (donc $\pi_0 < 2^{\pi_0} \leqslant \pi$), et $\pi'_0 = \pi$.

\dgrblock{dgrb1042}{One indexing category for an entire diagram}
\begin{dcomment}{exposition}{p220}{One indexing category for an entire diagram}
Grothendieck wants an approximation of a diagram which remembers its arrows.  Writing each value of $\varphi:J\to\Ind_\pi(C)$ as an ind-object separately does not provide this: an arrow can be represented only after enlarging the chosen indices, and a relation between arrows can require a further enlargement.  The comparison functor therefore asks for one $\pi$-filtered system of genuine $J$-diagrams in $C$.  Its essential surjectivity is a coherence assertion about presentations, not merely the existence of presentations of their values.

We shall use a cardinal convention which removes a recurrent ambiguity in these pages.  For an infinite cardinal $\rho$, ``large before $\rho$'' means that every diagram with at most $\rho$ arrows has a cocone.  This is modern $\rho^+$-filteredness.  An object called $\rho$-accessible here is correspondingly $\rho^+$-presentable.  We retain $\Ind_\rho$ for Grothendieck's completion, and reserve $\lambda$ for a regular cardinal when using modern terminology.  Thus $\Ind_\rho(C)$ is the completion under $\lambda$-filtered colimits with $\lambda=\rho^+$.  This convention applies even when $\rho$ itself is singular.

The precise sufficient condition established by the ensuing proof is as follows.  Suppose $C$ is essentially small, $|\operatorname{Mor}J|\leq\rho$, and $C$ admits colimits of chains of one nonzero limit ordinal $\alpha$ with $|\alpha|\leq\rho$.  Then
\[
 \Ind_\rho([J,C])\longrightarrow[J,\Ind_\rho(C)]
\]
is an equivalence.  One may take $\alpha=\omega$ in case (a).  For the intended case (b), take a smaller infinite cardinal $\theta<\rho$ and $\alpha=\theta^+$, assuming the stipulated $\theta^+$-filtered colimits of size at most $\theta^+$ exist in $C$.  The printed statement moves between $\pi_0$ as the completion bound and $\pi_0$ as this smaller auxiliary cardinal.  Corollary 4.9.5 and the formulation in 4.10 make the intended roles clear; the proof must keep the bound on $|\alpha|$ tied to the completion cardinal.  Footnote 94 should be read with that same bound.  This is a normalization of the statement, not an additional closure assumption on the intended theorem.

Full faithfulness has a short explanation which also identifies the cardinal restriction on $J$.  For $A,B:J\to\mathcal M$,
\[
 \operatorname{Nat}(A,B)=
 \operatorname{Eq}\left(
 \prod_{j\in J}\operatorname{Hom}(A_j,B_j)
 \rightrightarrows
 \prod_{u:j\to k}\operatorname{Hom}(A_j,B_k)
 \right).
\]
If every $A_j$ is $\rho^+$-presentable, this expression commutes with $\rho^+$-filtered colimits: every product has at most $\rho$ factors, and the equality conditions have the same bound.  Consequently the objects of $[J,C]$ remain presentable in $[J,\Ind_\rho C]$, so extension under filtered colimits is fully faithful.  The long construction which follows is needed only to produce enough such objects.  Its antecedent is the indexed-presentation criterion of \citet[I, 8.8.2--8.8.3]{SGA4I1972}; the present cardinal accounting is what makes that criterion usable for accessible categories.
\end{dcomment}


\dgrblock{dgrb1043}{The involution obstruction at every cardinal}
\begin{dcomment}{examples}{p220}{The involution obstruction at every cardinal}
The reservation in Remark 4.9.4 has a concrete answer.  The obstruction in SGA~4, I, Exercise 8.8.8 is an involution of an ind-object which cannot be reached equivariantly from any of the chosen approximating objects.  The same construction works at every infinite completion bound; it is the case $I=B(\mathbf Z/2)$ of \citet[Example 1.4, Construction 3.1 and Section 3.2]{Henry2025Ind}.  We give the calculation because it distinguishes the theorem's actual hypothesis from the much weaker fact that $\Ind_\rho(C)$ is accessible.

\textbf{Example (cardinal form of the involution obstruction).}
For every infinite cardinal $\rho$ there is a small category $C$ such that
\[
 \Ind_\rho([B(\mathbf Z/2),C])\longrightarrow
 [B(\mathbf Z/2),\Ind_\rho(C)]
\]
is not essentially surjective.

\textit{Proof.} Put $\lambda=\rho^+$.  The objects of $C$ are the ordinals $\beta<\lambda$.  There are no arrows $\beta\to\gamma$ when $\beta>\gamma$, there is only the identity when $\beta=\gamma$, and there are two arrows labelled $+1,-1$ when $\beta<\gamma$.  Composition multiplies labels.  These rules are associative and define a small category.  Let $y:C\to\widehat C$ be Yoneda, and form
\[
 F=\mathop{\operatorname{colim}}_{\beta<\lambda}y(\beta),
\]
using the arrows labelled $+1$ as transitions.  The ordinal $\lambda$ is $\lambda$-filtered: a set of fewer than $\lambda$ ordinals below $\lambda$ has an upper bound below $\lambda$.  Hence $F$ belongs to $\Ind_\rho(C)$.

For a fixed $\gamma<\lambda$, evaluation gives
\[
 F(\gamma)=\mathop{\operatorname{colim}}_{\beta<\lambda}
            C(\gamma,\beta)=\{+1,-1\}.
\]
Indeed the terms for $\beta>\gamma$ are these two-element sets and all subsequent transition maps are identities.  Precomposition with an arrow labelled $\varepsilon$ multiplies by $\varepsilon$.  Multiplication by $-1$ on every $F(\gamma)$ therefore defines a natural automorphism $s:F\to F$, with $s^2=1$.  It defines an object $(F,s)$ of the right-hand diagram category.

Every functor $B(\mathbf Z/2)\to C$ has trivial action, because every endomorphism in $C$ is an identity.  A map of actions $(y\gamma,1)\to(F,s)$ would correspond by Yoneda to an element of $F(\gamma)$ fixed by multiplication by $-1$.  There is no such element.  If $(F,s)$ were a filtered colimit of $C$-valued actions, each object of the nonempty indexing category would give exactly such an equivariant map to its colimit.  This contradiction proves the assertion.\hfill$\square$

For $\lambda=\omega$ the signed category is the abstract form of the asymmetric finite subsets used in SGA~4.  The uncountable replacement supplies the stronger filteredness required here.  Thus Grothendieck's question is settled affirmatively, and the same construction is now part of a general characterization of the indexing categories for which presentations always exist.  The chosen category also explains the missing hypothesis: a strictly increasing countable chain has no colimit in $C$.  Any cocone has an apex strictly above all its terms and one constant sign.  Reversing that sign gives a second cocone to the same apex, which cannot factor through the first because the apex has only its identity endomorphism.  This contradicts the universal property of a proposed colimit.
\end{dcomment}


\printedpage{80}
% head: A. Grothendieck : Dérivateurs, Ch. XVIII, Édition des Universités de Montpellier II et Paris VII

\textbf{4.9.6} (\footnote[96]{Résultats préliminaires (à la démonstration de 4.9.3).}). Le reste de la présente section 4.9 est consacré à la démonstration de 4.9.3, i.e. à la vérification, pour $J \longrightarrow \cM$, des conditions a), b) de 4.9.2. Dans toute la suite, on se fixe donc une
\[
\varphi \ : \ J \longrightarrow \cM \ = \ \Ind_\pi(C) \ . \tag*{(4.9.6.1)}
\]
On peut supposer $J$ réduite, donc petite et $\card J \leqslant \pi_0$. On va considérer le foncteur déduit du foncteur (4.9.2.1)
\[
\uHom(J,C)/\varphi \longrightarrow \prod_{j \in \Ob J} C/\varphi(j) \ , \tag*{(4.9.6.2)}
\]
dont les foncteurs en question se déduisent en composant avec les foncteurs projection
\[
\prod C/\varphi(j) \longrightarrow C/\varphi(j_0) \ . \tag*{(4.9.6.3)}
\]
Comme les catégories facteurs $C/\varphi(j)$ sont filtrantes (et même grandes devant $\pi$), leur produit est également filtrant, et même grand devant $\pi$ (immédiat), et de même
\msp{page 222}
pour les produits partiels. Il s'ensuit aussitôt que les projections (4.9.6.3) sont cofinales, en vertu du

\dgrblock{dgrb1044}{Why the cardinal must grow with the diagram}
\begin{dcomment}{exposition}{p221}{Why the cardinal must grow with the diagram}
Corollary 4.9.5 applies the representation theorem after enlarging the class of presentable objects.  Its usefulness lies precisely in not claiming that an arbitrary diagram is assembled from the originally chosen presentability rank.  The preceding involution example shows how that stronger assertion can fail.  At a sufficiently large rank the subcategory $\mathcal M_\pi$ contains the intermediate colimits needed to enforce the diagram's equations.

The exponent condition has two separate jobs.  If $\pi=c^{\pi_0}$ and $c\geq\max(\pi_0,|\operatorname{Mor}J|)$, then
\[
 \pi^{\pi_0}=(c^{\pi_0})^{\pi_0}
 =c^{\pi_0\cdot\pi_0}=c^{\pi_0}=\pi,
 \qquad \pi\geq 2^{\pi_0}>\pi_0.
\]
This gives enough room to carry out the accessibility enlargement established earlier in Chapter XVIII, and enough room to record all arrows and equations of $J$.  Conversely, a cardinal satisfying the displayed fixed-point condition is of the indicated form by taking $c=\pi$.  No generalized continuum hypothesis is being used.

For intuition take $\mathcal M=\mathbf{Set}$ and let $\pi$ dominate the number of arrows of $J$.  The sets of cardinal at most $\pi$ are the $\pi^+$-presentable objects.  Given at most $\pi$ chosen elements of the values of a $J$-diagram, close them under its structure maps.  At most $\pi\cdot\pi=\pi$ elements are added.  The resulting subdiagram has pointwise cardinal at most $\pi$.  The poset of these subdiagrams is $\pi^+$-filtered: the union of at most $\pi$ of them still has that bound.  Their union is the original diagram.  Moreover the equalizer formula in the preceding block proves that every such bounded subdiagram is $\pi^+$-presentable in $[J,\mathbf{Set}]$.  This is the elementary model for the corollary.  In a general accessible category there are no underlying subsets on which to perform this closure; the comma categories and their cofinal functors provide the replacement.
\end{dcomment}


\textbf{Lemme 4.9.6.1.} \textit{Soient $I$, $I'$ deux catégories avec $I'$ \textup{0}-connexe (p.ex. filtrante), alors le foncteur projection $pr_1 : I \times I' \longrightarrow I$ est cofinal.}

En effet, pour $i' \in \Ob I'$ \texttt{[plutôt} $i \in \Ob I$\texttt{]}, on a
\[
i\backslash(I \times I') \ \simeq \ (i\backslash I) \times I' \ ,
\]
or $i\backslash I$ et $I'$ étant 0-connexe, il en est de même de leur produit.

\textbf{Lemme 4.9.6.2} (\footnote[97]{Attention, il n'est pas vrai qu'un produit infini de catégories 0-connexes est 0-connexe !}). \textit{Soit $(I_\alpha)_{\alpha \in A}$ une famille de catégories} filtrantes \textit{(resp. grandes devant $\pi$). Alors $I_A \stackrel{\text{déf}}{=} \prod_{\alpha \in A} I_\alpha$ est filtrante (resp. grande devant $\pi$), et pour toute partie $B \subset A$, la projection canonique $I_A \longrightarrow I_B$ est cofinal.} % [sic] cofinal -> cofinale, agreement with « la projection »

(\footnote[98]{Cf. démonstration élégante ci-dessous, 4.9.6.6.}). La stabilité de la propriété “filtrante” ou “grande devant $\pi$” par AQT \texttt{[âne qui trotte]}. On a, si $B' = A\backslash B$, $I_A \simeq I_B \times I_{B'} \longrightarrow I_B$, et on applique le lemme 4.9.6.1 pour conclure que ce foncteur est cofinal, $I_{B'}$ étant filtrante, donc 0-connexe.

\msp{page 223}Le foncteur (4.9.6.3) étant cofinal, pour établir que (4.9.2.1) l'est, il suffit de prouver qu'il en est ainsi du foncteur (4.9.6.2), ce que nous ferons. Notons de plus que ledit foncteur est \textit{fidèle}, car il s'insère dans le carré commutatif
\[
\begin{tikzcd}[column sep=large, row sep=large]
\uHom(J,C)/\varphi \arrow[r] \arrow[d] & \prod C/\varphi(j) \arrow[d] \\
\uHom(J,C) \arrow[r] & C^{\Ob J}\ ,
\end{tikzcd}
\]
\printedpage{81}
% head: A. Grothendieck : Dérivateurs, Ch. XVIII, Édition des Universités de Montpellier II et Paris VII
où les flèches verticales et la deuxième flèche horizontale sont fidèles. En définitive, pour établir le théorème 4.9.3, nous allons étudier le foncteur \textit{fidèle} (4.9.6.2)
\[
\begin{aligned}
&f \ : \ \uHom(J,C)/\Phi \ \longrightarrow \ \textstyle\prod_{j \in J_0} C/\varphi(j) \\ % [sic] \Phi -> \varphi, the functor is written \varphi everywhere else
&\text{(où on a posé } J_0 = \Ob J,\ J_1 = \Fl J) \ ,
\end{aligned}
\tag*{$(*)$}
\]
et nous allons prouver :
\begin{itemize}[label={}]
\item[\textcircled{a}] Ce foncteur fidèle $(*)$ est \textit{cofinal.}
\item[\textcircled{b}] La catégorie source $\uHom(J,C)/\varphi$ est grande devant $\pi_0$.
\end{itemize}
Nous allons pour ceci dégager un lemme auxiliaire sur la cofinalité des foncteurs \textit{fidèles} :

\dgrblock{dgrb1045}{Products need bounds on witnesses}
\begin{dcomment}{examples}{p222}{Products need bounds on witnesses}
The product trick is the first way Grothendieck makes independent choices simultaneous.  If each $I_a$ is $\lambda$-filtered and $D:K\to\prod_a I_a$ has fewer than $\lambda$ arrows, choose a cocone on each projected diagram $D_a$.  These coordinate cocones give a cocone on $D$.  The empty diagram supplies nonemptiness.  Thus arbitrary small products of $\lambda$-filtered categories are $\lambda$-filtered, using the usual axiom of choice.  No bound on the number of factors is needed: the same small diagram is being solved in every coordinate.

The projection to a subproduct is final because its comma category at $(x_b)$ is
\[
 \left(\prod_{b\in B}(x_b\backslash I_b)\right)
 \times\left(\prod_{a\notin B}I_a\right).
\]
Both factors are nonempty and filtered.  This calculation includes the empty subproduct, where the target is the terminal category.  For the simpler projection $I\times I'\to I$, only connectedness and nonemptiness of $I'$ are needed, since $(i\backslash I)\times I'$ has the same connected components as $I'$.

Footnote 97 correctly warns against replacing filteredness by connectedness in an infinite product.  Here is a full example.  Let $P_n$ be the finite fence poset on vertices $0,1,\ldots,n$, with the successive relations
\[
 0<1>2<3>4<\cdots.
\]
Its underlying undirected comparability graph is a path, so $P_n$ is connected, and any zigzag from $0$ to $n$ has at least $n$ nonidentity steps.  Consider the two objects $a=(0)_n$ and $b=(n)_n$ of $\prod_{n\geq1}P_n$.  A zigzag with $m$ steps between them would project to a zigzag with at most $m$ steps from $0$ to $n$ in every coordinate.  Taking $n>m$ is impossible.  Thus the product is disconnected.  Filteredness avoids this failure because two objects always have a common target, giving a uniformly two-step zigzag in every coordinate.  The issue is a uniform finite witness, not merely the existence of coordinatewise witnesses.
\end{dcomment}


\msp{page 224}\textbf{Lemme 4.9.6.3} (\footnote[99]{Ce lemme fait partie du sorite général sur la cofinalité et les catégories filtrantes.}). \textit{Soit $f : I \longrightarrow I'$ un foncteur, et considérons les conditions :}
{\itshape
\begin{itemize}[label={}]
\item[a)] $I$ filtrante, $f$ cofinal, i.e. (en présence de $I$ filtrante) $f$ satisfait les conditions familières \rbox{F 1}, \rbox{F 2}.
\item[b)] $I'$ est filtrante, et les $i'\backslash I$ (pour $i' \in \Ob I'$) sont filtrantes.
\item[b$'$)] $I'$ \textup{0}-connexe, satisfait PF (2) \textup{(}\footnote[100]{p.ex. PF 2 est satisfaite si $I'$ est \textit{ordonnée}, sans avoir à supposer d'avance $I'$ filtrante, seulement 0-connexe.}\textup{)} (égalisation des double-flèches).
\item[b$''$)] $I'$ est filtrante, les $i'\backslash I$ pseudo-filtrantes, i.e. satisfont à la condition PF $1'$.
\item[c)] $I'$ filtrante, et $f$ satisfait à F 1, ainsi que la condition suivante.
\begin{itemize}[label={}]
\item[\rbox{F 3}] Pour tout couple $(i_0,i_1)$ d'objets de $I$ et toute flèche $f(i_0) \xrightarrow{u'} f(i_1)$, il existe un objet $i_2$ de $I$, et une flèche $v' : f(i_1) \longrightarrow f(i_2)$, tels que $v'$ et $v'u'$ soient relevables en des flèches $i_1 \xrightarrow{v} i_2$ et $i_0 \xrightarrow{w} i_2$ (i.e. il existe $i_2$ et $v : i_1 \longrightarrow i_2$, $w : i_0 \longrightarrow i_2$ tels qu'on ait $f(v)u' = f(w)$).
\end{itemize}
\end{itemize}
\[
\begin{tikzcd}
i_0 \arrow[rr, bend left=20, "w"] & i_1 \arrow[r, "v"] & i_2 % doubt resolved against PDF; see corrections.csv
\end{tikzcd}
\]
\[
\begin{tikzcd}
f(i_0) \arrow[r, "u'"] \arrow[rr, dotted, bend right=20, "f(w)"'] & f(i_1) \arrow[r, "f(v) = v'"] & f(i_2)
\end{tikzcd}
\]
\begin{itemize}[label={}]
\item[1)] On a les implications
\[
a \iff b \iff b' \implies b'' \iff c \ ,
\]
et si $f$ est \textup{fidèle} (p.ex. $I$ ordonnée), les cinq conditions a), b), b$'$), b$''$), c) sont équivalentes.
\end{itemize}
}

\printedpage{82}
% head: A. Grothendieck : Dérivateurs, Ch. XVIII, Édition des Universités de Montpellier II et Paris VII
{\itshape
\begin{itemize}[label={}]
\item[2)] \msp{page 225}Soit $\pi$ un cardinal infini, et considérons les conditions :
\begin{itemize}[label={}]
\item[$a_\pi$)] $I$ grande devant $\pi$, $f$ cofinal, i.e. F 1, F 2.
\item[$b_\pi$)] $I'$ grande devant $\pi$, les $i'\backslash I$ ($i' \in \Ob I'$) aussi.
\item[$b_\pi{}''$)] $I'$ grande devant $\pi$, les $i'\backslash I$ ($i' \in \Ob I'$) sont $\pi$-pseudo-filtrantes, i.e. satisfont PS $1_\pi{}'$. % [sic] PS 1_π' -> PF 1_π', the conditions are elsewhere called PF
\item[$c_\pi$)] $I'$ grande devant $\pi$, et $f$ satisfait F 1 et
\begin{itemize}[label={}]
\item[\rbox{F $3_\pi$}] \textup{(}\footnote[101]{\textbf{N.B.} La condition F $3_\pi$ est tautologiquement satisfaite si $f$ est \textit{pleinement fidèle.}}\textup{)}. Pour toute famille $(i_\alpha)_{\alpha \in A}$ d'objets de $I$, avec \textup{card} $A \leqslant \pi$, et $i \in \Ob I$, et une famille de flèches $u'_\alpha : f(i_\alpha) \longrightarrow f(i)$, existe $j \in \Ob I$ et une flèche $v' : f(i) \longrightarrow f(j)$, telle que $v'$ \textup{et} les $v'u'_\alpha$ proviennent de flèches de $I$, i.e. il existe des flèches $v : i \longrightarrow j$, $w_\alpha : i_\alpha \longrightarrow j$ telles que $f(v)u'_\alpha = f(w_\alpha)$ $\forall\ \alpha \in A$.
\end{itemize}
\end{itemize}
\[
\begin{tikzcd}[row sep=huge]
i_\alpha \arrow[d, dash] \arrow[rr, dotted, bend left=20, "w_\alpha"] & i \arrow[r, dotted, "v"] \arrow[d, dash] & j \arrow[d, dash] \\
f(i_\alpha) \arrow[r, "u'_\alpha"] & f(i) \arrow[r, "v' = f(v)"] & f(j)
\end{tikzcd}
\]
Ceci posé, on a
\[
a_\pi \iff b_\pi \implies c_\pi \implies b_\pi{}'' \ ,
\]
et si $f$ est fidèle, ou $I$ ordonné, les quatre conditions $a_\pi$, $b_\pi$, $b_\pi{}''$, $c_\pi$ sont équivalentes.
\end{itemize}
}
(\footnote[102]{Démonstration de la proposition donnée dans 4.9.6.7, page 227.}).

\dgrblock{dgrb1046}{What the cofinality lemma actually requires}
\begin{dcomment}{exposition}{p224}{What the cofinality lemma actually requires}
Grothendieck has reduced a problem about natural transformations to a faithful functor into a product of comma categories.  The new lemma is designed to exploit that faithfulness: one may arrange equality in the target and then reflect it back to the source.  Keeping the two operations separate makes both the proof and its limitations transparent.

Write $D_d=(d\downarrow f)$ for $d\in I'$, and say explicitly that a category is \emph{object-directed} if it is nonempty and any two objects admit maps to a common object.  This is the particular property called ``pseudo-filtrante'' in this passage.  It does not include equalization of parallel arrows and should not be confused with other uses of ``pseudofiltered'' in the literature.  The coherent finite form of the lemma is
\[
 \begin{split}
 &I\text{ filtered and }f\text{ final}\quad\Longleftrightarrow\quad
 I'\text{ filtered and every }D_d\text{ filtered},\\
 &I'\text{ filtered and every }D_d\text{ object-directed}
 \quad\Longleftrightarrow\quad I'\text{ filtered, F1 and F3}.
 \end{split}
\]
The first line implies the second.  If $f$ is faithful, the second also implies the first.  F1 means that every $d$ admits an arrow to some $f(i)$; F3 says that every arrow $f(i_0)\to f(i_1)$ becomes the image of an arrow out of $i_0$ after postcomposition with the image of an arrow out of $i_1$.

There is an actual omission in the printed condition b$'$: it requires only that $I'$ be connected and satisfy PF2, without mentioning the comma categories.  Taken literally it cannot be equivalent to a): the functor from a two-object discrete category to the terminal category is faithful and satisfies that printed b$'$, but its source is not filtered.  The missing clause is ``every $D_d$ is filtered'', as the proof and the explicit parenthesis on manuscript p.~236 confirm.  With that clause restored, connectedness and PF2 for $I'$ suffice; the argument below proves this slightly stronger formulation.

The nonempty requirement in b$''$ must likewise remain visible.  If PF1$'$ is read only as a condition on pairs of objects, the empty category satisfies it vacuously, and $\varnothing\to1$ disproves b$''\Rightarrow$c.  Reading it as common upper bounds for all finite families, including the empty family, supplies precisely the required convention.  The cardinal version explicitly quantifies over all families of cardinal at most $\pi$ and therefore already includes the empty one.  This is a clarification of the finite convention, not a counterexample to the intended criterion.
\end{dcomment}


\textbf{4.9.6.4. N.B.} Dans le cas qui nous intéresse, où $f$ est le foncteur $(*)$ (p. 223), on sait que $I'$ est grande devant $\pi$. On voudrait établir que $I$ l'est et que $f$ est cofinal, i.e. la condition $\mathrm{a}_\pi$. On le prouvera en prenant $\mathrm{b}_\pi$ (qui implique $\mathrm{a}_\pi$ puisque $f$ est fidèle). Il reste donc à
\msp{page 226}
prouver, sur le foncteur $f$ dans $(*)$, p. 223, les deux conditions suivantes :
\begin{itemize}[label={}]
\item[F 1)] Pour tout système de flèches dans $\cM$
\[
X(j) \xrightarrow{p(j)} \varphi(j) \ , \qquad j \in J_0 = \Ob J, \text{ les } X(j) \text{ dans } \Ob C \ ,
\]
il existe un foncteur $F \ : \ J \longrightarrow C$ et $F \xrightarrow{q} \varphi$, enfin un système de flèches $X(j) \xrightarrow{u(j)} F(j)$ dans les $C/\varphi(j)$, i.e. rendant commutatif les diagrammes % [sic] commutatif -> commutatifs, agreement with « les diagrammes »
\[
\begin{tikzcd}[column sep=small]
X(j) \arrow[rr, "u(j)"] \arrow[dr, "p(j)"'] & & F(j) \arrow[dl, "q(j)"] \\
& \varphi(j)\ . &
\end{tikzcd}
\]
\item[F $3_\pi$)] Pour tout système $(F_\alpha)_{\alpha \in A}$ ($\card A \leqslant \pi$) d'éléments dans $\uHom(C,J)$, et des flèches % [sic] \uHom(C,J) -> \uHom(J,C), the F_α are functors J -> C
\[
F_\alpha \xrightarrow{p_\alpha} \varphi \qquad \text{dans } \uHom(J,\cM),
\]
et pour $G$ dans $\uHom(J,C)$ et $G \xrightarrow{q} \varphi$ dans $\uHom(J,\cM)$, étant donné un système de flèches $u_\alpha(j)$,
\printedpage{83}
% head: A. Grothendieck : Dérivateurs, Ch. XVIII, Édition des Universités de Montpellier II et Paris VII
\[
\begin{tikzcd}[column sep=small]
F_\alpha(j) \arrow[rr, "u_\alpha(j)", "\text{comm.}"'] \arrow[dr, "p_\alpha(j)"'] & & G(j) \arrow[dl, "q(j)"] \\
& \varphi(j)\ , &
\end{tikzcd}
\qquad u_\alpha(j) \text{ flèche dans } C/\varphi(j) \ ,
\]
il existe $H \xrightarrow{q} \varphi$ dans $\uHom(J,C)$ et une flèche $G \xrightarrow{v} H$ dans $\uHom(J,C)/\varphi$, telle que pour tout $\alpha \in A$, le système de flèches composés
\msp{page 227}
\[
\begin{tikzcd}[row sep=huge]
F_\alpha(j) \arrow[r, "u_\alpha(j)", "\text{comm.}"'] \arrow[dr, dotted, "p_\alpha(j)"'] & G(j) \arrow[r, "v(j)", "\text{comm.}"'] \arrow[d, dotted, "q(j)"] & H(j) \arrow[dl, dotted, "q(j)"] \\
& \varphi(j) &
\end{tikzcd}
\]
(pour $j \in J_0 \in \Ob J$) provienne d'un \textit{foncteur} $w_\alpha : F_\alpha \longrightarrow H$. % [sic] J_0 \in \Ob J -> J_0 = \Ob J, J_0 was defined as Ob J
\end{itemize}
Pour prouver 4.9.3, nous sommes donc ramené à prouver les deux propositions F 1), F $3_\pi$) précédentes, sous l'une des deux conditions a) ou b) de 4.9.3.

\textbf{Corollaire 4.9.6.5} (de 4.9.6.3). \textit{Soit $(f_\alpha : I_\alpha \longrightarrow I'_\alpha)$ une famille de foncteurs cofinaux, les catégories source étant filtrantes. Alors le foncteur produit $f = \prod f_\alpha : I = \prod_\alpha I_\alpha \longrightarrow I' = \prod_\alpha I'_\alpha$ est cofinal.}

En effet, il faut vérifier que les catégories $i'\backslash I$, pour $i' \in \Ob I'$, sont 0-connexes. Or si $i' = (i'_\alpha)_{\alpha \in A}$, on a
\[
i'\backslash I \ \simeq \ \prod_\alpha (i'_\alpha\backslash I_\alpha) \ ,
\]
et par 4.9.6.3, 1°) a $\Longrightarrow$ c, les $i'_\alpha\backslash I_\alpha$ sont filtrantes, donc leur produit est filtrant
\msp{page 228}
(cf. 4.4.6.2), donc 0-connexe, q.e.d. % [sic] 4.4.6.2 -> 4.9.6.2, the lemma on products of filtrantes is 4.9.6.2

\textbf{Corollaire 4.4.6.6.} \textit{Soit $\pi$ un cardinal infini. Tout produit de catégories grandes devant $\pi$ est grand devant $\pi$.} % [sic] 4.4.6.6 -> 4.9.6.6, numbering of the section and footnote 98

Soit $I = \prod_{\alpha \in A} I_\alpha$, soit $B$ un ensemble de cardinal $\leqslant \pi$, il faut prouver que $I \xrightarrow{\mathrm{diag}_I} I^B$ est cofinal. Or ce foncteur s'identifie au produit des foncteurs similaires $I_\alpha \longrightarrow I_\alpha^B$, lesquels sont cofinaux, les catégories source étant de plus filtrantes. Donc on peut conclure par le corollaire 4.9.6.5.

\textbf{4.9.6.7.} \textsc{Démonstration de} 4.9.6.3. (Partie 1°.)

a $\Longrightarrow$ b. On sait déjà que a) implique que $I'$ est filtrante (SGA 4 I 8.1.3 b). Reste à prouver que pour tout $i' \in \Ob I'$, $i'\backslash I$ est filtrante (\footnote[103]{Court-circuit : on sait déjà que $i'\backslash I$ est 0-connexe, d'autre part elle est \textit{localement} filtrante, car coétale sur $I$ qui est filtrante, donc elle est filtrante.}). Deux conditions PF $1'$ et PF 2 à vérifier (en plus du fait que $i'\backslash I \neq \emptyset$, qui est connu par cofinalité de $f$, impliquant que $i'\backslash I$ est 0-connexe).
\begin{itemize}[label={}]
\item[1)] Soient $(i,u')$, $(j,v')$ deux éléments de $i'\backslash I$, il faut voir qu'ils sont majorés par un troisième élément. Or $I$ étant filtrante, $i$, $j$ sont majorés par $k$. On aurait fini donc si on savait que $f(u_1)u' = f(v_1)v'$, soit $w'$, car alors $(k,w')$ ferait l'affaire.
\end{itemize}

% ---- from der-ch18_p086-090.tex
\printedpage{84}
% head: A. Grothendieck : Dérivateurs, Ch. XVIII, Édition des Universités de Montpellier II et Paris VII

\[
\begin{tikzcd}
{} & i \arrow[drr, "u_1"] \arrow[dd, dash] & {} & {} \\
{} & {} & j \arrow[r, "v_1"] & k \arrow[dd, dash] \\
i' \arrow[r, "u'"] \arrow[drr, "v'"'] & f(i) \arrow[drr] & {} & {} \\
{} & {} & f(j) \arrow[r] \arrow[uu, dash, crossing over] & f(k)
\end{tikzcd}
\]

\begin{itemize}[label={}]
\item \msp{page 229}Mais on applique F 2 sur $f$ à la double-flèche en question $i' \rightrightarrows f(k)$, on trouve $w : k \longrightarrow \bar{k}$,
\[
\begin{tikzcd}
{} & k \arrow[r, "w"] \arrow[dd, dash] & \bar{k} \arrow[dd, dash] \\
{} & {} & {} \\
i' \arrow[r, shift left] \arrow[r, shift right] & f(k) \arrow[r] & f(\bar{k})
\end{tikzcd}
\]
telle que
\[
\underbrace{f(w)(f(u_1)u')}_{= f(wu_1)u'} = \underbrace{f(w)(f(v_1)v')}_{f(wv_1)v'} \ ,
\]
et remplaçant $k$ par $\bar{k}$, et $u_1$, $v_1$ par $wu_1$, $wv_1$, on trouve ce qu'on veut.
\end{itemize}

\dgrblock{dgrb1047}{The finite proof, with the lifting equation in its proper order}
\begin{dcomment}{exposition}{p228}{The finite proof, with the lifting equation in its proper order}
The proof is a calculus of eventual equality.  It is useful to complete it once because the same operations reappear, with a cardinal bound, in the construction of sections.  We use ``final'' for Grothendieck's right-cofinal functors.

Suppose first that $I$ is filtered and $f$ is final.  A basic consequence is the following form of F2: if $r,s:d\rightrightarrows f(i)$, there is $w:i\to k$ such that $f(w)r=f(w)s$.  Indeed $(i,r)$ and $(i,s)$ are joined by a finite zigzag in $(d\downarrow f)$.  A cocone in $I$ on the finite underlying graph makes the two composites from $d$ equal.  The two resulting maps from $i$ to its apex may differ; equalizing them in $I$ gives the single map $w$ asserted.  Together with F1, this proves that $I'$ is filtered: two objects first map into the image of $f$ and then to a common target; a parallel pair first maps to an object $f(i)$ and is then equalized by F2.

The same argument proves that every $D_d$ is filtered.  For two of its objects, choose a common target for their underlying objects of $I$, and then use F2 to make the two maps from $d$ agree.  For parallel arrows in $D_d$, equalize their underlying arrows in $I$.  The resulting arrow remains in the comma category because its structure map is obtained by composition.  This proves a)$\Rightarrow$b), including the target-filteredness which is invoked from SGA~4 in the text.

Now suppose only the corrected b$'$: $I'$ is nonempty and connected, has right equalizers for parallel pairs, and each $D_d$ is filtered.  Finality follows immediately, hence $I$ is nonempty and connected.  First obtain F3 by taking a common target for $(i_0,1)$ and $(i_1,u')$ in $D_{f(i_0)}$.  To equalize $a,b:i_0\rightrightarrows i_1$, equalize $f(a),f(b)$ in $I'$, use F1 to move the target into the image of $f$, and use F3 to replace that target arrow by an image $f(w)$.  The resulting arrows $wa,wb$ have the same image.  They are parallel arrows in $D_{f(i_0)}$ from $(i_0,1)$ to $(\operatorname{cod}w,f(wa))$, so filteredness of that comma category equalizes them in $I$.

For an outgoing span $i_1\leftarrow i_0\to i_2$, a common target in $D_{f(i_0)}$ supplies maps $v_1,v_2$ whose composites with the span have the same image.  The equalization just proved makes the composites themselves equal.  Thus every outgoing span completes to a commuting square.  In a connected category this, together with parallel-pair equalization, gives filteredness: move along a finite zigzag between two objects and complete each successive outgoing span to keep a common target.  This proves corrected b$'$)$\Rightarrow$a).  The converse is immediate from b).

Finally suppose c).  F1 makes each $D_d$ nonempty.  Given two objects of $D_d$, use filteredness of $I'$ to form a common target for their maps from $d$, use F1 to put this target in the image of $f$, and apply F3 twice to lift both arrows after a further enlargement.  Thus $D_d$ is object-directed.  Conversely, that property for $D_{f(i_0)}$ gives F3 from $(i_0,1)$ and $(i_1,u')$, while nonemptiness gives F1.  If $f$ is faithful, a parallel pair in a comma category can be equalized in $I'$, the equalizing arrow lifted by F3 after F1, and the resulting equality reflected by faithfulness.  Hence b) follows.

The printed formula on p.~231, $f(w)v'=f(v)$, reverses the factors and is not typed: with $v:i_2\to\bar i_2$, $w:i_1\to\bar i_2$ and $v':f(i_1)\to f(i_2)$ the correct equation is $f(v)v'=f(w)$.  The proof above uses this order.  This is an edition/source slip, not a failure of the argument.  Its companion typo $\xi_1=(i_0,v')$ must likewise be read as $(i_1,v')$.
\end{dcomment}


\begin{enumerate}[label=\arabic*), start=2]
\item PF (2) pour $i'\bs I$, i.e. on se donne un diagramme ci-contre \texttt{[ci-dessous]},
\[
\begin{tikzcd}
{} & i \arrow[r, shift left, "u_1"] \arrow[r, shift right, "u_2"'] \arrow[dd, dash] & j \arrow[r, dotted, "w"] \arrow[dd, dash] & j' \\
{} & {} & {} & {} \\
i' \arrow[r, "u'"] & f(i) \arrow[r, shift left, "f(u_1)"] \arrow[r, shift right, "f(u_2)"'] & f(j) \arrow[r, dotted, "f(w)"] & f(j') \ ,
\end{tikzcd}
\]
avec $f(u_1)u' = f(u_2)u'$, soit $v'$ (double-flèche dans $i'\bs I$ de $(i,u')$ dans $(j,v')$), on veut l'égaliser à droite. Or $I$ étant filtrant, il existe $w : j \longrightarrow j'$ telle que $wu_1 = wu_2$, or cette flèche définit une flèche $(j,v') \longrightarrow (j', f(w')v')$, qui égalise la double-flèche donnée. % [sic] f(w')v' -> f(w)v', the arrow just defined is w : j --> j'
\end{enumerate}

b $\Longrightarrow$ b$'$ est clair.

b) $\Longrightarrow$ a). Comme les $i'\bs I$ sont filtrantes, donc 0-connexes, $f$ est cofinal, il reste à prouver que $I$ est filtrant, donc à vérifier PF 0 ($I$ 0-connexe), PF 1, PF 2 pour $I$. $I$ est 0-connexe, car $I'$ l'est, et $f$ cofinal, i.e. $\mathbf{W}_0$-coasphérique, d'où PF 0.

\msp{page 230}Prouvons PF 1, i.e. que si
\begin{tikzcd}[row sep=tiny, column sep=small, baseline=(\tikzcdmatrixname-2-1.base)]
{} & i_1 \\
i_0 \arrow[ur] \arrow[dr] & {} \\
{} & i_2
\end{tikzcd}
est donné, on peut compléter en un carré commutatif.
\[
\begin{tikzcd}[row sep=small, column sep=small]
{} & i_1 \arrow[dr, dotted, "v_1"] & {} \\
i_0 \arrow[ur, "u_1"] \arrow[dr, "u_2"'] & {} & i_3 \\
{} & i_2 \arrow[ur, dotted, "v_2"'] & {}
\end{tikzcd}
\]
% foot: Version 17 décembre 2025

\printedpage{85}
% head: A. Grothendieck : Dérivateurs, Ch. XVIII, Édition des Universités de Montpellier II et Paris VII

\[
\begin{tikzcd}[row sep=small, column sep=small]
{} & f(i_1) \arrow[dr, "f(v_1)"] & {} \\
f(i_0) \arrow[ur, "f(u_1)"] \arrow[dr, "f(u_2)"'] & {} & f(i_3) \\
{} & f(i_2) \arrow[ur, "f(v_2)"'] & {}
\end{tikzcd}
\]

Mais considérons les objets $\xi_1 = (i_1, f(u_1))$, $\xi_2 = (i_2, f(u_2))$ de $f(i_0)\bs I$, et utilisant PF 1$'$ pour cette catégorie, on trouve $\xi_3 = (i_3, w')$ dans ladite catégorie, et des flèches de $\xi_1$, $\xi_2$ dans $\xi_3$, i.e. $v_1 : i_1 \longrightarrow i_3$, \texttt{[}$v_2 : i_2 \longrightarrow i_3$\texttt{,]} telles que $\underbrace{f(v_1)f(u_1)}_{f(v_1u_1)} = w'$, $\underbrace{f(v_2)f(u_2)}_{f(v_2u_2)} = w'$, donc
\[
f(v_1u_1) \ = \ f(v_2u_2) \ .
\]
Ainsi, on a une double-flèche $i_0 \overset{v_1u_1}{\underset{v_2u_2}{\rightrightarrows}} i_3$, égalisée par $f$. Si on prouve PF 2 sur $I$, on trouve $i_3 \xrightarrow{\alpha} \bar{i}_3$ telle que $\alpha(v_1u_1) = \alpha(v_2u_2)$, i.e. $(\alpha v_1)u_1 = (\alpha v_2)u_2$, et quitte à remplacer $i_3$ par $\bar{i}_3$, $v_1$, $v_2$ par $\alpha v_1$, $\alpha v_2$, on a gagné.

Reste à prouver PF 2 sur $I$, i.e. égaliser une double-flèche
\[
i_0 \ \overset{v_1}{\underset{v_2}{\rightrightarrows}} \ i_1 \ .
\]

\msp{page 231}Or $I'$ satisfaisant à PF 2, on peut égaliser la double-flèche $f(u_1)$, $f(u_2)$ par une flèche $f(i_1) \xrightarrow{v'} i'$ dans $I'$, et comme $i'\bs I$ non vide, on peut supposer que $i'$ est de la forme $f(i_2)$.
\[
\begin{tikzcd}
i_0 \arrow[r, shift left, "u_1"] \arrow[r, shift right, "u_2"'] \arrow[dd, dash] & i_1 \arrow[r] & i_2 \\
{} & {} & {} \\
f(i_0) \arrow[r, shift left, "f(u_1)"] \arrow[r, shift right, "f(u_2)"'] & f(i_1) \arrow[r, "v'"] & f(i_2) % doubt resolved against PDF; see corrections.csv
\end{tikzcd}
\]
Admettons pour un instant la condition F 3 de c), que nous allons appliquer à $(i_1, i_2, v')$. Il existe donc $i_2 \xrightarrow{v} \bar{i}_2$, $w : i_1 \longrightarrow \bar{i}_2$ tels que l'on ait $f(w)v' = f(v)$. Il s'ensuit que $f(w)$ égalise $f(u_1)$ et $f(u_2)$, i.e. qu'on a $f(wu_1) = f(wu_2)$, i.e. la double-flèche initiale est remplacée par $(\bar{u}_1 = wu_1, \bar{u}_2 = wu_2)$, telle que $f(\bar{u}_1) = f(\bar{u}_2)$. En somme, on est ramené au cas d'une double-flèche $(u_1, u_2)$ telle que $f(u_1) = f(u_2)$, soit $v'$. Mais alors soit $\xi_1 = (i_0, v')$, $\xi_0 = (i_0, \mathrm{id}_{f(i_0)})$, on a donc une double-flèche $(u_1, u_2) : \xi_0 \rightrightarrows \xi_1$ dans $f(i_0)\bs I$. Cette catégorie % [sic] \xi_1 = (i_0, v') -> \xi_1 = (i_1, v'), u_1, u_2 go from i_0 to i_1
\msp{page 232}est filtrante, donc cette double-flèche est égalisable à droite, ce qui implique aussitôt que la double-flèche $(v_1, v_2)$ \textit{dans $I$} est aussi égalisable à droite.

Donc il reste à prouver que b$'$) implique la condition F 3. Mais considérons les deux objets $\xi_0 = (i_0, \mathrm{id}_{f(i_0)})$, $\xi_1 = (i_1, u')$ de $f(i_0)\bs I$, et utilisons la condition de filtration PF 1 sur cette catégorie, donc on trouve un objet $\xi_2 = (i_2, v')$, et deux flèches
\begin{tikzcd}[row sep=tiny, column sep=small, baseline=(\tikzcdmatrixname-2-2.base)]
\xi_0 \arrow[dr, "w"] & {} \\
{} & \xi_2 \\
\xi_1 \arrow[ur, "v"'] & {}
\end{tikzcd}
, i.e. $i_0 \xrightarrow{w} i_2$
% foot: Version 17 décembre 2025
\printedpage{86}
% head: A. Grothendieck : Dérivateurs, Ch. XVIII, Édition des Universités de Montpellier II et Paris VII
et $i_1 \xrightarrow{v} i_2$ telles qu'on ait $f(w)\mathrm{id}_{i_0} = f(v)u' = v'$, d'où $f(w) = f(v)u'$, ce qui est la condition F 3.
\[
\begin{tikzcd}
i_0 \arrow[rr, dotted, bend left, "w"] \arrow[dd, dash] & i_1 \arrow[r, dotted, "v"] \arrow[dd, dash] & i_2 \arrow[dd, dash] \\
{} & {} & {} \\
f(i_0) \arrow[r, "u'"] \arrow[rr, dotted, bend right, "v'"'] & f(i_1) & f(i_2)
\end{tikzcd}
\]

Ainsi on a établi
\[
\text{a} \ \Longleftrightarrow \ \text{b} \ \Longleftrightarrow \ \text{b}' \ ,
\]
prouvons maintenant

a $\Longrightarrow$ c. On sait déjà que a) implique $I'$ filtrante, et il reste à établir F 3. Or $I$ étant filtrante, il existe des flèches
\begin{tikzcd}[row sep=tiny, column sep=small, baseline=(\tikzcdmatrixname-2-2.base)]
i_0 \arrow[dr, "w"] & {} \\
{} & i_2 \\
i_1 \arrow[ur, "v"'] & {}
\end{tikzcd}
. On
\msp{page 233}n'a pas forcément $f(w) = f(v)u'$, on a là une double-flèche $f(i_0) = f(i_1)$ dans $I'$. Mais par la condition F 2 sur $f$, il existe $i_2 \xrightarrow{\alpha} \bar{i}_2$ dans $I$ telle que $f(\alpha)$ égalise ladite double-flèche, donc $f(\underbrace{\alpha w}_{\bar{w}}) = f(\underbrace{\alpha v}_{\bar{v}})u'$, et quitte à remplacer $i_2$ par $\bar{i}_2$, et $w$, $v$ par $\bar{w} = \alpha w$, $\bar{v} = \alpha v$, on gagne. % [sic] f(i_0) = f(i_1) -> f(i_0) \rightrightarrows f(i_2), the double arrow f(w), f(v)u' goes from f(i_0) to f(i_2)

Reste à prouver (pour la partie $1^\circ$) que si $f$ est \textit{fidèle}, alors c) $\Longrightarrow$ b). Il reste donc à prouver que c) implique que pour $i'$ dans $I'$, $i'\bs I$ est filtrante. On sait déjà par F 1 que $i'\bs I$ est non vide, il reste à prouver les conditions PF 1$'$, PF 2 pour $i'\bs I$. Comme $I'$ filtrante, l'on peut compléter
\[
\begin{tikzcd}[row sep=small, column sep=small]
{} & f(i_1) \arrow[dr, dotted] & {} \\
i' \arrow[ur] \arrow[dr] & {} & j' \ , \\
{} & f(i_2) \arrow[ur, dotted] & {}
\end{tikzcd}
\]
et par F 1 on peut supposer $j'$ de la forme $f(i_3)$.
% foot: Version 17 décembre 2025

\printedpage{87}
% head: A. Grothendieck : Dérivateurs, Ch. XVIII, Édition des Universités de Montpellier II et Paris VII

\[
\begin{tikzcd}[column sep=small]
{} & i_1 \arrow[drr, dotted, "v_1\,?"] \arrow[dd, dash] & {} & {} & {} \\
{} & {} & {} & i_3 \arrow[r] \arrow[dd, dash] & \bar{i}_3 \arrow[dd, dash] \\
{} & f(i_1) \arrow[drr, "v_1'"] & i_2 \arrow[ur, dotted, "v_2\,?"'] \arrow[dd, dash, crossing over] & {} & {} \\
i' \arrow[ur, "u_1'"] \arrow[dr, "u_2'"'] & {} & {} & f(i_3) \arrow[r] & f(\bar{i}_3) \\
{} & {} & f(i_2) \arrow[ur, "v_2'"'] & {} & {}
\end{tikzcd}
\]
Appliquant F 3 à $i_1$, $i_3$, $v_1'$, on peut supposer (quitte à remplacer $i_3$ par $\bar{i}_3$)
\msp{page 234}que $v_1'$ provient d'un $v_1 : i_1 \longrightarrow i_3$. Appliquant à nouveau F 3 à $i_2$, $i_3$, $v_2'$, on peut supposer que $v_2'$ également provient d'un $v_2 : i_2 \longrightarrow i_3$. Alors $\xi_3 = (i_3, v_1'u_1' = v_2'u_2')$ est un majorant commun de $\xi_1 = (i_1, u_1')$, $\xi_2 = (i_2, u_2')$.

Prouvons enfin PF 2 pour $i'\bs I$, donc on a un diagramme ci-contre \texttt{[ci-dessous]} avec $f(u_1)u' = f(u_2)u'$ (double-flèche dans $i'\bs I$), on veut égaliser dans $i'\bs I$.
\[
\begin{tikzcd}
{} & i_0 \arrow[r, shift left, "u_1"] \arrow[r, shift right, "u_2"'] & i_1 \arrow[rr, bend left, "w"] & i_2 \arrow[r] & \bar{i}_2 \\
{} & {} & {} & {} & {} \\
i' \arrow[r, "u'"] & f(i_0) \arrow[r, shift left, "f(u_1)"] \arrow[r, shift right, "f(u_2)"'] & f(i_1) \arrow[r, dotted, "v'"] & f(i_2) \arrow[r] & f(\bar{i}_2)
\end{tikzcd}
\]
Mais $I'$ étant filtrante, on peut égaliser $(f(u_1), f(u_2))$ par $f(i_1) \longrightarrow j'$, et à cause de F 1 on peut supposer $j' = f(i_2)$. Appliquant encore F 3 à $(i_1, i_2, v')$, on trouve $i_1 \longrightarrow \bar{i}_2$ tel que $f(w)$ égalise $f(u_1)$, $f(u_2)$ :
\[
f(wu_1) \ = \ f(wu_2) \ .
\]
Mais $f$ étant \textit{fidèle}, cela implique $wu_1 = wu_2$, et on a donc égalisé la double-flèche de $i'\bs I$ par $w : (i_1, -) \longrightarrow (\bar{i}_2, -)$.

\subsection*{4.9.6.8. N.B.}

\begin{enumerate}[label=\arabic*$^\circ$)]
\item Si on ne suppose pas que $f$ soit fidèle, on voit que c) implique pourtant que les $i'\bs I$ satisfont PF 1$'$ (elles sont donc pseudo-filtrantes et 0-connexes). Donc $f$ est cofinal.

\msp{page 235}On voit, par le même argument, que $I$ satisfait PF 1$'$, i.e. est pseudo-filtrante. Ainsi, si on introduit la condition :
\begin{itemize}[label=\encadre{b$''$}]
\item $I'$ est filtrante, les $i'\bs I$ ($i \in \mathrm{Ob}\, I'$) sont pseudo-filtrantes (i.e. satisfont PF 1$'$) (donc $f$ cofinal).
\end{itemize}
On trouve les implications
\[
(\text{a} \ \Longleftrightarrow \ \text{b} \ \Longleftrightarrow \ \text{b}') \ \ \Longrightarrow \ \ (\text{b}'' \ \Longleftrightarrow \ \text{c}) \ ,
\]
et les conditions équivalentes c), b$''$) impliquent que $I$ est également pseudo-filtrante, et que pour toute double-flèche $\xi_0 \overset{u_1}{\underset{u_2}{\rightrightarrows}} \xi_1$ dans un $i'\bs I$ ou dans $I$, il existe une flèche $\xi_1 \xrightarrow{v} \xi_2$ dans la même catégorie, telle que $vu_1$ et $vu_2 \ : \ \xi_0 \rightrightarrows \xi_2$ aient \textit{même image}
% foot: Version 17 décembre 2025
\printedpage{88}
% head: A. Grothendieck : Dérivateurs, Ch. XVIII, Édition des Universités de Montpellier II et Paris VII
\textit{dans $I'$}. Donc si $f$ est fidèle, alors ces conditions c, c$'$ \texttt{[plutôt b$''$]} impliquent que $I$ et les $i'\bs I$ satisfont aussi la condition PF 2 de la double-flèche, donc sont filtrantes, donc dans ce cas les cinq conditions a) b) b$'$) c) c$'$) \texttt{[plutôt b$''$)]} sont équivalentes. Mais aussi lorsque $I$ est ordonnée, on sait [\,?] que la condition de la double-flèche dans $I$ est satisfaite automatiquement.

\item \msp{page 236}Notons aussi que dans la condition b$'$), on ne peut supprimer l'hypothèse que $I'$ satisfait PF 2 (en plus de $I'$ 0-connexe et les $i'\bs I$ filtrantes). Pour le voir, prenons $I' = I$, $f = \mathrm{id}$, donc l'hypothèse a ($\Longleftrightarrow$ b $\Longleftrightarrow$ b$'$) signifie que $I$ est filtrante, $f$ étant bien sûr cofinal. Mais b$'$ affaiblie signifie que $I$ est 0-connexe et les $i\bs I$ sont filtrantes, et on sait que ceci n'implique \textit{pas} que $I$ soit filtrante, même si on suppose $I$ pseudo-filtrante. Un exemple simple est donné par $I$ définie par
\[
\mathrm{Ob}\, I \ = \ \mathbf{N} \ , \qquad \mathrm{Hom}_I(i,j) \ = \ \left\{
\begin{array}{rcccc}
\emptyset & \text{si} & i & > & j \\[1ex]
\{\mathrm{id}_i\} & \text{si} & i & = & j \\[1ex]
\{\pm 1\} & \text{si} & i & < & j \ ,
\end{array}
\right.
\]
la composition des flèches se faisant par la multiplication des signes $\varepsilon_{k,j}\varepsilon'_{j,i} = (\varepsilon\varepsilon')_{k,i}$ (en posant $\mathrm{id}_i = (+1)_{i,i}$). Les catégories $i\bs I$ sont isomorphes entre elles et sont \textit{ordonnées} et pseudo-filtrantes, donc filtrantes, d'autre part, $I$ est pseudo-filtrante, pourtant elle n'est pas filtrante, car ne satisfait pas PF 2.

\item D'autre part, si $f$ n'est pas fidèle, alors c) (équivalente à $b''$) n'implique pas a), comme on voit en prenant $I' = e$ ; alors b$''$) équivaut à $I$ pseudo-filtrante, qui n'implique pas que $I$ soit filtrante.
\end{enumerate}

\subsection*{\msp{page 237}\textbf{4.9.6.9.} \textsc{Démonstration de 4.9.6.3, partie} $2^\circ$.}

\dgrblock{dgrb1048}{Connected slices cannot replace equalization}
\begin{dcomment}{examples}{p236}{Connected slices cannot replace equalization}
Grothendieck's signed category measures exactly the part of filteredness lost when one retains only common upper bounds.  Its objects are natural numbers, and for $n<m$ there are two arrows $n\to m$, labelled $+1$ and $-1$.  Composition multiplies signs.  Any two objects map to a larger integer, so the category is object-directed and connected.  But two opposite-sign arrows can never be equalized on the right: multiplying both by the same sign leaves them opposite.  It is therefore not filtered.

Fix $n$.  The objects of $n\backslash I$ are $1_n$ and pairs $(m,\varepsilon)$ for $m>n$.  Between $(m,\varepsilon)$ and $(k,\delta)$ with $m<k$ there is exactly one possible arrow, namely the arrow labelled $\delta\varepsilon^{-1}$.  At equal levels the only arrows are identities between equal objects.  Thus the slice is a poset, consisting of a bottom point followed by two points at each subsequent level, both below every point at a higher level.  Two objects have an upper bound at any strictly larger level; hence the slice is filtered.  Shifting levels shows that these slices are mutually isomorphic.  Taking $f=1_I$ gives exactly the counterexample promised on p.~236: connectedness of $I'$ and filteredness of all $D_d$ do not imply filteredness of $I$ unless PF2 is retained for $I'$.

There is an even smaller illustration of the role of faithfulness.  For a nontrivial group $G$, the functor $BG\to1$ satisfies F1 and F3: there is only one target arrow, and one can choose both lifting arrows to be the identity.  Its source is object-directed, but not filtered, since $1$ and any $g\ne1$ cannot be equalized by multiplying on the left by an invertible element.  The functor is final because $BG$ is connected.  Finality alone is consequently weaker than having a filtered source.  The same example works for every cardinal version of F3.
\end{dcomment}


C'est essentiellement “la même” démonstration que pour la partie $1^\circ$).

a$_\pi$) $\Longrightarrow$ b$_\pi$). On peut court-circuiter la démonstration directe, en regardant, pour un ensemble $A$ de cardinal $\leqslant \pi$, le diagramme
\[
\begin{tikzcd}
I \arrow[r, "\delta_A"] \arrow[d, "f"'] & I^A \arrow[d, "f^A"] \\
I' \arrow[r, "\delta_{A'}"] & I'^A
\end{tikzcd}
\]
\texttt{[plutôt $\delta'_A$]}. Comme $f$ est cofinal, $I$ filtrante, $f^A$ est aussi cofinal (4.9.8.5), d'autre part, $I$ grande devant $\pi$ signifie que $\delta_A$ est cofinal ($\forall$ $A$ avec $\mathrm{card}\, A \leqslant \pi$). Donc $f^A\delta_A = \delta_{A'}f$ est cofinal, donc $\delta_{A'}f$ et $f$ le sont, donc $\delta_{A'}$ aussi, ce qui signifie que $I'$ est grande devant $\pi$.

Reste à prouver que pour $i'$ dans $I'$, $i'\bs I$ est grande devant $\pi$. Cela résulte du fait que $i'\bs I \xrightarrow{p} I$ est coétale, que $i\bs I$ est 0-connexe, et du

\textbf{Lemme 4.9.6.9.1.} \textit{Soit $I$ une catégorie filtrante (resp. grande devant $\pi$). Alors pour toute $J$ sur $I$ coétale sur $I$ et $0$-connexe, $J$ est filtrante (resp. grande devant $\pi$).}
% foot: Version 17 décembre 2025

% ---- from der-ch18_p091-095.tex
\printedpage{89}
% head: A. Grothendieck : Dérivateurs, Ch. XVIII, Édition des Universités de Montpellier II et Paris VII
% foot: Version 17 décembre 2025

Le cas filtrant est déjà connu (cf. II \dots). Pour le cas “grande devant $\pi$”, on va utiliser le

\textbf{Corollaire 4.9.6.9.2.} \textit{Soit $J$ une catégorie. Pour que $J$ soit grande devant $\pi$, il}
\msp{page 238}
\textit{faut et il suffit qu'elle soit $0$-connexe, et satisfasse PF$_\pi$ 2 (cf. page 81), ainsi que la variante suivante de PF$_\pi$ 1$'$.}
\begin{itemize}[leftmargin=4em]
\item[\textit{PF$_\pi$ 1)}] \textit{Pour tout objet $x_0$ de $J$, $x_0\backslash J$ satisfait PF$_\pi$ 1$'$, i.e. pour toute famille de flèches $x_0 \xrightarrow{u_\alpha} y_\alpha$ ($\alpha \in A$, $\card A \leqslant \pi$), existe une flèche $x_0 \xrightarrow{v} z$ et des flèches $v_\alpha : y_\alpha \longrightarrow z$, avec $v_\alpha u_\alpha = v \ \forall\, \alpha \in A$.}
\end{itemize}

\dgrblock{dgrb1049}{A connected category of elements inherits filteredness}
\begin{dcomment}{exposition}{p237}{A connected category of elements inherits filteredness}
The shortcut through a coétale functor expresses a useful structural fact.  In modern language the functor $p:E\to I$ is a discrete opfibration: for every $e\in E$ and arrow $p(e)\to i$ there is a unique lift with source $e$.  Equivalently, $E$ is the category of elements of a covariant set-valued functor on $I$.  The projection $(d\downarrow f)\to I$ is of this form, associated with $i\mapsto I'(d,f(i))$.

Assume $I$ is $\lambda$-filtered and $E$ is nonempty and connected.  We give a direct proof that $E$ is $\lambda$-filtered.  Any finite zigzag in $E$ has a common target: choose a cocone on its projected finite diagram in $I$ and lift its legs.  The lifted endpoint objects coincide along each edge by uniqueness of lifting, and hence along the entire zigzag.  This proves that any two objects of $E$ have a common target.  A parallel pair in $E$ is equalized by lifting an equalizing arrow in $I$; uniqueness of lifting identifies the resulting composites.  Thus $E$ is filtered.

For a family $(e_a)$ of fewer than $\lambda$ objects, choose a base object $e_0$ and, using the finite case, a common target $z_a$ of $e_0$ and $e_a$.  The projected arrows $p(e_0)\to p(z_a)$ have a common compatible target in $I$ by $\lambda$-filteredness.  Lift these arrows from the $z_a$.  Their targets are all the same object, because the composites starting at $e_0$ are unique lifts of the same base arrow.  It receives arrows from every $e_a$.  A family of fewer than $\lambda$ parallel arrows is equalized by the same lifting argument as before.  Common upper bounds and simultaneous equalization now give cocones on arbitrary diagrams with fewer than $\lambda$ arrows.  This proves Lemma 4.9.6.9.1 completely.

The proof also explains Corollary 4.9.6.9.2.  Connectedness plus the common-target condition for every family issuing from one object first gives the finite common-target property by following zigzags.  For a larger family, join each member to a fixed base object and then apply the stipulated condition to those arrows from the base.  Together with simultaneous equalization, this is exactly $\lambda$-filteredness.  Connectedness cannot be discarded: a disjoint union of two terminal categories satisfies every condition localized at an individual object but has no cocone on the diagram consisting of one object in each component.
\end{dcomment}


\textsc{Démonstration}. Si $J$ est grande devant $\pi$, alors PF$_\pi$ 1) résulte aussitôt de la conjonction de PF$_\pi$ 1$'$ et PF$_\pi$ 2, d'où le “il faut”. Inversement, supposons que $J$ satisfasse PF$_\pi$ 1, PF$_\pi$ 2, et soit 0-connexe, il faut prouver qu'elle satisfait PF$_\pi$ 1$'$. Mais on sait déjà que les conditions en question, qui impliquent PF 1, PF 2, 0-connexe, impliquent que $J$ est filtrante. Donc si on a une famille $(x_\alpha)_{\alpha\in A}$, $\card A \leqslant \pi$, choisissant un $x_0 \in \Ob J$, on peut pour tout $\alpha$ trouver $y_\alpha$ qui majore $x_0$ et $x_\alpha$, d'où des flèches
\begin{tikzcd}[row sep=0.4em, column sep=1em, baseline=(\tikzcdmatrixname-2-2.base)]
x_\alpha \arrow[dr] & \\
& y_\alpha \\
x_0 \arrow[ur] &
\end{tikzcd}, on peut leurs appliquer PF$_\pi$ 1, d'où des flèches $y_\alpha \longrightarrow z$, et en composant des $x_\alpha \longrightarrow z$, q.e.d. % [sic] leurs appliquer -> leur appliquer, grammar

Prouvons maintenant le lemme. Il suffit (par le corollaire précédent) de prouver que $J$ satisfait à PF$_\pi$ 1 et PF$_\pi$ 2. Or PF$_\pi$ 1 est immédiat, puisque $\forall\, x_0$ dans $J$, $x_0\backslash J \xrightarrow{\ \sim\ } y_0\backslash I$, où $y_0$ est
\msp{page 239}
l'image de $x_0$ dans $I$. Quant à PF$_\pi$ 2, si on a des $x_0 \trois{u_\alpha} x_1$ dans $J$, d'où des $g(x_0) \trois{g(u_\alpha)} g(x_1)$, on égalise dans $I$ par $g(x_1) \xrightarrow{w} y$, mais comme $g$ est coétale, $\exists\, x_1 \xrightarrow{v} x_2$ dans $J$ telle que $g(v) = w$, donc on a $g(vu_\alpha) = g(vu_\beta)$ pour $\alpha, \beta \in A$, et comme $g$ est fidèle, on en conclut $vu_\alpha = vu_\beta$, donc $v$ égalise $(u_\alpha)_{\alpha\in A}$, q.e.d.

Prouvons b$_\pi$) $\Longrightarrow$ a$_\pi$) : Si $I'$ et les $i'\backslash I$ sont grandes devant $\pi$, alors $I$ l'est aussi (et $f$ est cofinal, les $i'\backslash I$ étant 0-connexes, car filtrantes). Comme on sait déjà que $I$ est 0-connexe ($I'$ l'étant et $f$ cofinal), il reste, par 4.9.6.9.2, à vérifier PF$_\pi$ 1 et PF$_\pi$ 2.

\msp{page 240}
La démonstration donnée dans 1° pour le cas similaire filtrant (implication b$'$ $\Longrightarrow$ a, pages 231-232), s'applique mutatis mutandis. (La propriété F 3 a déjà été établie dans la partie 1°), comme conséquence de b, et peut donc être utilisée librement.)

Donc on a prouvé a$_\pi$ $\Longleftrightarrow$ b$_\pi$. Prouvons a$_\pi$ $\Longrightarrow$ c$_\pi$. On sait déjà que a$_\pi$ implique que $I'$ est grande devant $\pi$, et bien sûr la condition F 1, reste à établir la validité de la condition F 3$_\pi$. Par F 3, $\forall\, \alpha \in A$ on peut trouver $j_\alpha$ dans $I$ et $w^0_\alpha : i_\alpha \longrightarrow j_\alpha$, $v^0_\alpha : i \longrightarrow j$ tels que % [sic] v^0_\alpha : i \longrightarrow j -> v^0_\alpha : i \longrightarrow j_\alpha, the following diagram has v^0_\alpha : i \to j_\alpha
\[
f(v_\alpha)u'_\alpha \ = \ f(w_\alpha) \ .
\]

\printedpage{90}
% head: A. Grothendieck : Dérivateurs, Ch. XVIII, Édition des Universités de Montpellier II et Paris VII
% foot: Version 17 décembre 2025

\[
\begin{tikzcd}[column sep=large, row sep=huge]
i_\alpha \arrow[rr, bend left=20, "w^0_\alpha"] \arrow[d, dash] & i \arrow[r, "v^0_\alpha"] \arrow[d, dash] & j_\alpha \arrow[r, "\lambda_\alpha"] \arrow[d, dash] & j \arrow[d, dash] \\
f(i_\alpha) \arrow[r, "u'_\alpha"] & f(i) \arrow[r, "f(v_\alpha)"] & f(j_\alpha) \arrow[r] & f(j)
\end{tikzcd}
\]

Comme $I$ satisfait PF 1$_\pi$ (cf. 4.9.6.9.2), il existe $h$ dans $I$ et des $\lambda_\alpha : j_\alpha \longrightarrow j$, tels que $\lambda_\alpha v^0_\alpha = \lambda$ soit indépendant de $\alpha$, alors on a pour tout $\alpha \in A$ % [sic] h -> j, the arrows \lambda_\alpha go to j
\[
f(\lambda)u'_\alpha \ \Bigl( \ = \ f(\lambda_\alpha)\underbrace{f(v_\alpha)u'_\alpha}_{f(w_\alpha)} \ \Bigr) \ = \ f(\underbrace{\lambda_\alpha w^0_\alpha}_{w_\alpha}) \ ,
\]
donc $\lambda : i \longrightarrow j$ et les $w_\alpha = \lambda_\alpha w^0_\alpha : i_\alpha \longrightarrow j$ satisfont aux conditions voulues.

c$_\pi$ $\Longrightarrow$ b$''_\pi$. Il faut prouver que si c$_\pi$ est satisfait, alors les $i'\backslash I$ satisfont PF 1$'_\pi$.

\msp{page 241}
Soit donc $\xi_\alpha = (i_\alpha, u'_\alpha)$ ($\alpha \in A$, $\card A \leqslant \pi$) une famille d'éléments de $i'\backslash I$, comme $I'$ satisfait à PF 1$_\pi$ (cf. 4.9.6.9.2), il existe $j'$ dans $I'$ et des flèches $u''_\alpha : f(i_\alpha) \longrightarrow j'$, telles que $u''_\alpha u'_\alpha = u'$ indépendant de $\alpha$. Comme $f$ satisfait F 1, on peut supposer que $j'$ est de la forme $f(j)$. Appliquant F 3$_\pi$ aux $i_\alpha$, $j$, $u''_\alpha$, on trouve $v : j \longrightarrow k$ et des $w_\alpha : i_\alpha \longrightarrow k$, tels que
\[
f(v)u''_\alpha \ = \ f(w_\alpha) \qquad \forall\, \alpha \in A \ ,
\]
d'où
\[
f(w_\alpha)u'_\alpha \ = \ f(v)\,\underbrace{u''_\alpha u'_\alpha}_{u'} \ = \ f(v)u' \qquad \text{indépendant de $\alpha$,}
\]
soit $w'$, donc les $w_\alpha$ définissent une flèche des $\xi_\alpha$ dans $\xi = (k, w')$, q.e.d.

\dgrblock{dgrb1050}{Simultaneous lifting at the cardinal bound}
\begin{dcomment}{exposition}{p240}{Simultaneous lifting at the cardinal bound}
The cardinal proof requires one simultaneous enlargement.  Repeating F3 separately for an infinite family would give unrelated targets, so the common-domain form of filteredness is essential.  Put $\lambda=\pi^+$.  For a functor $f:I\to I'$, with $I'$ $\lambda$-filtered, let $\mathrm{F3}_\pi$ have exactly the meaning stated on p.~225.  Then F1 and $\mathrm{F3}_\pi$ are equivalent to every $(d\downarrow f)$ being nonempty and admitting common targets for all families of at most $\pi$ objects.  The reverse implication, not displayed in the printed implication chain, follows by applying this property to $(i,1)$ and the objects $(i_a,u'_a)$ of $(f(i)\downarrow f)$.  There are at most $\pi$ of these objects because $\pi$ is infinite.

For the forward implication, a family $(i_a,r_a:d\to f(i_a))$ first acquires a compatible target in $I'$.  By F1 its target may be replaced by $f(j)$.  Apply $\mathrm{F3}_\pi$ to all arrows $f(i_a)\to f(j)$ at once.  The resulting maps $i_a\to k$ make the composites from $d$ equal, so they are a common target in the comma category.  If $f$ is faithful, a family of at most $\pi$ parallel comma arrows is equalized by first equalizing their images in $I'$, moving the target into the image of $f$, and lifting the resulting arrow by the ordinary F3.  Faithfulness reflects the common image.  The comma categories are therefore $\lambda$-filtered.

We also fill in the cardinal descent which the source leaves to ``mutatis mutandis''.  Suppose $I'$ and every $(d\downarrow f)$ are $\lambda$-filtered.  For a diagram $D:K\to I$ with at most $\pi$ arrows, choose a cocone $fD\to d$ in $I'$ and then a map $d\to f(i)$ by F1.  The resulting family $fD(k)\to f(i)$ lifts simultaneously by $\mathrm{F3}_\pi$ to maps $w_k:D(k)\to j$.  Their naturality equations hold after applying $f$.  For each arrow $a:k\to l$ of $K$, the pair $w_k,w_lD(a)$ is therefore a parallel pair in $(fD(k)\downarrow f)$ with target $(j,f(w_k))$.  Each such pair can be equalized in $I$ by filteredness of that comma category.  There are at most $\pi$ choices of an equalizing map out of $j$; filteredness of $(f(j)\downarrow f)$ gives a compatible common successor for them.  Composing $w_k$ with its map from $j$ makes every naturality equation hold.  This is the required cocone on $D$.

Finally, if $I$ is $\lambda$-filtered and $f$ final, its comma categories are connected discrete opfibrations over $I$, hence $\lambda$-filtered by the previous block.  The target $I'$ inherits $\lambda$-filteredness by choosing image targets for all objects of a small diagram, and then using the filtered comma categories to impose its at most $\pi$ relations; equivalently, the diagonal argument printed on p.~237 applies.  These arguments prove all intended cardinal implications, with no need to assume $\pi$ regular.  The only regular cardinal used in the modern formulation is $\pi^+$.
\end{dcomment}


b$''_\pi$ $\Longrightarrow$ b$_\pi$ si $f$ est fidèle. On est ramené à prouver que si b$''_\pi$ est satisfait, $f$ fidèle, alors les catégories $i'\backslash I$, dont on sait déjà qu'elles satisfont PF 1$_\pi$, satisfont également PF 2$_\pi$. Soient donc $\xi = (i, u')$, $\eta = (j, v')$ deux objets de $i'\backslash I$, et $(u_\alpha)_{\alpha\in A}$ une famille de flèches $\xi \longrightarrow \eta$ (avec $\card A \leqslant \pi$), donc des flèches dans $I$ telles que $f(u_\alpha)u' = v' \ \forall\, \alpha \in A$.
\[
\begin{tikzcd}[column sep=large]
i \arrow[r, shift left=2, "u_\alpha"] \arrow[r] \arrow[r, shift right=2] & j \arrow[r, dotted, "v"] & k
\end{tikzcd}
\]
\[
\begin{tikzcd}[column sep=large]
i' \arrow[r, "u'"] \arrow[rr, bend right=15, "v'"'] & f(i) \arrow[r, shift left=2] \arrow[r] \arrow[r, shift right=2] & f(j) \arrow[r] & f(k)
\end{tikzcd}
\]

Comme $I'$ est grande devant $\pi$, on peut égaliser par $f(j) \longrightarrow k'$, et comme $k'\backslash I \neq \emptyset$, on peut supposer $k'$ de la forme $f(k)$. Or
\msp{page 242}
d'après l'implication b$''$ $\Longrightarrow$ a $\Longrightarrow$ c de la partie 1°) (cf. page 235) et comme b$''_\pi$ $\Longrightarrow$ b$''$, on sait déjà que $f$ satisfait à F 3. En l'appliquant à $j$, $k$ et $f(j) \longrightarrow f(k)$, on peut
\printedpage{91}
% head: A. Grothendieck : Dérivateurs, Ch. XVIII, Édition des Universités de Montpellier II et Paris VII
% foot: Version 17 décembre 2025
supposer que $f(j) \longrightarrow f(k)$ provient de $v : j \longrightarrow k$. On aura donc $f(v)f(u_\alpha) = f(v)f(u_\beta)$ $\forall\, \alpha, \beta \in A$, donc $f(vu_\alpha) = f(vu_\beta)$. Comme $f$ est fidèle, cela implique que $vu_\alpha = vu_\beta$, i.e. $v$ égalise les $u_\alpha$. D'ailleurs, $v$ peut être considérée comme une flèche \textit{dans} $i'\backslash I$, de $(j, v')$ dans $(k, f(v)v')$. Q.e.d.

\textbf{4.9.6.10.} Quand $I$ est \textit{ordonné}, alors les quatre conditions a$_\pi$, b$_\pi$, c$_\pi$, $b''_\pi$ sont équivalentes, car les $i'\backslash I$ sont ordonnés, et il est claire qu'alors b$''_\pi$ $\Longleftrightarrow$ b$_\pi$. % [sic] claire -> clair, grammar

\msp{page 243}

\dgrblock{dgrb1051}{Replacing faithfulness by eventual equality in its fibres}
\begin{dcomment}{development}{p242}{Replacing faithfulness by eventual equality in its fibres}
The proof isolates a weaker hypothesis than faithfulness.  In approximation problems, two different arrows may become equal after applying the comparison functor, yet already become equal at a later stage of the source.  That is sufficient for the same conclusion.

\textbf{Proposition (relative eventual equalization).}\ours
Let $\pi$ be infinite, $I'$ be $\pi^+$-filtered, and $f:I\to I'$ satisfy F1 and $\mathrm{F3}_\pi$.  Then the following are equivalent: $I$ is $\pi^+$-filtered; and every family $(a_t:x\to y)_{t\in T}$ with $|T|\leq\pi$ and with $f(a_t)$ independent of $t$ is equalized by some arrow $y\to z$ in $I$.  Under these conditions $f$ is final and each $(d\downarrow f)$ is $\pi^+$-filtered.

\textit{Proof.} Necessity follows immediately from filteredness.  Conversely, the preceding simultaneous-lifting argument shows that every comma category is nonempty and has common targets for at most $\pi$ objects.  For a family of at most $\pi$ parallel arrows in such a comma category, equalize their images in $I'$, use F1 to move the target into the image of $f$, and apply F3 to replace the resulting arrow by an image $f(w)$.  The composites $wa_t$ now have identical images.  The asserted eventual-equality condition gives a further arrow $h$ equalizing them in $I$.  It is also an arrow of the comma category when the new structure map is defined by composition.  Thus the comma category is $\pi^+$-filtered.  The cardinal descent proved above gives $\pi^+$-filteredness of $I$, and nonempty connected comma categories give finality.\hfill$\square$

Faithfulness implies the extra condition with $y\to z$ the identity.  A functor from an ordered category has the same property for the still simpler reason that any parallel arrows are already equal.  The proposition describes exactly what the printed proof uses when neither simplification is available.  We checked the cofinality criteria in SGA~4, I, 8.1 and the filtered-colimit criteria in \citet[Section 4.19]{StacksFiltered}; we did not find this particular cardinal formulation with the relative kernel condition stated there.  The argument is a rearrangement of Grothendieck's proof, offered as a usable extension rather than a priority claim.
\end{dcomment}


\section*{4.10\quad Le théorème de représentation indicielle : démonstration}

On va d'abord reformuler le théorème 4.9.3 sous forme plus générale, en termes de catégories \textit{cofibrées}. Soit
\[
\mathcal{D} \longrightarrow J \tag*{(4.10.1)}
\]
un foncteur cofibrant, où $J$ est équivalent à une petite catégorie, et de même pour les catégories fibres $\mathcal{D}_j$ ($j \in \Ob J$). On se donne un cardinal infini $\pi$. Soit
\[
\mathcal{M}_j \ = \ \Ind_\pi(\mathcal{D}_j) \qquad (j \in \Ob J) \ , \tag*{(4.10.2)}
\]
alors les $\mathcal{M}_j$ sont les fibres d'une catégorie cofibrée $\mathcal{N}$ sur $J$, que je désigne aussi par $\Ind_\pi(\mathcal{D}/J)$, donc on pose
\[
\mathcal{N} \ = \ \Ind_\pi(\mathcal{D}/J) \ . \tag*{(4.10.3)}
\]
On a donc une inclusion pleinement fidèle de catégories cofibrées sur $J$
\[
\mathcal{D} \hookrightarrow \mathcal{N} \ ,
\]
d'où une inclusion
\[
\uGamma(\mathcal{D}/J) \longrightarrow \uGamma(\mathcal{N}/J) \ . \tag*{(4.10.4)}
\]

\msp{page 244}
Dans le cas où $\mathcal{D}$ est de la forme $C \times J$, alors $\mathcal{N}$ est de la forme $\mathcal{M} \times J$, avec $\mathcal{M} = \Ind_\pi(C)$, et l'inclusion (4.10.3) n'est autre (à isomorphisme près) que
\[
\underline{\operatorname{Hom}}(J, C) \longrightarrow \underline{\operatorname{Hom}}(J, \mathcal{M}) \ .
\]
Revenant au cas général, on voit encore que :
\begin{enumerate}[label={}, leftmargin=3em]
\item[1°)] Le deuxième membre de (4.10.3) est stable par $\varinjlim$ grandes devant $\pi$, lesquelles se calculent fibre à fibre.
\item[2°)] Le premier membre de (4.10.3) est formé d'objets $\pi$-accessibles du second.
\end{enumerate}
D'où un foncteur \textit{pleinement} fidèle
\[
\Ind_\pi(\uGamma(\mathcal{D}/J)) \hookrightarrow \uGamma(\mathcal{N}/J) \ . \tag*{(4.10.4)}
\]
\setcounter{footnote}{103}%
(\footnote{\texttt{[Les numéros 4.10.1 à 4.10.4 sont utilisés deux fois.]}}).

\dgrblock{dgrb1052}{Sections turn many approximation problems into one}
\begin{dcomment}{exposition}{p243}{Sections turn many approximation problems into one}
Grothendieck now replaces a diagram in one category by a family of categories whose objects are transported along the arrows of the indexing category. The gain is substantial: a single theorem will approximate a diagram, a compatible family of objects, or a section of a fibred construction. His \emph{cofibrant functor} means a Grothendieck opfibration, not a cofibration in a model category. For an arrow $a:i\to j$, choose transport $a_*:\mathcal D_i\to\mathcal D_j$. A section consists of objects $X_i$ and maps $a_*X_i\to X_j$, satisfying unit and composition equations. These maps need not be invertible or cocartesian. The coherence isomorphisms of a cleavage are understood throughout.

To compare his cardinal conventions with current ones, put $\kappa=\pi^+$. A category ``large before $\pi$'' is $\kappa$-filtered: every diagram with at most $\pi$ arrows admits a cocone. An object called $\pi$-accessible here is $\kappa$-presentable. Thus $\Ind_\pi\mathcal D_i$ in these pages is the usual $\Ind_\kappa\mathcal D_i$. This shift matters even when $\pi=\aleph_0$: Grothendieck is then asking for countable-diagram cocones, not merely finite-diagram cocones.

The completion of the opfibration is obtained by applying $\Ind_\kappa$ to every fibre and every transport functor. Its sections may be approximated pointwise without difficulty. The theorem asks whether those approximations can be made into sections before taking their colimit. The condition on one ordinal-shaped colimit supplies exactly the operation needed to turn successive approximations to the structure maps into an actual section.

There is a size qualification to the preliminary assertion on p.~244. Pointwise $\kappa$-presentable sections are $\kappa$-presentable when $|\operatorname{Mor}J|<\kappa$, as in the theorem immediately following. Without it the assertion is false. Take $J$ discrete with $|J|\geq\kappa$ and fibres $\mathrm{Set}=\Ind_\kappa(\mathrm{Set}_{<\kappa})$. The constant singleton family is pointwise presentable. Index a diagram by the $<\kappa$-sized subsets $A\subset J$, putting a singleton in coordinates of $A$ and the empty set elsewhere. This index is $\kappa$-filtered and the colimit is the constant singleton, but no stage receives a map from that family. The presentability claim therefore uses the forthcoming size bound; this does not affect Theorem~4.10.1 itself.
\end{dcomment}


\printedpage{92}
% head: A. Grothendieck : Dérivateurs, Ch. XVIII, Édition des Universités de Montpellier II et Paris VII
% foot: Version 17 décembre 2025

\textbf{Théorème 4.10.1.} \textit{Supposons qu'on ait}
\[
\mathrm{card}\,\textup{réd}\, J \ \leqslant \ \pi \ , \tag*{(4.10.1.1)}
\]
\textit{et qu'il existe un ordinal limite $\alpha$, tel que :}
\begin{enumerate}[label={}, leftmargin=3em]
\item[\textit{1)}] $\card \alpha \leqslant \pi$.
\item[\textit{2)}] \textit{Les limites inductives de}
\msp{page 245}
\textit{type $I_\alpha = \{\beta \mid \beta < \alpha\}$ sont représentables dans les catégories $\mathcal{D}_j$ ($j \in \Ob J$), et pour toute flèche $\alpha : j \longrightarrow j'$ dans $J$ \texttt{[trop de $\alpha$]}, $\alpha^* : \mathcal{D}_j \longrightarrow \mathcal{D}_{j'}$ commute auxdites limites.}
\end{enumerate}
\textit{Alors le foncteur canonique pleinement fidèle (4.10.4) est une équivalence de catégories.}

Ceci va être contenu dans le résultat un peu plus précis :

\textbf{Corollaire 4.10.2.} \textit{Les hypothèses étant celles de 4.10.1, avec en plus}
\[
\card J \ \leqslant \ \pi \tag*{(4.10.2.1)}
\]
\textit{(qui équivaut à (4.10.1.1) si $J$ est réduite, cas auquel on peut évidemment se réduire dans 4.10.1). Soit $\varphi \in \Ob \uGamma(\mathcal{N}/J)$, et considérons le foncteur canonique} fidèle \textit{:}
\[
\uGamma(\mathcal{D}/J)/\varphi \longrightarrow \prod_{j\in J} C_j/\varphi(j) \ . \tag*{(4.10.2.2)}
\]
\textit{Alors la catégorie source $\uGamma(\mathcal{D}/J)/\varphi$ est grande devant $\pi$, et le foncteur envisagé est cofinal.}

\msp{page 246}
De ceci, le théorème résulte de façon essentiellement soritale, tout comme dans le cas particulier 4.9.3 précédent (cf. page 223) (\footnote{Un peu court\,! Il faudrait expliciter l'argument avec soin – il y en a pour deux ou trois pages\,!}).

Pour établir le corollaire, comme la catégorie but de (4.10.2.2) est grande devant $\pi$, comme produit de telles catégories (cf. 4.9.6.6, page 228), et le foncteur fidèle, il suffit en vertu de 4.9.6.3, 2°, de vérifier les conditions F 1 et F 3$_\pi$ (page 225) pour le foncteur envisagé.

On va interpréter le foncteur (4.10.2.2) en termes d'une catégorie cofibrée
\[
\mathcal{D}_\varphi \longrightarrow J \ , \tag*{(4.10.5)}
\]
dont les fibres sont les $\mathcal{D}_j/\varphi(j)$. On note que si $\alpha : j \longrightarrow j'$ est une flèche de $J$, alors on a
\[
\left\{
\begin{aligned}
&\alpha_* : C_j \ \longrightarrow \ C_{j'} \\
&\text{et} \ \ \alpha_*(\varphi(j)) \ \longrightarrow \ \varphi(j') \ ,
\end{aligned}
\right.
\]

\printedpage{93}
% head: A. Grothendieck : Dérivateurs, Ch. XVIII, Édition des Universités de Montpellier II et Paris VII
% foot: Version 17 décembre 2025

\msp{page 247}
d'où un foncteur composé
\[
C_j/\varphi(j) \ \longrightarrow \ C_{j'}/\alpha^*(\varphi(j)) \ \longrightarrow \ C_{j'}/\varphi(j') \ ,
\]
avec transitivité à isomorphisme canonique près etc. On note que l'on a un isomorphisme canonique
\[
\uGamma(\mathcal{D}_\varphi/J) \ \simeq \ \uGamma(\mathcal{D}/J)/\varphi \ ,
\]
de sorte que le foncteur fidèle canonique (4.10.2.2) s'identifie au foncteur
\[
\uGamma(\mathcal{D}_\varphi/J) \longrightarrow \prod_{j\in\Ob J} (\mathcal{D}_\varphi)_j \ . \tag*{(4.10.6)}
\]

\dgrblock{dgrb1053}{From cofinal approximation to the equivalence}
\begin{dcomment}{exposition}{p246}{From cofinal approximation to the equivalence}
Grothendieck explicitly asks that the passage from the cofinality assertion to the representation theorem be written out. We do so because it also explains which information the section category retains. Let $\varphi$ be a section of $\mathcal N=\Ind_\kappa(\mathcal D/J)$ and let
\[
 I=\Gamma(\mathcal D/J)/\varphi.
\]
After replacing the fibres and $J$ by small equivalent categories, this category is small: its objects comprise a small family of section data and maps in locally small fibres. Corollary~4.10.2 makes $I$ $\kappa$-filtered and makes the evaluation functor $I\to\mathcal D_j/\varphi_j$ final for every $j$. Indeed it is the composite of the asserted final functor to the product and a final product projection. The canonical diagram $I\to\Gamma(\mathcal D/J)$ has a cocone to $\varphi$. At $j$ its colimit is
\[
 \mathop{\mathrm{colim}}_{(D\to\varphi)\in I}D_j
 \simeq\mathop{\mathrm{colim}}_{(d\to\varphi_j)\in\mathcal D_j/\varphi_j}d
 \simeq\varphi_j.
\]
These isomorphisms are compatible with every structure map, because the original cocone consists of section morphisms. Filtered colimits of sections are pointwise: the completed transport functors preserve them. Consequently the colimit section is $\varphi$, proving essential surjectivity.

For full faithfulness, let $D$ be a section with values in $\mathcal D_j$. Maps from $D$ to a section $Y$ form the equalizer
\[
 \operatorname{Hom}_{\Gamma}(D,Y)\longrightarrow
 \prod_j\operatorname{Hom}(D_j,Y_j)
 \rightrightarrows
 \prod_{a:i\to j}\operatorname{Hom}(a_*D_i,Y_j).
\]
Every object appearing on the left of these fibrewise hom functors belongs to the corresponding $\mathcal D_j$ and is therefore $\kappa$-presentable. Both products have fewer than $\kappa$ factors. Since $\kappa$-filtered colimits of sets commute with $\kappa$-small limits, this equalizer preserves those colimits. The usual ind hom formula now gives full faithfulness of the comparison functor. This proves the omitted passage without assuming that arbitrary pointwise morphisms are natural.

The auxiliary opfibration $\mathcal D_\varphi$ records the same argument more economically. Its transport sends $(d\to\varphi_i)$ to $a_*d\to a_*\varphi_i\to\varphi_j$, and a section of it is precisely a section $D$ together with $D\to\varphi$. A colimit of an ordinal chain in such a slice is computed in $\mathcal D_j$, with its induced map to $\varphi_j$. Lemma~4.10.4 ensures that this remains the same colimit after embedding in the ind-category. Thus no extra colimit hypothesis is lost in passing to the slices.
\end{dcomment}


Le corollaire 4.10.2 sera alors une conséquence du résultat plus général suivant :

\textbf{Corollaire 4.10.3.} \textit{Soit $\mathcal{E} \longrightarrow J$ un foncteur cofibrant, et $\pi$ un cardinal infini. On suppose :}
\begin{enumerate}[label={}, leftmargin=3em]
\item[\textit{a)}] $\card J \leqslant \pi$.
\item[\textit{b)}] \textit{Les catégories fibres $\mathcal{E}_j$ sont équivalentes à des petites catégories, et grandes devant $\pi$.}
\item[\textit{c)}] \textit{Il existe un ordinal limite tel que :}
\begin{enumerate}[label={}, leftmargin=3em]
\item[\textit{1°)}] $\card \alpha \leqslant \pi$.
\item[\textit{2°)}] \msp{page 248}\textit{Les limites inductives de type $I_\alpha$ sont représentables dans les catégories $\mathcal{E}_j$, et les foncteurs cochangement de base $\lambda_* : \mathcal{E}_j \longrightarrow \mathcal{E}_{j'}$ y commutent. (C'est le cas p.ex. si les limites inductives filtrantes} dénombrables \textit{sont représentables dans les $\mathcal{E}_j$, et si les $\alpha_*$ y commutent\,; ou s'il existe un cardinal $\pi_0 \leqslant \pi$ et un cardinal $\pi'_0 > \pi_0$, tel que les limites inductives grandes devant $\pi_0$ et de cardinal $\leqslant \pi'_0$ soient représentables dans les $\mathcal{E}_j$, et que les $\alpha_*$ y commutent.)}
\end{enumerate}
\end{enumerate}
\textit{Sous ces conditions, le foncteur canonique fidèle}
\[
\uGamma(\mathcal{E}/J) \longrightarrow \prod_{j\in\Ob J} \mathcal{E}_j \tag*{(4.10.3.1)}
\]
\textit{est cofinal, et la catégorie source est grande devant $\pi$.}

Il est $\pm$ immédiat que les hypothèses de 4.10.1 impliquent que celles de 4.10.3 sont satisfaites pour $\mathcal{E} = \mathcal{D}_\varphi \longrightarrow J$ (4.10.5). Pour
\msp{page 249}
ceci, on utilise le

% marginal vertical bar printed in the left margin beside the following statement
\textbf{Lemme 4.10.4.} \textit{Soient $C$ une catégorie, et $\pi$ un cardinal infini. Alors le foncteur d'inclusion $C \longrightarrow \mathcal{M} \overset{\textup{déf}}{=} \Ind_\pi C$ commute aux $\varinjlim$ de cardinal $\leqslant \pi$.}

Soit en effet $I$ telle que $\card I \leqslant \pi$, et $(x_i)_{i\in I}$ un système inductif sur $C$, y ayant une limite $x$. Soit d'autre part $X = \varinjlim\limits_{I'} \xi_{i'}$ un objet de $\Ind_\pi C$, avec $I'$ grande devant $\pi$. On a

% ---- from der-ch18_p096-100.tex
\printedpage{94}
% head: A. Grothendieck : Dérivateurs, Ch. XVIII, Édition des Universités de Montpellier II et Paris VII
donc
\[
\begin{aligned}
\Hom_{\mathcal{M}}(x,X) &\simeq \varinjlim_{I'} \Hom(x,\xi_{i'})\\
&\simeq \varinjlim_{I'} \varprojlim_{I} \Hom(x_i,\xi_{i'})\\
&\overset{\text{SGA 4 I 9.8}}{\simeq} \varprojlim_{I} \underbrace{\varinjlim_{I'} \Hom(x_i,\xi_{i'})}_{\simeq \Hom(x_i,X)}\\
&\simeq \varprojlim_{I} \Hom(x_i,X)\ ,
\end{aligned}
\]
ce qui prouve que $x$ est aussi une $\varinjlim$ des $x_i$ dans $\mathcal{M}$ tout entier.

\msp{250}
Ainsi, on est réduit à prouver 4.10.3. Appliquant 4.9.6 2°, on est ramené à prouver, pour le foncteur envisagé, les conditions F 1 et F $3_\pi$ (page 225).

Dans les arguments qui suivent, on va utiliser les conditions a), b) de 4.10.3, mais non, pour l'instant, la condition c.

\textbf{Lemme 4.10.5.} \textit{Sous les conditions de 4.10.3 a) b) (à l'exclusion de c), pour que (4.10.3.1) soit cofinal et $\Gamma(\mathcal{E}/J)$ grande devant $\pi$, il faut et il suffit que le foncteur envisagé satisfasse à la condition F 1, rappelée ci-dessous, et à la variante F$'$ 1 de celle-ci énoncée ci-dessous :}
\begin{itemize}[leftmargin=3em]
\item[F\,1)] {\itshape Pour tout système $(X(j))_{j\in\Ob J}$ d'éléments $X(j)\in\Ob\mathcal{E}_j$, existe une \textup{section} $G$ de $\mathcal{E}$ sur $J$, et un système de flèches
\[
X(j) \longrightarrow G(j) \qquad (j\in ObJ)\ . \tag*{\textup{(4.10.5.1)}}
\]}
\item[F$'$1)] {\itshape Sous les conditions de \textup{F\,1} précédent, supposons de plus une section $F$ de $\mathcal{E}$ sur $J$ \msp{251} et un système de flèches
\[
F(j) \longrightarrow X(j) \qquad (j\in\Ob J)\ ,
\]
alors $G$ et le système de flèches (4.10.5.1) peut être choisi de telle façon que les flèches composées
\[
F(j) \longrightarrow X(j) \longrightarrow G(j)
\]
définissent un homomorphisme de sections
\[
F \longrightarrow G
\]
(dans la catégorie $\uGamma(\mathcal{E}/J)$).}
\end{itemize}
(\footnote[106]{\textbf{N.B.} Les conditions F 1 et F$\,'$\,1 équivalent à dire que le foncteur envisagé satisfait aux conditions F1 et F3, vu encore (en vertu de 4.9.6.3, 1°) que $\uGamma(\mathcal{E}/J)$ est \textit{filtrante} sans plus, et le foncteur fidèle (4.10.3.1) cofinal. Cela implique déjà que $\uGamma(\mathcal{E}/J)$ est grande devant $\pi$.}).

\dgrblock{dgrb1054}{Why one compatible family controls many families}
\begin{dcomment}{exposition}{p250}{Why one compatible family controls many families}
The force of Lemma~4.10.5 is that a construction with one prescribed section suffices to establish the simultaneous lifting condition for at most $\pi$ sections. The fibres supply enough room to remove all failures of naturality before the section construction starts.

Suppose sections $F_a$ and $F$ are given, together with arbitrary component maps $u_{a,j}:F_{a,j}\to F_j$, where $|A|\leq\pi$. For an arrow $b:i\to j$, compare the two maps from $b_*F_{a,i}$ to $F_j$: one uses the structure of $F_a$ followed by $u_{a,j}$, the other uses $b_*u_{a,i}$ followed by the structure of $F$. There are at most $\pi$ such pairs at each $j$. Since $\mathcal E_j$ is $\kappa$-filtered, they can all be equalized by one map $F_j\to X_j$. This use of transport is implicit in the source's arrows in the total category.

Apply F$'1$ to the section $F$ and the maps $F_j\to X_j$. It produces a section $G$ and maps $X_j\to G_j$ for which $F\to G$ is natural. Each $F_a\to G$ is now natural as well: its naturality square is the composite of the already equalized square ending at $X_j$ and the naturality square for $F\to G$. Thus F$3_\pi$ holds. The faithful cofinality criterion gives both $\kappa$-filteredness of $\Gamma(\mathcal E/J)$ and finality of evaluation.

Conversely, suppose evaluation is final and its source is $\kappa$-filtered. F1 follows. Given $F_j\to X_j$, apply F1 to obtain $X_j\to H_j$ for some section $H$. The ordinary F3 condition for the resulting pointwise arrow $F\to H$ then gives a section map $H\to G$ whose composite from $F$ is natural. This is F$'1$. The equivalence asserted in the lemma is therefore exact; it is special to this evaluation functor and the bound on the number of naturality equations.
\end{dcomment}


\textsc{Démonstration}. Il faut voir que F 1 \& F$'$ 1 équivaut à F 1 \& F $3_\pi$. L'implication $\Longleftarrow$ est $\pm$ immédiate, c'est de $\Longrightarrow$ que nous aurons besoin, donc
\[
\text{F 1 \& F}'\text{ 1} \quad\Longrightarrow\quad \text{F } 3_\pi\ .
\]
% foot: Version 17 décembre 2025

\printedpage{95}
% head: A. Grothendieck : Dérivateurs, Ch. XVIII, Édition des Universités de Montpellier II et Paris VII
Pour ceci, réénonçons F $3_\pi$ dans le contexte présent. On se donne une famille $(F_\alpha)_{\alpha\in A}$, $\card A\leqslant\pi$, d'objets de $\uGamma(\mathcal{E}/J)$, un objet $F$ de la même catégorie, et pour tout $j$ dans $\Ob J$, un système de flèches
\[
F_\alpha(j) \xrightarrow{u_\alpha(j)} F(j)\ ,
\]
mais sans supposer que pour $\alpha$ fixé, $j$ variable

\msp{252}
cela définisse un homomorphisme de sections $F_\alpha\longrightarrow F$. On cherche une section $G$ de $\mathcal{E}$ sur $J$, et un homomorphisme de sections $F\longrightarrow G$, tels que pour tout $\alpha$, le système des composés
\[
F_\alpha(j) \longrightarrow F(j) \longrightarrow G(j)
\]
définisse un homomorphisme de sections $F_\alpha\longrightarrow G_\alpha$. Soit donc $\lambda : j\longrightarrow j'$ une flèche de $J$, considérons les carrés (\textit{non commutatifs}) % [sic] F_\alpha -> G_\alpha -> F_\alpha -> G, the section sought is G (cf. p. 95 below: F_\alpha -> G)
\[
\begin{tikzcd}
F_\alpha(j) \arrow[r,"u_\alpha(j)"] \arrow[d,"F_\alpha(\lambda)"'] \arrow[dr,phantom,"!"] & F(j) \arrow[d,"F(\lambda)"] \\
F_\alpha(j') \arrow[r,"u_\alpha(j')"] & F(j')
\end{tikzcd}
\tag*{$(*)$}
\]
(indexés par $\alpha\in A$). Pour $j'$ fixé dans $J$, le nombre de ces carrés est majoré par $\card\Ob(J/j')\times\card A\leqslant\pi\cdot\pi=\pi$. Il s'ensuit, comme $\mathcal{E}_{j'}$ est grande devant $\pi$, qu'il existe un objet $X(j')$ de $\mathcal{E}_{j'}$, et une flèche $v(j') : F(j')\longrightarrow X(j')$ de telle façon que pour tout $\alpha\in A$, et tout $\lambda\in\Ob J/j'$, le carré déduit de $(*)$

\msp{253}
en composant avec $F(j')\longrightarrow X(j')$ soit commutatif. On fait ce choix pour tout objet $j'$ de $J$. Appliquons alors F$'$ 1 à $F$ et au système de flèches $F(j)\xrightarrow{v(j)}X(j)$, on trouve $G$ et un système de flèches
\[
w(j)\ :\ X(j) \longrightarrow G(j)
\]
de telle façon que les flèches composées
\[
\begin{tikzcd}[column sep=small]
w(j)v(j) : & F(j) \arrow[rr] \arrow[dr,dotted] & & G(j) \\
& & X(j) \arrow[ur,dotted] &
\end{tikzcd}
\]
définissent un homomorphisme de sections. Je dis que cet homomorphisme a la propriété voulue pour F $3_\pi$, i.e. que pour tout $\alpha\in A$, le système des composés
\[
F_\alpha(j) \longrightarrow F(j) \longrightarrow G(j)
\]
(pour $j\in\Ob J$ variable) définit un homomorphisme de sections $F_\alpha\longrightarrow G$. En effet, il suffit pour cela de considérer, pour une flèche $\lambda$ de $J$, le
% foot: Version 17 décembre 2025
\printedpage{96}
% head: A. Grothendieck : Dérivateurs, Ch. XVIII, Édition des Universités de Montpellier II et Paris VII
\msp{254}
diagramme
\[
\begin{tikzcd}
F_\alpha(j) \arrow[r,"u_\alpha(j)"] \arrow[d,"F_\alpha(\lambda)"'] \arrow[dr,phantom,"!"] & F(j) \arrow[r,"v(j)"] \arrow[d,"F(\lambda)"'] & X(j) \arrow[r,"w(j)"] \arrow[d,phantom,"\text{comm.}"] & G(j) \arrow[d,"G(\lambda)"] \\
F_\alpha(j') \arrow[r,"u_\alpha(j')"] & F(j') \arrow[r,"v(j')"] & X(j') \arrow[r,"w(j')"] & G(j')\ .
\end{tikzcd}
\]
Le lemme 4.10.5 étant prouvé, tout revient donc à prouver la condition F 1, et son complément F$'$ 1, qui se prouvera dans la même haleine.

\textbf{4.10.6.} On se donne donc le système d'objets
\[
X(j)\ \in\ \Ob\mathcal{E}_j \qquad (j\in\Ob J)\ , \tag*{(4.10.6.1)}
\]
et le cas échéant, un objet $F$ de $\uGamma(\mathcal{E}/J)$, et un système de flèches
\[
F(j) \xrightarrow{u_j} X(j)\ , \tag*{(4.6.10.2)} % [sic] (4.6.10.2) -> (4.10.6.2), numbering of 4.10.6; cited as (4.10.6.2) below
\]
et on se propose de construire $G$ de $\uGamma(\mathcal{E}/J)$, et un système de flèches
\[
X(j) \xrightarrow{w(j)} G(j)\ , \tag*{(4.10.6.3)}
\]
de telle façon que (si on a la donnée supplémentaire des $u(j)$) les flèches composées
\[
F(j) \xrightarrow{u(j)} X(j) \xrightarrow{v(j)} G(j) % [sic] v(j) -> w(j), the arrow X(j) -> G(j) is w(j) in (4.10.6.3)
\]

\msp{255}
soient fonctorielles en $j$.

On va faire une construction transfinie sur les ordinaux $\alpha$ tels que
\[
\card\alpha\ \leqslant\ \pi
\]
d'un système inductif
\[
(X_\alpha(j))_\alpha
\]
d'objets de $\prod_{j\in\Ob J}\mathcal{E}_j$. On va d'abord ignorer la donnée éventuelle (4.10.6.2). On posera
\[
X_0(j)\ =\ X(j)\ , \tag*{(4.10.6.4)}
\]
et on va donner la prescription pour passer de $X_\alpha$ à $X_{\alpha+1}$. La construction inclura aussi, pour toute flèche $\lambda : j'\longrightarrow j$ de $J$ et tout $\alpha$, la donnée de
\[
X_\alpha(j) \xrightarrow{t_\alpha(\lambda)} X_{\alpha+1}(j)\ , \qquad \text{flèche (dans $\mathcal{E}$) au dessus de $\lambda$,} \tag*{(4.10.6.5)} % [sic] X_\alpha(j) -> X_\alpha(j'), the source of an arrow over \lambda : j' -> j (cf. p. 97: X_\alpha(j') -> X_{\alpha+1}(j))
\]
et ces flèches, et les flèches de transition $X_\alpha(j)\xrightarrow{\tau_\alpha(j)}X_{\alpha+1}(j)$, doivent
% foot: Version 17 décembre 2025
\printedpage{97}
% head: A. Grothendieck : Dérivateurs, Ch. XVIII, Édition des Universités de Montpellier II et Paris VII
\msp{256}
satisfaire aux trois conditions suivantes :
\begin{itemize}[leftmargin=3em]
\item[(i)] Si $\lambda$ est une flèche identique $\mathrm{id}_J : J\longrightarrow J$, alors \fbox{$t_\alpha(\lambda)=\tau_\alpha(j)$.} % [sic] id_J : J -> J -> id_j : j -> j, the identity arrow of an object j of J
\item[(ii)] Dans le diagramme
\[
\begin{tikzcd}
X_\alpha(j') \arrow[r,"\tau_\alpha(j')"] \arrow[dr,"t_\alpha(\lambda)"] & X_{\alpha+1}(j') \arrow[r] \arrow[dr,"t_{\alpha+1}(\lambda)"] \arrow[d,phantom,"\bullet"] & X_{\alpha+2}(j') \arrow[r] \arrow[dr] & X_{\alpha+3}(j') \arrow[r] \arrow[dr,dotted] & {} \\
X_\alpha(j') \arrow[r] & X_{\alpha+1}(j') \arrow[r,"\tau_{\alpha+1}(j')"] & X_{\alpha+2}(j') \arrow[r,"\tau_{\alpha+2}(j')"] & X_{\alpha+3}(j') \arrow[r] & {,}
\end{tikzcd}
\] % [sic] bottom row printed with j' (X_\alpha(j'), X_{\alpha+1}(j'), \tau_{\alpha+1}(j'), X_{\alpha+2}(j'), \tau_{\alpha+2}(j'), X_{\alpha+3}(j')) -> j, the arrows t(\lambda) go to the objects at j (cf. the text below: \tau_{\alpha+1}(j), \tau_{\alpha+2}(j))
$\tau_{\alpha+2}(j)$ “rend le carré $\bullet$ commutatif”, i.e. égalise $t_{\alpha+1}(\lambda)\tau_\alpha(j')$ et $\tau_{\alpha+1}(j)t_\alpha(\lambda)$ :
\[
\boxed{\tau_{\alpha+2}(j)t_{\alpha+1}(\lambda)\tau_\alpha(j')\ =\ \tau_{\alpha+2}(j)\tau_{\alpha+1}(j)t_\alpha(\lambda)\ .}
\]
\item[(iii)] Si on a deux flèches consécutives $\lambda : j''\longrightarrow j'$, $\mu : j'\longrightarrow j$, on a donc le diagramme
\[
\begin{tikzcd}
X_\alpha(j'') \arrow[r] \arrow[dr,"t_\alpha(\lambda)"] \arrow[ddr,"t_\alpha(\mu\lambda)"] & X_{\alpha+1}(j'') \arrow[r] \arrow[dr] & X_{\alpha+2}(j'') \arrow[r] \arrow[dr] & X_{\alpha+3}(j'') \\
X_\alpha(j') \arrow[r] & X_{\alpha+1}(j') \arrow[r] \arrow[dr,"t_{\alpha+1}(\mu)"] & X_{\alpha+2}(j') \arrow[r] \arrow[dr] & X_{\alpha+3}(j') \\
X_\alpha(j) \arrow[r] & X_{\alpha+1}(j) \arrow[r,"\tau_{\alpha+1}(j)"] & X_{\alpha+2}(j) \arrow[r,"\tau_{\alpha+2}(j)"] & X_{\alpha+3}(j)\ ,
\end{tikzcd}
\]
\[
\boxed{\tau_{\alpha+2}(j)t_{\alpha+1}(\mu)t_\alpha(\lambda)\ =\ \tau_{\alpha+2}(j)\tau_{\alpha+1}(j)t_\alpha(\mu\lambda)}
\]
\end{itemize}

\msp{257}
Ainsi, pour définir les $X_{\alpha+1}(j)$, il faut tenir compte de toutes ces contraintes. Donc :
\begin{itemize}[leftmargin=3em]
\item[1°)] Pour $j$ fixé, $X_{\alpha+1}(j)$ doit être but des flèches
\[
X_\alpha(j') \xrightarrow{t_\alpha(\lambda)} X_{\alpha+1}(j)\ ,
\]
pour les $\lambda\ :\ j'\longrightarrow j$ dans $\Ob J/j$ -- y compris pour $\lambda=\mathrm{id}_j$, auquel cas $t_\alpha(\lambda)\ :\ X_\alpha(j)\longrightarrow X_{\alpha+1}(j)$ jouera le rôle du morphisme de transition $\tau_\alpha(j)$.
\item[2°)] Si $\alpha$ est de la forme $\alpha'+2$, i.e. s'il n'est ni ordinal limite, ni successeur d'un ordinal limite, alors on exige de plus, pour les morphismes de transition $\tau_\alpha(j)$, d'égaliser les couples de flèches prévus dans (ii) et (iii) ci-dessus.
\end{itemize}
Comme l'ensemble de ces données dans 1°), ou de ces contraintes dans 2°),
\msp{258}
est de cardinal $\leqslant\pi$, et que la catégorie $\mathcal{E}_j$ est grande devant $\pi$, on voit qu'on peut bel et bien trouver $X_{\alpha+1}(j)$ et des $t_\alpha(\lambda)\ :\ X_\alpha(j')\longrightarrow X_{\alpha+1}(j)$, de façon à satisfaire (si $\alpha=\alpha'{+}2$) aux contraintes rappelées dans 2°. Cela achève la description de la construction de $X_{\alpha+1}(\bullet)$, en termes de $X_\beta(\bullet)$ $(\beta\leqslant\alpha)$ et des $t_\beta(\lambda)$ $(\beta<\alpha)$.

Il faut enfin définir $X_\alpha(\bullet)$ pour $\alpha$ un ordinal limite. Cette fois, on lui demande seulement de recevoir les $X_\beta(\bullet)$ pour $\beta<\alpha$. Comme l'ensemble $I_\alpha$ de ces $\beta$ est de cardinal $\leqslant\pi$, et que les $\mathcal{E}_j$ sont grandes devant $\pi$, on peut bel et bien trouver les $X_\alpha(j)$.
% foot: Version 17 décembre 2025

\printedpage{98}
% head: A. Grothendieck : Dérivateurs, Ch. XVIII, Édition des Universités de Montpellier II et Paris VII
\msp{259}
On a donc prouvé l'essentiel du

\dgrblock{dgrb1055}{A transfinite compatibility condition must be added}
\begin{dcomment}{development}{p258}{A transfinite compatibility condition must be added}
Grothendieck's iterative construction has the right mechanism: add all needed arrows, and then move farther along the filtered system until the equations hold. At an internal limit ordinal, however, the displayed conditions on consecutive stages no longer compare the new arrows with all earlier ones. The inference on p.~262 from condition (ii) to a morphism of ind-objects needs a repair.

\textbf{A counterexample to that inference.}\ours\ Let $J=B(\mathbb Z/2)$ and take the constant opfibration with fibre a small skeleton of countable sets. This fibre has a terminal object, hence is $\aleph_1$-filtered, and it admits countable ordinal colimits. Set $X_\beta=\{0,1\}$ with identity transition maps for every countable ordinal $\beta$. For the nonidentity group element $s$, set $t_\beta(s)=\mathrm{id}$ when $\beta<\omega$, and set $t_\beta(s)$ equal to the transposition when $\beta\geq\omega$. Put $t_\beta(1)=\mathrm{id}$ throughout. Conditions (i)--(iii) hold: consecutive indices never cross from a finite ordinal to $\omega$, and on either side the maps form a representation of $\mathbb Z/2$. But along the chain of length $\omega+\omega$ the induced maps from its constant stages would have to be both the identity and the transposition. They define no morphism of its ind-objects. This disproves the intermediate inference, not the section theorem.

\textbf{Strengthened construction.}\ours\ Under 4.10.3(a),(b), fix a nonzero limit ordinal $\delta$ with $|\delta|\leq\pi$. Write $r_{\gamma\beta}:X_\beta\to X_\gamma$ for the transitions, and write a map above $a:i\to j$ as $t_\beta(a):a_*X_\beta(i)\to X_{\beta+1}(j)$. The objects and maps can be chosen so that, for every $\beta<\gamma<\delta$,
\[
 t_\gamma(a)\,a_*r_{\gamma\beta}
 =r_{\gamma+1,\beta+1}\,t_\beta(a),                       \tag{B1}
\]
and so that the unit condition holds and the composition equations become equal after a later transition. Prescribed maps $F_j\to X_0(j)$ from a section can also be made natural after the first step.

Here is the full induction. Having constructed the system through stage $\gamma$, first choose in each fibre a common target $Y_j$ of $X_\gamma(j)$ and all $a_*X_\gamma(i)$ for arrows $a:i\to j$. This uses at most $\pi$ objects. For every $\beta<\gamma$, the proposed two sides of (B1) are now parallel maps into $Y_j$. Postcompose the chosen maps by one equalizing arrow, simultaneously for every such pair. There are at most $\pi\cdot\pi=\pi$ pairs.

At this same step impose all previously formed composition equations. For $a:i\to j$, $b:j\to k$ and $\beta+2\leq\gamma$, compare, in $\mathcal E_k$, the maps obtained from
\[
 t_{\beta+1}(b)\,b_*t_\beta(a)
 \quad\hbox{and}\quad t_\beta(ba)
\]
by transport to $X_\gamma(k)$ and then to $Y_k$. Again there are at most $\pi$ pairs; add them to the equalizations already required. Coherence isomorphisms identify $(ba)_*$ with $b_*a_*$. Choose the identity transport map to be the transition itself. At the first step also equalize $t_0(a)a_*u_i$ and $r_{10}(j)u_jF(a)$ for every $a$, where $u_j:F_j\to X_0(j)$ is the optional datum. Filteredness supplies a single further target accommodating all these equations. Call it $X_{\gamma+1}(j)$.

At a limit stage $\gamma<\delta$, choose a cocone on the entire earlier chain in each fibre; this defines $X_\gamma$ and the transitions $r_{\gamma\beta}$. Its diagram has at most $\pi$ arrows. No colimit is required at this stage. At the next successor step (B1) is imposed for every earlier $\beta$, including those before this limit. This is precisely the additional condition absent from the printed construction. All equations already imposed survive postcomposition, so the induction is consistent.

Suppose now that the $\delta$-chain colimits exist and transport preserves them. Put $G_j=\operatorname{colim}_{\beta<\delta}X_\beta(j)$. Equation (B1) makes the maps $a_*X_\beta(i)\to G_j$ a compatible cocone; preservation gives $a_*G_i\to G_j$. The composition equations hold on each stage after some further stage, and hence hold in $G$. The unit equations hold by construction. Thus $G$ is a section receiving $X_0$, and the first-step equations make $F\to G$ natural when prescribed. This proves F1 and F$'1$, completing the general ordinal case of the theorem. For $\delta=\omega$, the source's consecutive-stage argument already suffices.

The parenthesis in 4.10.3(c) has a related endpoint issue. If $\pi_0=\pi$, its hypotheses need not supply an ordinal of cardinal at most $\pi$. They still suffice for the conclusion: take $\delta=\pi_0^+$, which is at most $\pi'_0$ and is large before $\pi_0$. Every proper initial segment has cardinal at most $\pi$, so the induction above runs through all stages below $\delta$; only the final colimit has the larger size. Alternatively, imposing $\pi_0<\pi$ makes the printed implication to condition (c) valid as stated.

We checked the related constructions in SGA~4, I, 9.21 and the modern limit results cited below. We did not find this particular defect in the displayed XVIII argument discussed there. The strengthened induction is supplied as a proof repair, with no claim of priority.
\end{dcomment}


\textbf{Lemme 4.10.7.} \textit{Sous les conditions de 4.10.3 a) b), et pour un système $X(\bullet)\in\Ob\prod_{j\in\Ob J}\mathcal{E}_j$ donné, on peut trouver un système inductif $(X_\alpha(i))_\alpha$ dans $\prod\mathcal{E}_j$, indexé par les ordinaux $\alpha$ telle que} $\card\alpha\leqslant\pi$, \textit{et pour toute flèche $\lambda : j'\longrightarrow j$ dans $J$ des flèches $t_\alpha(\lambda) : X_\alpha(j')\longrightarrow X_\alpha(j)$, de telle façon que l'on ait les conditions suivantes.} % [sic] (X_\alpha(i))_\alpha -> (X_\alpha(j))_\alpha, the index is j \in Ob J; [sic] X_\alpha(j') -> X_\alpha(j) -> X_\alpha(j') -> X_{\alpha+1}(j), cf. (4.10.6.5) and p. 97
\begin{itemize}[leftmargin=3em]
\item[\textit{1°)}] \textit{$X_0(\bullet)=X(\bullet)$.}
\item[\textit{2°)}] \textit{Les conditions (i) (ii) (iii) (p. 256).}
\item[\textit{3°)}] \textit{Si on se donne $F$ dans \textup{Ob}$\,\uGamma(\mathcal{E}/J)$, et des flèches $u(j) : F(j)\longrightarrow X(j)=X_0(j)$, on peut supposer que pour toute flèche $\lambda : j'\longrightarrow j$ dans $J$, on a}
\[
\begin{tikzcd}
F(j') \arrow[r,"u(j')"] \arrow[d,"F(\lambda)"'] & X_0(j') \arrow[dr,"t_0(\lambda)"] & \\
F(j) \arrow[r,"u(j)"] & X_0(j) \arrow[r,"\tau_0(j)"] & X_1(j)
\end{tikzcd}
\qquad\qquad \text{commutatif.}
\tag*{$(*)$}
\]
\end{itemize}

\msp{260}
Il reste seulement à prouver ce dernier complément 3°. C'est une contrainte uniquement sur la construction de $X_1(\bullet)$ et des $t_0(\lambda)$. Supposons donc que l'on ait construit $X_1(\bullet)$ sans tenir compte de cette contrainte. Pour $j$ fixé, il faut donc “rendre commutatif” tous ces carrés, i.e. toutes ces double-flèches $(t_0(\lambda)u(j'),\tau_0(j)u(j)F(\lambda))$ de but $X_1(j)$. Comme l'ensemble de ces double-flèches est de cardinal $\card\Ob\mathcal{E}/j\leqslant\pi$, on peut trouver $X_1(j)\longrightarrow X'_1(j)$ dans $\mathcal{E}_j$, égalisant toutes ces double-flèches, donc quitte à remplacer $X_1(j)$ par $X'_1(j)$, on gagne.

\msp{261}
Il est temps enfin de faire appel à l'hypothèse c) de 4.10.3. Soit donc $\alpha$ un ordinal limite, comme stipulé dans cette condition. Comme $\card\alpha\leqslant\pi$, a fortiori $\card\beta\leqslant\pi$ si $\beta<\alpha$, on a donc un système inductif
\[
(X_\beta(\bullet))_{\beta\in I_\alpha}
\]
dans $\prod_{j\in J}\mathcal{E}_j$. Comme les $\varinjlim$ de type $I_\alpha$ existent dans $\mathcal{E}_j$, on peut donc considérer
\[
G(j)\ =\ \varinjlim_{\beta\in I_\alpha} X_\beta(j)\ \in\ \Ob\mathcal{E}_j\ .
\]
Si $\lambda : j'\longrightarrow j$ est une flèche de $J$, on en déduit une $\lambda$-flèche
\[
G(\lambda)\ :\ G(j)=\varinjlim X_\beta(j) \longrightarrow G(j')\simeq\varinjlim X_{\beta+1}(j')
\]
par $G(\lambda)=\varinjlim t_\beta(\lambda)$. Il faut
\msp{262}
faire attention que les $t_\beta(\lambda)$ ne définissent pas un homomorphisme de systèmes projectifs, mais par la condition (ii), p. 256, ils définissent un homomorphisme d'ind-objets, donc
% foot: Version 17 décembre 2025

% ---- from der-ch18_p101-105.tex
\printedpage{99}
% head: A. Grothendieck : Dérivateurs, Ch. XVIII, Édition des Universités de Montpellier II et Paris VII
% foot: Version 17 décembre 2025

un homomorphisme sur les $\varinjlim$ (\footnote{Attention, c'est ici qu'on utilise tacitement le fait que le foncteur $\lambda_* : \mathcal{E}_{j'} \longrightarrow \mathcal{E}_j$ commute aux $\varinjlim$ de type $I_\alpha$, ce qui implique que les inclusions $\mathcal{E}_j \into \mathcal{E}$ commutent à ce type de $\varinjlim$ .}). La condition (i) assure que $\lambda \longmapsto G(\lambda)$ transforme identités en identités, et la condition (iii) que l'on a $G(\mu\lambda) = G(\mu)G(\lambda)$. Donc $G$ apparaît comme une section de $\mathcal{E}$ sur $J$. On a d'autre part des flèches
\[
X(j) \defeq X_0(j) \longrightarrow G(j) = \varinjlim_{\beta\in I_\alpha} X_\beta(j) \ ,
\]
ce qui prouve F (1). D'autre part, la condition 3°) du lemme prouve la validité de F$'$ (1). Cela achève la démonstration de 4.10.3, et par là de 4.10.1 et de 4.9.3.

\msp{page 263}

\dgrblock{dgrb1056}{A section is a module with an extra structure map}
\begin{dcomment}{examples}{p262}{A section is a module with an extra structure map}
The section theorem becomes concrete over the walking arrow. Let $R\to S$ be a homomorphism of rings and let its transport functor be $S\otimes_R-:R\text{-}\mathrm{Mod}\to S\text{-}\mathrm{Mod}$. A section is a triple $(M,N,\theta)$ with $\theta:S\otimes_RM\to N$. A map of sections is a pair $(f,g)$ satisfying $g\theta=\theta'(1\otimes f)$. Thus the condition imposed by the section category is one familiar commutative square.

This category is equivalent to the category of left modules over
\[
 T=\begin{pmatrix}R&0\\ S&S\end{pmatrix},
\]
where the lower-left copy of $S$ is the $(S,R)$-bimodule supplied by the ring map. The module corresponding to $(M,N,\theta)$ is $M\oplus N$, with action
\[
 \begin{pmatrix}r&0\\s&t\end{pmatrix}(m,n)
   =(rm,\theta(s\otimes m)+tn).
\]
Bilinearity of $\theta$ makes this action associative; conversely the two diagonal idempotents of $T$ split any $T$-module into $M$ and $N$, and its lower-left action recovers $\theta$. These constructions are inverse on objects and maps.

Choose an infinite $\pi\geq|R|,|S|$ and put $\kappa=\pi^+$. A triple is the union of its subtriples whose two underlying sets have cardinal at most $\pi$. To construct one containing specified subsets of at most $\pi$ elements, take the $R$-submodule $M_0$ they generate, then the $S$-submodule $N_0$ generated by the prescribed elements and the image of $S\otimes_RM_0$. Both sets still have cardinal at most $\pi$, since $\pi$ is infinite. Any collection of at most $\pi$ such subtriples is contained in another, by the same construction. Their poset is therefore $\kappa$-filtered and their colimit is the original triple.

Each small triple is $\kappa$-presentable. Indeed its $T$-module has at most $\pi$ elements, hence a presentation with at most $\pi$ generators and relations: use all its elements as generators and all relations in the free module, whose underlying set also has cardinal at most $\pi$. A map from this presentation into a $\kappa$-filtered colimit factors at one stage, and its relations and the equality of two such maps can be witnessed at one later stage. This directly verifies the presentability assertion. The construction illustrates both tasks in Grothendieck's proof: enlarge the objects enough to contain the structure map, then impose its equations.
\end{dcomment}


\textbf{Corollaire 4.10.8} (du théorème 4.10.1). \textit{Soit $\mathcal{N} \longrightarrow J$ un foncteur cofibrant, tel que les fibres soient des catégories accessibles (4.8), les foncteurs cochangement de base $\lambda_*$ : $\mathcal{N}_{j'} \longrightarrow \mathcal{N}_j$ accessibles, enfin $J$ équivalente à une petite catégorie. Alors la catégorie $\mathcal{M} = \uGamma(\mathcal{N}/J)$ est accessible. Plus précisément, soit $\pi_0$ un cardinal infini tel que les catégories $\mathcal{N}_j$ soient $\pi_0$-accessibles, et que les foncteurs $\lambda_*$ soient $\pi_0$-accessibles, et soit $\pi$ un cardinal infini tel que l'on ait :}
\begin{enumerate}[leftmargin=*]
\item[\cercle{a}] $\pi^{\pi_0} = \pi$.
\item[\cercle{b}] $\forall\ \lambda : j' \longrightarrow j$ \textit{dans} $J$, $\lambda_* : \mathcal{E}_j \longrightarrow \mathcal{E}_{j'}$ \textit{applique} $(\mathcal{E}_j)_{\pi_0}$ \textit{dans} $(\mathcal{E}_{j'})_{\pi_0}$.
\item[\cercle{c}] $\cardred J \leqslant \pi$.
\end{enumerate}
\textit{Alors $\mathcal{M}$ est $\pi$-accessible, et plus précisément, on a}
\[
\mathcal{M} \ \simeq \ \Ind_\pi(\mathcal{M}_\pi) \ ,
\]
\textit{où $\mathcal{M}_\pi$ est la sous-catégorie pleine de $\mathcal{M}$ formée des sections $F$ de $\mathcal{M}$ sur $J$ telles que l'on ait $F(j) \in (\mathcal{E}_j)_\pi \ \forall\ j \in \operatorname{Ob} J$,} \msp{page 264} \textit{$(\mathcal{E}_j)_\pi$ désignant la sous-catégorie pleine de $\mathcal{E}_j$ formée des objets $\pi$-accessibles de $\mathcal{E}_j$.}

En effet, l'hypothèse au début implique l'existence de $\pi_0$ et de $\pi$ satisfaisant les conditions dites. Les conditions a) et b) sur $\pi$, plus le fait que les foncteurs $\lambda_*$ soient $\pi_0$-accessibles, impliquent que la sous-catégorie strictement pleine $\mathcal{N}_\pi$ de $\mathcal{N}$ dont les fibres sont les $(\mathcal{N}_j)_\pi$, est une sous-catégorie cofibrée, et que de plus $\mathcal{N}$ est $J$-équivalent à $\Ind_\pi(\mathcal{N}_\pi/J)$. D'autre part, les fibres de $\mathcal{N}_\pi/J$ sont stables par $\varinjlim$ grandes devant $\pi_0$ et de cardinal $\leqslant \pi$, donc la condition “ordinale” de 4.10.1 est vérifiée (en y faisant $\mathcal{D} = \mathcal{N}_\pi$). Comme de plus $\operatorname{card} J \leqslant \pi$, on en conclut par 4.10.1 que $\mathcal{M} \simeq \Ind_\pi(\mathcal{M}_\pi)$.

\msp{page 265}

\dgrblock{dgrb1057}{The earlier language of cardinal filtrations}
\begin{dcomment}{context}{p264}{The earlier language of cardinal filtrations}
Grothendieck is returning to a problem already treated in the foundations of SGA~4: how a category assembled from small objects retains a small supply of generators after one imposes compatible structure. Expos\'e I first develops ind-objects and simultaneous presentations in Section~8, then cardinal filtrations and categories of sections in Section~9. The concrete antecedents are 8.8.2--8.8.5 and 9.22--9.24 \citep{SGA4I1972}. The references to SGA in this chapter therefore identify specific constructions, not merely a historical affinity.

The present argument makes a useful separation within that method. One cardinal bounds the arrows, equations, and chosen approximations. Another colimit shape supplies the completion of the successive corrections. Once those two roles have been isolated, a theorem about ordinary diagrams becomes a theorem about sections, because the correction is performed fibre by fibre and then checked against transport. The distinction between the category of all sections and its cartesian or cocartesian part must survive that passage: invertibility of a structure map is additional structure to enforce, not a consequence of ordinary naturality.

This is also a way to read the source's repeated requests to write out apparently formal steps. The same filtered system must witness an object, its maps, and their relations. A proof which merely chooses representatives of the separate objects has not yet constructed that system. The cofinality criterion and the iteration are precisely the two parts which supply it. They remain ordinary categorical arguments here; the homotopical relevance comes from applying them to categories of diagrams, sections, and specified classes of arrows.
\end{dcomment}


\textbf{Corollaire 4.10.9.} \textit{Soit $\mathcal{N} \longrightarrow J$ une “cofibration en paratopos” (\footnote{terme ambigu, cf. la double explication.}), i.e. un foncteur cofibrant tel que les fibres soient des catégories paratopos, et les foncteurs cochangement de base soient des foncteurs images inverses (resp. images directes) pour des morphismes de paratopos, i.e. qu'ils commutent aux (petites) $\varinjlim$ (resp. aux petites $\varprojlim$ ). Supposons de plus $J$ équivalente à une petite catégorie. Alors $\uGamma(\mathcal{N}/J)$ est une catégorie paratopos.}

En effet, c'est une catégorie accessible en vertu de 4.10.8, d'autre part l'hypothèse sur les $\lambda_*$ implique aussitôt que $\uGamma(\mathcal{N}/J)$ est stable par petites $\varprojlim$ (resp. par petites $\varprojlim$
\printedpage{100}
% head: A. Grothendieck : Dérivateurs, Ch. XVIII, Édition des Universités de Montpellier II et Paris VII
% foot: Version 17 décembre 2025
\texttt{[plutôt} $\varinjlim$ \texttt{?]}). Donc $\uGamma(\mathcal{N}/J)$ est une catégorie paratopos en vertu de 4.8.9 b $\Longrightarrow$ a (resp. c $\Longrightarrow$ a).

\textbf{Corollaire 4.10.10.} \textit{Soit $\mathcal{N} \longrightarrow J$ une catégorie} fibrée \textit{(non cofibrée!), telle que les fibres soient des catégories paratopos et les foncteurs} \msp{page 266} \textit{changement de base soient des foncteurs images directes pour des morphismes de paratopos (i.e. des foncteurs accessibles, commutant aux petites $\varprojlim$ ). Supposons de plus $J$ équivalente à une petite catégorie. Alors $\uGamma(\mathcal{N}/J)$ est une catégorie paratopos.}

\dgrblock{dgrb1058}{Paratoposes and the direction of pointwise limits}
\begin{dcomment}{exposition}{p265}{Paratoposes and the direction of pointwise limits}
The application to paratoposes uses the characterization established earlier in this chapter: an accessible category with all small limits, equivalently an accessible category with all small colimits, is locally presentable. ``Paratopos'' here should not be read as asserting the additional exactness properties of a Grothendieck topos. Module categories, including the preceding example, already lie in the intended range.

In the direct-image alternative we retain the words ``for a morphism of paratoposes'': the transports are actual right adjoints and are therefore accessible. Their bare preservation of limits, mentioned in the parenthesis, is not by itself the hypothesis needed to invoke 4.10.8. Indeed, to see accessibility of a right adjoint $R$, take a small strongly generating family of presentable objects $a$ in its target. Choose a regular cardinal bounding the presentability ranks of both the $a$ and their images $La$ under its left adjoint. Testing the filtered-colimit comparison for $R$ by $\operatorname{Hom}(a,R-)\simeq\operatorname{Hom}(La,-)$ now makes it a bijection for every $a$, hence an isomorphism. This gives one common rank. The next corollary explicitly includes accessibility.

For ordinary sections of an opfibration, limits are pointwise whenever the fibres have them, with no preservation requirement on transport. Given a diagram of sections $X^d$, put $L_j=\lim_d X^d_j$. For $a:i\to j$, the maps $a_*L_i\to a_*X^d_i\to X^d_j$ form a cone and therefore induce $a_*L_i\to L_j$. Their unit and composition equations follow after composing with each limit projection. A cone of sections factors uniquely through the $L_j$, and the same projections verify naturality of the factorization. This proves the limit universal property in the section category.

By contrast, constructing colimits pointwise uses preservation of those colimits by $a_*$, since the needed arrow is $a_*\operatorname{colim}X^d_i\to\operatorname{colim}X^d_j$. Thus the colimit-preserving case of 4.10.9 gives pointwise colimits as well as limits. In the limit-preserving case, pointwise limits suffice, together with accessibility, for local presentability. Both printed occurrences of $\varprojlim$ in the proof are defensible. The editors' tentative replacement of the second by $\varinjlim$ is not needed and does not follow from the stated preservation hypothesis. We retain the printed text and its question rather than turning a tentative editorial suggestion into a mathematical correction.

For a fibration the directions reverse. A section carries maps $X_i\to a^*X_j$. Colimits are pointwise without a preservation condition, while pointwise limits require that $a^*$ preserve them. Consequently 4.10.10 follows either from accessibility of sections and pointwise limits, or by replacing each right-adjoint inverse-image functor by its left adjoint, obtaining the equivalent description by an opfibration. The arrow labels in its printed proof are inconsistent with the direction $a:i\to j$; the adjunction is correctly $a_!:\mathcal N_i\rightleftarrows\mathcal N_j:a^*$.
\end{dcomment}


En effet, l'hypothèse implique que pour toute flèche $\lambda : j' \longrightarrow j$ dans $J$, le foncteur $\lambda_*$ : $\mathcal{N}_{j'} \longrightarrow \mathcal{N}_j$ admet un adjoint à gauche $\lambda^* : \mathcal{N}_j \longrightarrow \mathcal{N}_{j'}$, d'où s'ensuit que $\mathcal{N} \longrightarrow J$ est aussi cofibrant, les foncteurs de cochangement de base commutant aux petites $\varinjlim$ . Donc on peut appliquer 4.8.9, qui donne le résultat voulu.

\textbf{Question 4.10.11.} \textit{Soit $\mathcal{N} \longrightarrow J$ un foncteur} fibrant, \textit{tel que les $\mathcal{N}_j$ soient des catégories accessibles, et les foncteurs $\lambda^*$ de changement de base accessibles, $J$ étant équivalente à une petite catégorie. Je} \msp{page 267} \textit{conjecture que sous ces conditions, $\uGamma(\mathcal{N}/J) \defeq \Hom(J,\mathcal{N})$ est accessible, et qu'il en est de même de $\Hom\,\underline{\operatorname{cart}}_J(J,\mathcal{N})$.}

C'est ce qui est prouvé dans SGA 4 I 9.24 sous des hypothèses légèrement plus fortes (à l'aide de l'ancienne notion trop forte à mon gré, de filtrations cardinales). Je serais étonné que la démonstration (un peu longuette) de loc. cit. ne s'applique pas au cas général. Mais on aimerait avoir un énoncé qui coiffe celui-ci et 4.10.8, dans le cas d'un foncteur $\mathcal{N} \longrightarrow J$ qui ne serait supposé ni fibrant ni cofibrant, mais \ldots(???).

\section*{4.11\quad Application à l'existence des limites dans une catégorie de la forme $\Ind_{\pi_0}(C)$.}

\textbf{Proposition 4.11.1.} \textit{Soit} \msp{page 268} \textit{$C$ une catégorie équivalente à une petite catégorie, $\pi$ un cardinal infini, posons $\mathcal{M} = \Ind_\pi(C)$. Soit $\Phi$ un ensemble de catégories de cardinal réduit $\leqslant \pi$. On se propose d'examiner si $\mathcal{M}$ est “stable par $\varinjlim$ de type $\Phi$” (i.e. stable par $\varinjlim$ de type $I$, quel que soit $I \in \Phi$). Il suffit pour cela que $C$ est stable par $\varinjlim$ de type $\Phi$, et que de plus il existe un ordinal limite $\alpha$ avec $\operatorname{card}\alpha \leqslant \pi$ et tel que $C$ soit stable par $\varinjlim$ de type $I_\alpha$ (p.ex. tel que $I_\alpha \in \Phi$). Inversement, si $\mathcal{M}$ est stable par $\varinjlim$ de type $\Phi$, et si on suppose de plus $C$ karoubienne (i.e. stable par facteurs directs), alors $C$ est stable par $\varinjlim$ de type $\Phi$. Dans l'un et l'autre cas, le foncteur d'inclusion $C \into \mathcal{M}$ commute aux $\varinjlim$ envisagées.}

\dgrblock{dgrb1059}{The question about fibred sections is settled}
\begin{dcomment}{status}{p267}{The question about fibred sections is settled}
Grothendieck asks for accessibility of ordinary and cartesian sections under accessible inverse-image functors. Both assertions are now proved. The first construction is an oplax limit of the pseudofunctor $J^{\mathrm{op}}\to\mathrm{Acc}$ classifying the fibration; the second is its pseudolimit. After the ranks have been enlarged uniformly, \citet[Theorems~8.1, 8.3 and Corollary~9.2]{Positselski2024Limits} give these conclusions and pointwise descriptions of the presentable objects. The hypotheses on a common smaller chain-colimit cardinal and on preservation of presentables in the pseudolimit theorem are essential to that precise fixed-rank statement. This supplies a direct modern answer to 4.10.11, rather than merely an analogy with it.

Here is how the categorical construction encodes the question. Form the product $P=\prod_j\mathcal N_j$. For each $a:i\to j$, a potential section needs an arrow $X_i\to a^*X_j$. An inserter adds all these arrows to an object of $P$. Equifiers then impose the equations for identity arrows and composable pairs. The result is the category of ordinary sections. To obtain cartesian sections, add inverses to the structure maps and equify the two composites with the identity maps. The coherence isomorphisms of the pseudofunctor occur in the composition equations. This accounts for the actual section category, rather than a category in which compatibility is imposed only up to an unspecified isomorphism.

The required common cardinals can be chosen because there are only set-many fibres and transport functors. First choose a regular $\lambda$ for which all fibres are accessible and all transports preserve $\lambda$-filtered colimits. A $\lambda$-chain is itself $\lambda$-filtered, so these chain colimits exist and are preserved. Then choose a sufficiently large regular $\kappa>\lambda$, with the appropriate sharp cardinal inequality, so that all fibres are $\kappa$-accessible, $J$ is $\kappa$-small, and all transports preserve $\kappa$-presentables when needed. The uniformization mechanism developed in 4.12 explains the last step.

SGA~4, I, 9.22--9.24 already treats ordinary and cartesian sections, using cardinal filtrations and stronger closure assumptions \citep{SGA4I1972}. The new formulation separates the smaller colimits used to construct compatible structure from the larger filtered colimits used to present arbitrary objects. The wish for an encompassing theorem is realized by closure of accessible categories under the relevant products, inserters, equifiers and pseudolimits. It does not assert closure under arbitrary strict equalizers of accessible functors, nor does it give a hypothesis-free theorem for an arbitrary functor to $J$.
\end{dcomment}


(\footnote{Expliciter ici les corollaires des types 4.13.20, 4.13.21 (p. 341).}).

\textsc{Démonstration}. a) Condition suffisante. On applique 4.9.3 pour conclure que
\printedpage{101}
% head: A. Grothendieck : Dérivateurs, Ch. XVIII, Édition des Universités de Montpellier II et Paris VII
% foot: Version 17 décembre 2025
\msp{page 269}
pour toute catégorie $J$ de cardinal réduit $\leqslant \pi$ (en particulier pour $J \in \Phi$), tout foncteur $\varphi : J \longrightarrow \mathcal{M} = \Ind_\pi(C)$ admet une représentation indicielle $\Phi : J \times I \longrightarrow C$, avec $I$ grande devant $\pi$. Ceci, plus la stabilité de $C$ par $\varinjlim$ de type $J$, suffit à établir l'existence de la $\varinjlim$ de $\varphi$.

La réciproque provient du fait que si $C$ est karoubienne, alors son image essentielle dans $\mathcal{M}$ n'est autre que $\mathcal{M}_\pi$, laquelle est stable dans $\mathcal{M}$ par toute $\varinjlim$ de cardinal $\leqslant \pi$ représentable dans $\mathcal{M}$. Si $\mathcal{M}$ est stable par $\varinjlim$ de type $\Phi$, il en est de même de $\mathcal{M}_\pi$, donc de $C$, q.e.d.

\textbf{Corollaire 4.11.2.} \textit{Pour que $\mathcal{M}$ soit stable par $\varinjlim$ filtrantes (resp. par petites $\varinjlim$ ), il suffit que $C$ le soit par $\varinjlim$ filtrantes (resp. par $\varinjlim$ )} de cardinal $\leqslant \pi$, \textit{et c'est aussi nécessaire si $C$ est karoubienne.}

Supposons $C$ stable par $\varinjlim$ filtrantes \msp{page 270} (resp. par $\varinjlim$ de cardinal $\leqslant \pi$, alors par 4.11.1 $\mathcal{M}$ est stable par ce même type de limites inductives. D'autre part, il est stable par $\varinjlim$ grandes devant $\pi$. Il suffit donc de prouver le

\dgrblock{dgrb1060}{Transporting colimits through an ind-presentation}
\begin{dcomment}{exposition}{p269}{Transporting colimits through an ind-presentation}
The representation theorem pays for itself immediately. Once a diagram $X:J\to\Ind_\kappa C$ is written as $X_j=\operatorname{colim}_{i\in I}D_{i,j}$ with $I$ $\kappa$-filtered and every $D_i$ a genuine $C$-valued diagram, the candidate colimit is forced:
\[
 Q=\operatorname{colim}_{i\in I}\left(\operatorname{colim}_{j\in J}^{C}D_{i,j}\right).
\]
The inner colimits are preserved by $C\to\Ind_\kappa C$ because $|\operatorname{Mor}J|<\kappa$; the outer one exists by definition of the completion. For any $Y$ in the completion,
\[
 \operatorname{Hom}(Q,Y)
 \simeq\lim_i\lim_j\operatorname{Hom}(D_{i,j},Y)
 \simeq\lim_j\operatorname{Hom}(X_j,Y).
\]
The middle interchange is an interchange of two limits, so it needs no extra exactness hypothesis. This proves 4.11.1's sufficiency, including the prescribed cocone.

The preservation assertion of Lemma~4.10.4 is equally specific. If $c=\operatorname{colim}_{j\in J}^{C}c_j$ and $Y=\operatorname{colim}_{a\in A}y_a$, then
\[
 \operatorname{Hom}(c,Y)
 \simeq\operatorname{colim}_a\lim_j\operatorname{Hom}(c_j,y_a)
 \simeq\lim_j\operatorname{colim}_a\operatorname{Hom}(c_j,y_a).
\]
Here the second interchange does use $\kappa$-filteredness of $A$ and $\kappa$-smallness of $J$. This distinction explains why no general colimit-preservation property of the ordinary Yoneda embedding is being claimed.

For necessity, an existing $\kappa$-small colimit of $\kappa$-presentables is again $\kappa$-presentable: its hom functor is a $\kappa$-small limit of hom functors, and the same interchange applies. The $\kappa$-presentable objects of $\Ind_\kappa C$ are retracts of objects of $C$. To see the retract statement, express a presentable $X$ as a filtered colimit of representables and factor $\mathrm{id}_X$ through one stage. If $C$ is idempotent complete, that retract lies in $C$, proving the converse. Idempotent completeness is used at precisely this point, not in the construction of $Q$.
\end{dcomment}


\textbf{Lemme 4.11.3.} \textit{Soit $\mathcal{M}$ une catégorie, $\pi$ un cardinal. Supposons que $\mathcal{M}$ soit stable par $\varinjlim$ filtrantes de cardinal $\leqslant \pi$ (resp. par $\varinjlim$ de cardinal $\leqslant \pi$, resp. par sommes des cardinal $\leqslant \pi$),} et \textit{par $\varinjlim$ grandes devant $\pi$. Alors $\mathcal{M}$ est stable par petites $\varinjlim$ filtrantes (resp. par petites $\varinjlim$ , resp. par petites sommes).} % [sic] des cardinal -> de cardinal, agreement (cf. "de cardinal" elsewhere in the statement)

Ce lemme se prouve comme 4.4.14, page 97.

Reste la réciproque dans 4.11.2, qui résulte de l'assertion similaire dans 4.11.1.

\textbf{Corollaire 4.11.4.} \textit{Quand $\mathcal{M} = \Ind_\pi C$ est-il stable par sommes ? S'il existe un ordinal $\alpha$ tel que $\operatorname{card}\alpha \leqslant \pi$, et que $C$ soit stable par $\varinjlim$ de type $I_\alpha$, alors il suffit que $C$ soit stable par sommes de cardinal $\leqslant \pi$. Cette condition est aussi nécessaire si $C$ est karoubienne.}

(\footnote{Pour autres variantes de 4.11.1, cf. 4.13.19 (p. 340) et corollaires.}).

Se prouve comme 4.11.2, en utilisant 4.11.3.

\msp{page 271}

\dgrblock{dgrb1061}{Building an arbitrary colimit from bounded pieces}
\begin{dcomment}{exposition}{p270}{Building an arbitrary colimit from bounded pieces}
The cardinal enlargement in Lemma~4.11.3 can be made quite explicit. Assume $\pi$ is infinite, as it is in the application, and put $\kappa=\pi^+$. Every small category $J$ is the union of its subcategories with at most $\pi$ arrows. Their poset $P$ is $\kappa$-filtered: the subcategory generated by the union of at most $\pi$ such subcategories still has at most $\pi$ arrows, since only finite composites are added. For a diagram $X:J\to\mathcal M$, suppose all the colimits $Q_K=\operatorname{colim}_{K}X$ exist for $K\in P$. A cocone from the diagram $K\mapsto Q_K$ to $Y$ is exactly a cocone from $X$ to $Y$. Every object and arrow of $J$ occurs in some $K$, and two occurrences can be compared in a larger one. Thus $\operatorname{colim}_{K\in P}Q_K$ is the desired colimit over $J$.

When $J$ is filtered and only filtered colimits are assumed for the bounded pieces, take instead its filtered subcategories of size at most $\pi$. Every set of at most $\pi$ arrows is contained in one: repeatedly add cocones for pairs of objects and equalizers for parallel pairs, then close under composition. Countably many repetitions suffice; their union is filtered and still has at most $\pi$ arrows. The same construction bounds a family of at most $\pi$ such subcategories. Their indexing poset is therefore $\kappa$-filtered, and the preceding cocone argument applies unchanged. For coproducts use the poset of subsets of the indexing set of size at most $\pi$.

This also makes 4.11.4 stronger than its stated form. If $C$ has coproducts indexed by sets of size at most $\pi$, then $\Ind_\kappa C$ has all small coproducts without an additional ordinal-colimit assumption. For a bounded family of ind-objects, take the product of their $\kappa$-filtered presentation categories, form the coproduct of the resulting objects of $C$ at each index, and then take its filtered colimit. The hom universal property proves this is the coproduct. The preceding bounded-subset construction gives arbitrary coproducts. The converse under idempotent completeness follows from the presentability argument on p.~269. Thus the omission of the word ``limit'' before ``ordinal'' in 4.11.4 creates no counterexample to its conclusion.
\end{dcomment}


\textbf{Proposition 4.11.5} (\footnote{?, démonstration pas achevée.}). \textit{Soit $\Phi$ un ensemble de catégories $I$ équivalentes à des catégories finies, et “rigides”, i.e. pour tout $i \in \operatorname{Ob} I$, $\Hom_I(i,i) = \{\mathrm{id}_i\}$. Soient $C$ une catégorie équivalente à une petite catégorie, $\pi$ un cardinal infini, $\mathcal{M} = \Ind_\pi C$. Pour que $\mathcal{M}$ soit stable par $\varinjlim$ (resp. par $\varprojlim$ ) de type $\Phi$, il} suffit \textit{que $C$ le soit. Dans le cas non respé, si $C$ est karoubienne, cette condition est aussi nécessaire.}

Cette dernière assertion est contenue dans l'assertion similaire dans 4.11.1. Pour la première assertion, cela résulte de la façon habituelle du

\textbf{Lemme 4.11.6} (\footnote{?, démonstration non explicitée est peut-être faux!}). \textit{Soit $J$ une catégorie équivalente à une catégorie finie, $J$ rigide, $C$ une catégorie équivalente à une petite catégorie, $\pi$ un cardinal infini. Alors tout foncteur}
\printedpage{102}
% head: A. Grothendieck : Dérivateurs, Ch. XVIII, Édition des Universités de Montpellier II et Paris VII
% foot: Version 17 décembre 2025
\textit{$J \longrightarrow \mathcal{M} \defeq \Ind_\pi C$ admet une représentation indicielle $J \times I \longrightarrow C$, avec $I$ grande devant $\pi$, en d'autres termes, le} \msp{page 272} \textit{foncteur canonique pleinement fidèle}
\[
\Ind_\pi(\Hom(J,C)) \ \longrightarrow \ \Hom(J,\mathcal{M})
\]
\textit{déduit de l'inclusion $C \into \mathcal{M}$, d'où $\Hom(J,C) \into \Hom(J,\mathcal{M})$ (cf. 4.9), est une équivalence de catégories.} (\textbf{N.B.} J'ignore s'il en est ainsi sans hypothèse de rigidité sur $J$.)

\textsc{Démonstration}. Cela se prouve comme le résultat similaire pour $\Ind(C)$ au lieu de $\Ind_\pi C$, cf. SGA 4 I 8.8.5.

\msp{page 273}

\dgrblock{dgrb1062}{Finite rigid diagrams really do admit simultaneous presentations}
\begin{dcomment}{exposition}{p272}{Finite rigid diagrams really do admit simultaneous presentations}
The hesitation attached to 4.11.6 concerns a true statement. After taking a skeleton, a finite category with only identity endomorphisms admits an ordering of its objects in which every nonidentity arrow increases the order. Indeed a cycle would produce an endomorphism, and the identity-only condition would make every arrow in the cycle invertible, contradicting skeletal distinctness. The finite induction in SGA~4, I, 8.8.5 therefore extends to the present cardinal bound \citep{SGA4I1972}.

We give the induction with the required simultaneous constraints. Fix $X:J\to\Ind_\kappa C$ and let $S=C^J/X$. To show that $S$ is $\kappa$-filtered, take a diagram $H:K\to S$ with $|\operatorname{Mor}K|<\kappa$. Construct a cocone $H\to D$ one vertex of $J$ at a time. At a vertex $j$, include in $C/X_j$ every $H_k(j)$, with the arrows prescribed by $K$, and every previously chosen $D(i)$ indexed by an incoming nonidentity arrow $a:i\to j$, mapped to $X_j$ through $X(a)$. Include also the maps $H_k(i)\to D(i)$ and $H_k(i)\to H_k(j)$, and, for composable incoming arrows, the maps between previously chosen $D(i)$'s. This is a diagram with fewer than $\kappa$ arrows. A cocone in the $\kappa$-filtered category $C/X_j$ gives $D(j)$, the arrows $D(a)$, and the components $H_k(j)\to D(j)$. The imposed relations say exactly that composition is respected and that every $H_k\to D$ is natural. Since no nonidentity arrow leaves $j$ for an earlier vertex, subsequent choices cannot alter this work. The induction supplies the desired cocone, including nonemptiness when $K$ is empty.

The same induction can begin with arbitrary objects $c_j\to X_j$, adding them as isolated objects to the diagrams at the corresponding vertices. It proves the object-lifting condition for evaluation $S\to\prod_j C/X_j$. To see the required connectedness, take two lifts of a fixed family. Repeat the induction with the two section diagrams and with the additional pairs of arrows from each $c_j$; equalizing these pairs in $C/X_j$ gives a common successor in the relevant comma category. Evaluation is therefore final. The density argument already given on p.~246 proves the equivalence $\Ind_\kappa(C^J)\simeq(\Ind_\kappa C)^J$.

The modern sharp result is \citet[Theorem~1.3]{Henry2025Ind}: for all categories $C$, this equivalence holds exactly when $J$ is essentially $\kappa$-small and well-founded. Finite rigid categories are such shapes. Thus the source's further question about dropping rigidity has a negative answer in general; a finite category with a nonidentity endomorphism is already outside the universal statement. The condition concerns a theorem for every $C$, not the impossibility of successful representation for one particular $C$.
\end{dcomment}


% a vertical rule is printed in the left margin beside Corollaire 4.11.7
\textbf{Corollaire 4.11.7} (\footnote{sous réserve de validité de 4.11.5.}). \textit{Pour que $\mathcal{M} = \Ind_\pi C$ soit stable par $\varinjlim$ (resp. par $\varprojlim$ )} finies, \textit{il} suffit \textit{que $C$ le soit. Et c'est aussi nécessaire dans le cas non respé, lorsqu'on suppose $C$ karoubienne.}

\section*{4.12\quad Foncteurs accessibles entre catégories accessibles}

\textbf{Proposition 4.12.1} (\footnote{Remonterait à 4.8.}). \textit{Soient $\mathcal{M}$, $\mathcal{N}$ deux catégories accessibles (4.8), et $f : \mathcal{M} \longrightarrow \mathcal{N}$ un foncteur accessible. Alors il existe des cardinaux infinis $\pi_0$, $\pi_1$ tels que $\pi_1 \geqslant \pi_0$ et que pour tout cardinal $\pi \geqslant \pi_1$, on ait}
\[
f(\mathcal{M}_\pi) \ \subset \ \mathcal{N}_{\pi^{\pi_0}} \ , \tag*{(4.12.1)} % [sic] (4.12.1) -> (4.12.1.1), the formula is cited below as (4.12.1.1) and the next one is numbered (4.12.1.2)
\]
\textit{et en particulier}
\[
f(\mathcal{M}_\pi) \ \subset \ \mathcal{N}_\pi \qquad \textit{si } \pi = \pi^{\pi_0}. \tag*{(4.12.1.2)}
\]
\textit{De façon plus précise, on peut prendre pour $\pi_0$ un cardinal utile pour $\mathcal{M}$} \msp{page 274} \textit{et tel que $f$ soit $\pi_0$-accessible (p.ex. le plus petit cardinal de ce type). Soit alors $\pi_1$ un cardinal tel qu'on ait $f(\mathcal{M}_{\pi_0}) \subset \mathcal{N}_{\pi_1}$, $\pi_1 \geqslant \cardred \mathcal{M}_{\pi_0}$, et enfin tel que $\mathcal{N}_{\pi_1}$ soit} \texttt{[une]} \textit{sous-catégorie strictement génératrice de $\mathcal{N}$. Le couple $(\pi_0,\pi_1)$ satisfait aux conditions énoncées.}

\dgrblock{dgrb1063}{Why the finite limit and colimit corollaries follow}
\begin{dcomment}{exposition}{p273}{Why the finite limit and colimit corollaries follow}
The preceding representation establishes both clauses of 4.11.5, but they use different universal-property arguments. For colimits, the iterated-colimit calculation on p.~269 applies directly. For limits, the inclusion $C\to\Ind_\kappa C$ preserves every existing limit in $C$: the ordinary Yoneda embedding preserves limits, and the completion is a full subcategory of the presheaf category.

Let $D:J\to\Ind_\kappa C$ have a simultaneous presentation $D_j=\operatorname{colim}_{i\in I}D_{i,j}$, with $J$ finite rigid. If $C$ has $J$-limits, put $L_i=\lim_j^C D_{i,j}$ and $L=\operatorname{colim}_i L_i$. For a representable $c\in C$, compactness and filtered-limit commutation in sets give
\[
 \operatorname{Hom}(c,L)
 \simeq\operatorname{colim}_i\lim_j\operatorname{Hom}(c,D_{i,j})
 \simeq\lim_j\operatorname{Hom}(c,D_j).
\]
For an arbitrary ind-object $Z=\operatorname{colim}_a c_a$, apply $\lim_a$ to this identity and commute the two limits. This proves $\operatorname{Hom}(Z,L)\simeq\lim_j\operatorname{Hom}(Z,D_j)$, hence $L=\lim_j D_j$.

Although an arbitrary finite category need not be rigid, finite limits are generated by finite products and equalizers, and finite colimits by finite coproducts and coequalizers. The corresponding shapes are finite rigid categories. The empty shape provides the terminal or initial object, and the parallel-pair shape provides equalizers or coequalizers. This proves Corollary~4.11.7 without silently applying 4.11.6 to a finite category with endomorphisms. The converse for finite colimits uses idempotent completeness as before; no analogous converse for finite limits is asserted.
\end{dcomment}


\textsc{Démonstration}. On aura en effet
\[
\mathcal{N}_1 \ \xleftarrow{\ \sim\ } \Ind_{\pi_0}(\mathcal{M}_{\pi_0})
\]
et $f$ commutant aux $\varinjlim$ grandes devant $\pi_0$, sera connu (à isomorphisme unique près) par sa restriction à $\mathcal{M}_{\pi_0}$, laquelle se factorise par $\mathcal{N}_{\pi_1}$,
\[
f|\mathcal{M}_{\pi_0} : \mathcal{M}_{\pi_0} \ \longrightarrow \ \mathcal{N}_{\pi_1} \ .
\]
Pour prouver (4.12.1.1), posant $\pi' = \pi^{\pi_0} \geqslant \pi$, et notant que $\pi'^{\pi_0} = \pi_0$ \texttt{[plutôt} $\pi'^{\pi_0} = \pi'$\texttt{]},
\[
f(\mathcal{M}_\pi) \ \subset \ f(\mathcal{M}_{\pi'}) \ \overset{?}{\subset} \ \mathcal{N}_{\pi'} \ ,
\]
\printedpage{103}
% head: A. Grothendieck : Dérivateurs, Ch. XVIII, Édition des Universités de Montpellier II et Paris VII
% foot: Version 17 décembre 2025
on voit qu'il suffit de prouver (4.12.1.2) (pour $\pi'$). Donc on peut supposer $\pi = \pi^{\pi_0}$ et prouver (4.12.1.2). Mais comme $\pi_1 \geqslant \cardred \mathcal{M}_{\pi_0}$, on sait que si $\pi$ est tel que $\pi \geqslant \pi_1$, $\pi^{\pi_0} = \pi$, alors $\mathcal{M}_\pi$ est formé des $\varinjlim\limits_{I} x_i$, les $x_i \in \mathcal{N}_{\pi_0}$, et $I$ grande devant $\pi_0$, et $\operatorname{card} I \leqslant \pi$.

\msp{page 275}

On a alors ($f$ étant $\pi_0$-accessible)
\[
f(x) \ = \ \varinjlim_{I} f(x_i) \ ,
\]
les $f(x_i)$ sont dans $\mathcal{N}_{\pi_1}$, donc dans $\mathcal{N}_\pi$, or $\mathcal{N}_\pi$ stable par $\varinjlim$ de cardinal $\leqslant \pi$ (puisque $\mathcal{N}$ est engendrée par $\mathcal{N}_{\pi_1}$, formée d'objets $\pi_1$-accessibles et a fortiori $\pi$-accessibles, cf. SGA 4 I 9.9). D'où $f(x) \in \mathcal{N}_\pi$, q.e.d.

\textbf{Corollaire 4.12.2.} \textit{Supposons que $\pi_0$ soit aussi utile pour $\mathcal{N}$, et que $\pi_1 \geqslant \cardred \mathcal{N}_{\pi_0}$. Soit $\pi$ un cardinal $\geqslant \pi_1$, et tel que $\pi^{\pi_0} = \pi$. Alors $f$ s'identifie, à équivalence de catégories et à isomorphisme de foncteurs près, au foncteur}
\[
\Ind(f_\pi) \ : \ \Ind_\pi(\mathcal{M}_\pi) \ \longrightarrow \ \Ind_\pi(\mathcal{N}_\pi) \ ,
\]
\textit{où $f_\pi : \mathcal{M}_\pi \longrightarrow \mathcal{N}_\pi$ est le foncteur induit par $f$.}

\textbf{Proposition 4.12.3} (\footnote{remonte à 4.7.}). \textit{Soient $C$ une catégorie équivalente à une petite catégorie, $\mathcal{N}$ une $\mathfrak{U}$-catégorie, $\pi_0$ un cardinal infini.} \msp{page 276} \textit{On suppose que $\mathcal{N}$ satisfait $L_{\pi_0}$ (stabilité par $\varinjlim$ grandes devant $\pi_0$). Soit $\mathcal{M} = \Ind_{\pi_0} C$, $\Hom_{!\,\pi_0}(\mathcal{M},\mathcal{N})$ la sous-catégorie strictement pleine de $\Hom(\mathcal{M},\mathcal{N})$ formée des $f$ : $\mathcal{M} \longrightarrow \mathcal{N}$ $\pi_0$-accessibles. Alors le foncteur restriction}
\[
\Hom_{!\,\pi_0}(\mathcal{M},\mathcal{N}) \ \longrightarrow \ \Hom(C,\mathcal{N}) \tag*{(4.12.3.1)}
\]
\textit{est une équivalence de catégories. On trouve un foncteur quasi-inverse comme composé}
\[
\Hom(C,\mathcal{N}) \ \longrightarrow \ \Hom(\underbrace{\Ind_{\pi_0}(C)}_{\mathcal{M}},\Ind_{\pi_0}(\mathcal{N})) \ \longrightarrow \ \Hom(\mathcal{M},\mathcal{N}) \ , \tag*{(4.12.3.2)}
\]
\textit{où la deuxième flèche est le foncteur composition $F \longmapsto L \circ F$ avec le foncteur limite inductive $L : \Ind_{\pi_0}\mathcal{N} \longrightarrow \mathcal{N}$. Le foncteur composé (4.11.3.2) a une image contenue dans $\Hom_{!\,\pi_0}(\mathcal{M},\mathcal{N})$ et induit une équivalence entre cette dernière catégorie et $\Hom(C,\mathcal{N})$.} % [sic] (4.11.3.2) -> (4.12.3.2), the composed functor is the display just labelled (4.12.3.2)

\dgrblock{dgrb1064}{An accessible functor eventually respects a common size bound}
\begin{dcomment}{exposition}{p275}{An accessible functor eventually respects a common size bound}
The point of 4.12.1 is to replace the qualitative word ``accessible'' by a controlled estimate on the images of small objects. Choose a useful $\pi_0$ for $\mathcal M$ such that $f$ preserves colimits indexed by categories large before $\pi_0$. The essentially small category $\mathcal M_{\pi_0}$ has only set-many images under $f$. Their accessibility ranks can therefore be bounded by one $\pi_1$, enlarged also to dominate the size of a skeleton and a generating rank for $\mathcal N$.

For $\pi\geq\pi_1$ satisfying $\pi^{\pi_0}=\pi$, the earlier presentation theorem expresses every $x\in\mathcal M_\pi$ as a colimit of objects $x_i\in\mathcal M_{\pi_0}$ over a category $I$ large before $\pi_0$ with $|\operatorname{Mor}I|\leq\pi$. Since $f$ preserves this colimit,
\[
 f(x)\simeq\operatorname{colim}_{i\in I}f(x_i).
\]
Every $f(x_i)$ is $\pi$-accessible, and a colimit of at most $\pi$ such objects is $\pi$-accessible. The latter assertion follows by writing its hom functor as a $\pi$-bounded limit and commuting it with the filtered colimits used to define $\pi$-accessibility. Thus $f(\mathcal M_\pi)\subset\mathcal N_\pi$. For general $\pi$, put $\rho=\pi^{\pi_0}$. Cardinal arithmetic gives $\rho^{\pi_0}=\rho$, and monotonicity gives $\mathcal M_\pi\subset\mathcal M_\rho$, proving the advertised bound $f(\mathcal M_\pi)\subset\mathcal N_\rho$.

The printed $\mathcal N_1\simeq\Ind_{\pi_0}(\mathcal M_{\pi_0})$ on p.~274 must mean $\mathcal M\simeq\Ind_{\pi_0}(\mathcal M_{\pi_0})$, and the $x_i$ later belong to $\mathcal M_{\pi_0}$, not $\mathcal N_{\pi_0}$. These are present in the edition and are retained as source slips. The proof itself is the familiar uniformization argument. Applied to a small family of functors, one first bounds all their required ranks, then chooses a common fixed point; this is the step used for sections and cartesian sections.

Once the same $\pi$ is useful for both categories, preservation of the appropriate filtered colimits says that $f$ is determined by $f_\pi:\mathcal M_\pi\to\mathcal N_\pi$. Corollary~4.12.2 identifies it with the ind-extension of that small functor. The identification includes natural transformations, as the following universal property makes precise.
\end{dcomment}


\msp{page 277}

\textit{En d'autres termes, on peut dire que le foncteur d'inclusion canonique}
\[
C \ \longrightarrow \ \Ind_{\pi_0} C \ = \ \mathcal{M}
\]
\textit{est} 2-universel \textit{pour les foncteurs de $C$ dans des catégories $\mathcal{N}$ satisfaisant $L_{\pi_0}$, celles-ci étant considérées comme $0$-objets d'une $2$-catégorie où les $1$-objets sont les foncteurs $\mathcal{N} \longrightarrow \mathcal{N}'$ $\pi_0$-accessibles, i.e. commutent aux $\varinjlim$ grandes devant $\pi_0$.}

Cet énoncé est l'analogue, pour $C$ petite quelconque, de 4.4.16 (page 101), valable pour $C$ stable par $\varinjlim$ de cardinal $\leqslant \pi_0$. Il y a bien sûr d'autres variantes, en relation avec la commutation à d'autres types de $\varinjlim$ ( $\varinjlim$ filtrantes, sommes, limites inductives finies etc.) envisagées dans la section 4.11 précédente.

La démonstration de 4.12.3 est soritale.

% ---- from der-ch18_p106-110.tex
\printedpage{104}
% head: A. Grothendieck : Dérivateurs, Ch. XVIII, Édition des Universités de Montpellier II et Paris VII
% foot: Version 17 décembre 2025
\msp{page 278}
\textbf{Corollaire 4.12.4.} \textit{Soient $\M$ et $\N$ deux catégories accessibles, et soit $\pi_0$ un cardinal infini tel que $\M$ soit $\pi_0$-accessible, i.e. que $\pi_0$ soit un cardinal utile pour $\M$. Soit $\pi_1 = \sup(\pi_0,\allowbreak \card\red\M_{\pi_0})$. Alors pour tout cardinal $\pi \geqslant \pi_1$ utile pour $\M$ (et en particulier, pour tout $\pi \geqslant \pi_1$ tel que $\pi^{\pi_0} = \pi$), la catégorie des foncteurs $\M \longrightarrow \N$ qui sont $\pi$-accessibles est elle-même accessible.}

\dgrblock{dgrb1065}{The free filtered-colimit completion includes natural transformations}
\begin{dcomment}{exposition}{p277}{The free filtered-colimit completion includes natural transformations}
Grothendieck's ``2-universal'' formulation records more than existence of an extension. Given $g:C\to\mathcal N$, where $\mathcal N$ admits $\kappa$-filtered colimits, define
\[
 \overline g(X)=\operatorname{colim}_{(c\to X)\in C/X}g(c),
 \qquad X\in\Ind_\kappa C.
\]
The category $C/X$ is small and $\kappa$-filtered. A map $X\to Y$ induces $C/X\to C/Y$, and hence a map between the displayed colimits. Identity maps and composition are respected by uniqueness in the colimit universal property. For $X=c$ representable, $(c,\mathrm{id})$ is terminal in $C/c$, so $\overline g(c)\simeq g(c)$.

This extension preserves $\kappa$-filtered colimits. For a $\kappa$-filtered diagram $X_i$ with colimit $X$, consider pairs $(i,c\to X_i)$. Their category is $\kappa$-filtered: fewer than $\kappa$ objects and equations can first be moved to a common index $i$, using presentability of the representables, and then bounded inside $C/X_i$. The functor sending a pair to $c\to X$ is final. Every such map factors through some $X_i$, and two factorizations become equal at a later index; the same argument with fewer than $\kappa$ constraints verifies the required compatibility. Therefore the displayed formula for $\overline g(X)$ is the iterated colimit of the formulas for $\overline g(X_i)$, proving preservation.

If $F$ already preserves these colimits, the canonical expression $X=\operatorname{colim}_{C/X}c$ forces $F(X)\simeq\overline{F|_C}(X)$. Finally a transformation $g\to h$ induces $\overline g\to\overline h$ by applying its components to each presentation. Conversely a transformation between colimit-preserving extensions is forced by its restriction to $C$, since the maps from the presentation objects are a colimit cocone. Restriction is consequently fully faithful and essentially surjective, exactly the equivalence of categories stated in 4.12.3. This is the cardinal-bounded counterpart of the ordinary ind universal property in SGA~4, I, 8.7 \citep{SGA4I1972}.
\end{dcomment}


En effet, on aura $\M \simeq \Ind_\pi\M_\pi$, d'où par 4.12.3
\[
\Hom_{!\,\pi}(\M,\N) \xrightarrow{\sim} \Hom(\M_\pi,\N) ,
\]
or $\Hom(\M_\pi,\N)$ est accessible par 4.9.5 (p. 220). On peut préciser d'ailleurs ce point. Supposons $\pi_0$ choisi tel que $\pi_0$ soit utile non seulement pour $\M$, mais pour $\N$ également. Rappelons que l'on a $\card\Fl\M_\pi \leqslant 2^\pi$ (cf. 4.7.7, page 173).

\msp{page 279}
Posons donc
\[
\pi' = 2^\pi ,
\]
on a donc $\pi'^{\pi_0} = 2^{\pi\pi_0} = 2^\pi = \pi'$, et $\card\red\M_\pi \leqslant \pi'$. Donc 4.9.5 nous donne que
\[
\Hom(\M_\pi,\N) \simeq \Ind_{\pi'}(\Hom(\M_\pi,\M_{\pi'})) , % [sic] \M_{\pi'} -> \N_{\pi'}, the second argument should lie in N, as in Corollaire 4.12.5
\]
d'où :

\dgrblock{dgrb1066}{Fixing an accessibility rank before forming a functor category}
\begin{dcomment}{exposition}{p278}{Fixing an accessibility rank before forming a functor category}
Grothendieck first fixes the amount of colimit preservation demanded of a functor. This makes its domain effectively small: a functor preserving the relevant filtered colimits is determined, including its natural transformations, by its values on a small category of presentable objects. Write $\lambda=\pi^+$ and $A=\mathcal M_\pi$. Whenever $\mathcal N$ admits $\lambda$-filtered colimits, restriction gives
\[
 \operatorname{Fun}_{\lambda\text{-filt}}(\mathcal M,\mathcal N)
 \simeq \operatorname{Fun}(A,\mathcal N).
\]
The inverse is the extension $F(x)=\varinjlim_{(a\to x)\in A/x}F(a)$. The same formula extends a natural transformation, and its uniqueness follows because the cocone $a\to x$ presents $x$. Thus the equivalence concerns categories of functors, rather than merely their isomorphism classes. This is the ordinary categorical form of the universal property of $\Ind$; compare \citet[Remark 2.3(1)]{ChornyRosicky2012}.

There is one hypothesis to restore in the proof as printed. Proposition 4.12.3 assumes $\mathcal N$ satisfies $L_\pi$, whereas the first statement of 4.12.4 chooses $\pi_0$ using $\mathcal M$ alone. The displayed restriction equivalence does not follow without the required colimits in $\mathcal N$. Choosing $\pi_0$ useful for both categories supplies this hypothesis, exactly as Grothendieck does in the paragraph leading to 4.12.5. This repairs the proof and gives the form used subsequently; we do not infer a counterexample to the broader assertion merely from this omitted hypothesis.

The cardinal $2^\pi$ in 4.12.5 has a different job from $\pi$. The latter controls how a functor is extended; the former is large enough to control all the objects and arrows in its small domain simultaneously. If a skeleton of $A$ has at most $2^\pi$ arrows and $\pi\geq\pi_0$, then
\[
 (2^\pi)^{\pi_0}=2^{\pi\pi_0}=2^\pi.
\]
Corollary 4.9.5 therefore applies to $\operatorname{Fun}(A,\mathcal N)$. The $\mathcal M_{\pi'}$ printed in the preliminary display on p.~279 must be read mathematically as $\mathcal N_{\pi'}$, as the following corollary confirms. Also, ``under the conditions of 4.12.4'' retains the requirement that $\pi$ be useful for $\mathcal M$; the shorter inequality $\pi\geq\pi_1$ in 4.12.5 does not by itself replace that requirement.
\end{dcomment}


\textbf{Corollaire 4.12.5.} \textit{Sous les conditions de 4.12.4, si $\pi_0$ est utile pour $\M$} et \textit{pour $\N$, alors pour $\pi \geqslant \pi_1$, $\Hom_{!\,\pi}(\M,\N)$ est $\pi'$-accessible, où $\pi' = 2^\pi$, de façon plus précise, on a une équivalence de catégories}
\[
\Hom_{!\,\pi}(\M,\N) \simeq \Ind_{\pi'}(\Hom(\M_\pi,\N_{\pi'})) .
\]

\textbf{Remarque 4.12.6.} Si $\M$, $\N$ sont accessibles, la catégorie $\Hom_{\mathrm{acc}}(\M,\N)$ des foncteurs accessibles n'est par toujours accessible. Si $\N$ satisfait $\mathrm{L}_{\pi_0}$, cette catégorie satisfait $\mathrm{L}_{\pi_0}$, % [sic] par toujours -> pas toujours, negation "n'est pas" intended
et les $\varinjlim$ grandes devant $\pi_0$ se calculent argument par argument.

\msp{page 280}
Si on pouvait trouver une sous-catégorie (équivalente à une petite catégorie) $\cH_\pi$ de $\cH = \Hom_{\mathrm{acc}}(\M,\N)$, pour un cardinal $\pi \geqslant \pi_0$ telle que tout objet $F$ de $\cH$ s'écrive $\varinjlim\limits_{I} F_i$, les $F_i$ dans $\cH_\pi$, $I$ grande devant $\pi$, alors si les objets de $\cH_\pi$ sont des foncteurs $\pi'$-accessibles de $\M$ dans $\N$, il en serait de même de tout $F$ dans $\cH$. Donc il existerait un cardinal infini $\pi'$ tel que tout foncteur accessible de $\M$ dans $\N$ soit $\pi'$-accessible. Mais sauf erreur, si p.ex. $\M = \N = \mathrm{Ens}$, il n'existe pas de tel cardinal. Je ne veux pas m'arrêter ici sur cette question, et faire un contre-exemple.

\msp{page 281}

\dgrblock{dgrb1067}{Accessible endofunctors of sets have no common bound}
\begin{dcomment}{examples}{p280}{Accessible endofunctors of sets have no common bound}
Grothendieck's hesitation on p.~280 can be settled by the most elementary representable functors. For each set $B$, consider
\[
 H_B(X)=\operatorname{Hom}(B,X)=X^B.
\]
If $\lambda$ is regular and $|B|<\lambda$, this functor preserves $\lambda$-filtered colimits. A map $B\to\varinjlim_iX_i$ is represented by choosing a representative for each of its fewer than $\lambda$ values and then a common stage. If two representatives define the same map, fewer than $\lambda$ equalities have to be witnessed, and filteredness places all their witnesses at a further common stage. These are respectively surjectivity and injectivity of the comparison map.

The bound cannot be made independent of $B$. Fix an infinite cardinal $\pi$, let $|B|>\pi$, and order the subsets $T\subseteq B$ of cardinal at most $\pi$ by inclusion. Their poset is $\pi^+$-directed: the union of at most $\pi$ such subsets still has cardinal at most $\pi$. Their colimit is $B$. But the image of
\[
 \varinjlim_{T\in[B]^{\leq\pi}} T^B\longrightarrow B^B
\]
consists precisely of the maps with image of cardinal at most $\pi$. It omits $\mathrm{id}_B$. Thus $H_B$ is accessible, for example at the regular cardinal $|B|^+$, but is not $\pi^+$-accessible.

This proves the stronger assertion about the category of all accessible endofunctors. Suppose that category were $\mu$-accessible and choose a small set $\mathcal G$ of $\mu$-presentable generators. There is a regular cardinal $\lambda$ such that every $G\in\mathcal G$ preserves $\lambda$-filtered colimits. Pointwise small colimits of accessible set-valued functors remain accessible: choose one common preservation rank for the small diagram, and interchange its colimit with a filtered colimit in $\mathrm{Set}$. Consequently these pointwise colimits are also colimits in the full category of accessible functors. Every object, being a $\mu$-filtered colimit of objects of $\mathcal G$, would therefore preserve $\lambda$-filtered colimits. The preceding example contradicts this conclusion. Natural transformations cause no size difficulty: between two accessible functors they are determined on a common small subcategory of sufficiently presentable sets.

This completes the counterexample requested by the source. The broader replacement of small generating sets by generating classes is the subject of \citet[Definition 2.1 and Examples 2.2]{ChornyRosicky2012}; that theory gives a context for the size obstruction, rather than making the ordinary accessibility assertion true.
\end{dcomment}


\section*{4.13\quad Ensembles accessibles d'objets d'une catégorie accessible}

\textbf{Définition 4.13.1.} Soient $\M$ une catégorie accessible, $S$ une partie de $\Ob\M$, $\uS$ la sous-catégorie pleine de $\M$ définie par $S$. On dit que $S$ est une \textit{partie accessible} de $\Ob\M$, ou encore que $\uS$ est une \textit{sous-catégorie “accessible dans $\M$”}, si
\printedpage{105}
% head: A. Grothendieck : Dérivateurs, Ch. XVIII, Édition des Universités de Montpellier II et Paris VII
% foot: Version 17 décembre 2025
\begin{enumerate}
\item[a)] $\uS$ est strictement pleine, i.e. pour $x$, $y$ deux objets isomorphes de $\M$, $x \in S$ implique $y \in S$,
\item[b)] $\uS$ est accessible,
\item[c)] le foncteur d'inclusion $\uS \longrightarrow \M$ est accessible.
\end{enumerate}
Soient $\pi_0$ un cardinal utile pour $\M$, i.e. tel que $\M \simeq \Ind_{\pi_0}(\M_{\pi_0})$, soit $\pi_1 = \sup(\pi_0, \card\red\M_{\pi_0})$. Soit d'autre part $\pi'_0$ un cardinal tel que

\msp{page 282}
$\uS$ (supposé accessible dans $\M$) soit $\pi_0$-\textit{accessible}, i.e. $\pi'_0$-utile pour $\uS$, et de plus tel % [sic] \pi_0-accessible -> \pi'_0-accessible, the following "i.e. \pi'_0-utile" refers to \pi'_0
que $f : \uS \hookrightarrow \M$ soit $\pi'_0$-accessible. Soit $\pi'_1 \geqslant \pi'_0$ un cardinal tel que $\uS_{\pi'_0} \subset \M_{\pi'_1}$, $\pi'_1 \geqslant \card\red\uS_{\pi'_0}$. Il résulte de 4.12.1 que pour tout cardinal
\[
\pi \geqslant \pi'_1 , \qquad \text{tel que } \pi^{\pi'_0} = \pi, \tag*{$(*)$}
\]
on doit avoir
\[
\uS_\pi \subset \M_\pi , \qquad \text{donc } \uS_\pi \subset \M_\pi \cap \uS,
\]
d'ailleurs, comme $\uS \hookrightarrow \M$ est $\pi'_0$-accessible et a fortiori $\pi$-accessible, et pleinement fidèle, il s'ensuit aussitôt que $\M_\pi \cap S \subset S_\pi$, d'où
\[
S_\pi = \M \cap S . % [sic] \M \cap S -> \M_\pi \cap S, the preceding lines concern \M_\pi \cap S
\]

\msp{page 283}
D'ailleurs, la condition $(*)$ implique que $\pi$ est utile pour $\uS$, donc
\[
\uS \xleftarrow{\sim} \Ind_\pi(\uS_\pi) .
\]
$\Big($ D'ailleurs, \textit{si} $\pi'_0 \geqslant \pi_0$, la relation $\pi^{\pi'_0} = \pi$ implique a fortiori $\pi^{\pi_0} = \pi$, et si $\pi'_1 \geqslant \pi_1$, on voit que $(*)$ implique que $\pi$ est utile aussi pour $\M$, i.e.
\[
\M \xleftarrow{\sim} \Ind_\pi(\M_\pi) .
\]
D'autre part, le foncteur d'inclusion $\uS \hookrightarrow \M$ se déduit essentiellement de $\uS_\pi \hookrightarrow \M_\pi$ en passant aux catégories $\Ind_\pi$. (\footnote[116]{\texttt{[Le contenu de cette parenthèse est]} “facultatif”.}). $\Big)$

\dgrblock{dgrb1068}{Accessibility of a full subcategory is a relative assertion}
\begin{dcomment}{exposition}{p282}{Accessibility of a full subcategory is a relative assertion}
Grothendieck now turns from functors to classes of objects specified by a property. His terminology ``accessible in $\mathcal M$'' includes three separate requirements: closure under isomorphism, accessibility of the full subcategory itself, and accessibility of its inclusion. In current language this is a replete, accessible, accessibly embedded full subcategory. Intrinsic accessibility alone does not control how its objects sit in the ambient category. This distinction will account for the extra bound needed in 4.13.2(c).

The equality of presentable objects on p.~282 is particularly useful. Let $i:\mathcal S\hookrightarrow\mathcal M$ be fully faithful and preserve $\lambda$-filtered colimits. If $s\in\mathcal S$ is $\lambda$-presentable in $\mathcal M$, then for every $\lambda$-filtered diagram $D$ in $\mathcal S$,
\[
 \mathcal S(s,\varinjlim D)
 \cong\mathcal M(is,\varinjlim iD)
 \cong\varinjlim\mathcal M(is,iD)
 \cong\varinjlim\mathcal S(s,D).
\]
Thus ambient presentability implies intrinsic presentability. The reverse implication requires the inclusion to preserve presentable objects as well as colimits. Uniformization supplies this after increasing the cardinal and bounding the ambient ranks of a small generating family. With those choices,
\[
 \mathcal S_\pi=\mathcal S\cap\mathcal M_\pi.
\]
The missing subscript on the last $\mathcal M$ in the printed formula is therefore substantive notation, not a new description of $S_\pi$.

The subsequent correspondence between retract-closed subcategories of $\mathcal M_\pi$ and subcategories generated from them by $\pi^+$-filtered colimits follows from one familiar retract argument. If a $\pi^+$-presentable object $s$ is a filtered colimit of objects $a_i$ in the chosen small subcategory, its identity factors through some $a_i\to s$, so $s$ is a retract of $a_i$. Conversely presentability is preserved by retracts. The natural closure in 4.13.3 is precisely the closure which makes this argument an equality of full subcategories. In 4.13.4 the printed words ``non pleinement fid\`ele'' contradict both 4.13.2 and this proof: the canonical functor under discussion is fully faithful.
\end{dcomment}


Ceci prouve entre autres le “il faut” dans la proposition qui suit :

\textbf{Proposition 4.13.2.} \textit{Soient $\M$ une catégorie accessible, $\pi_0$ un cardinal infini utile pour $\M$, $S$ une partie de} $\Ob\M$, \textit{$\uS$ la sous-catégorie pleine de $\M$ correspondante.}

\msp{page 284}
\begin{enumerate}
\item[\textit{a)}] \textit{Pour que $S$ soit accessible, il faut et il suffit qu'il existe un cardinal infini $\pi$ et une partie $S(\pi)$ de $\M_\pi$, tels que $S$ soit formé des objets de $\M$ qui sont isomorphes à un objet de la forme $\varinjlim\limits_{i\in I} x_i$, avec les $x_i$ dans $S(\pi)$ et $I$ grande devant $\pi$. En d'autres termes, il faut et il suffit qu'il existe un cardinal infini $\pi_1$ et une sous-catégorie pleine % [sic] \pi_1 -> \pi, the rest of the statement uses \pi
$\uS(\pi)$ de $\M_\pi$, tels que $\uS$ soit l'image essentielle du foncteur canonique pleinement fidèle}
\[
\Ind_\pi(\uS(\pi)) \longrightarrow \M \qquad \bigl(\xleftarrow{\sim} \Ind_\pi(\M)\ \text{(\footnotemark[117])}\bigr) , \tag*{(4.13.2.1)} % [sic] \Ind_\pi(\M) -> \Ind_\pi(\M_\pi), footnote 117 refers to \pi utile pour \M, i.e. \M \simeq \Ind_\pi(\M_\pi)
\]
\footnotetext[117]{si $\pi$ utile pour $\M$.}
\textit{prolongeant l'inclusion canonique $\uS(\pi) \hookrightarrow \M$ (cf. 4.12.3).}
\item[\textit{b)}] \printedpage{106}
% head: A. Grothendieck : Dérivateurs, Ch. XVIII, Édition des Universités de Montpellier II et Paris VII
% foot: Version 17 décembre 2025
\textit{Si $S$ est accessible, on peut trouver des cardinaux $\pi'_0$, $\pi'_1$ tels que $\pi'_1 \geqslant \pi'_0 \geqslant \pi_0$ (où $\pi_0$ est un cardinal donné utile pour $\M$)}
\msp{page 285}
\textit{et tels que tout cardinal $\pi \geqslant \pi'_1$ avec $\pi^{\pi'_0} = \pi$ satisfasse aux conditions de c), en prenant $S(\pi) = S \cap \Ob\M_\pi$, i.e. $\uS(\pi) = \uS \cap \M_\pi$. Sous ces conditions, $\pi$ est utile pour $\uS$} et \textit{pour $\M$, $S(\pi)$ est une sous-catégorie} strictement \textit{pleine de $\M_\pi$, stable dans $\M_\pi$ par facteurs directs (ou ce qui revient au même, $\M_\pi$ étant karoubienne, une sous-catégorie karoubienne de $\M_\pi$).}
\item[\textit{c)}] \textit{On peut prendre $\pi'_0$ le plus petit cardinal $\geqslant \pi_0$ utile pour $\uS$ et tel que $\uS \hookrightarrow \M$ soit $\pi'_0$-accessible, et $\pi'_1 = \sup(\pi'_0, \card\red\uS_{\pi'_0})$.}
\end{enumerate}

La partie b) ainsi que c) a déjà été prouvée, ce qui implique le “il faut” dans la partie a). Il reste à prouver le “il suffit” dans a). Mais l'hypothèse implique que
\[
\uS \xleftarrow{\sim} \Ind_\pi(S_\pi) ,
\]
ce qui implique que $\uS$ est accessible et même $\pi$-accessible, et elle implique aussi que $\uS$ est une sous-catégorie \textit{strictement} pleine de $\M$. Il reste à voir

\msp{page 286}
\texttt{[que]} le foncteur d'inclusion est accessible, en fait il est $\pi$-accessible, puisque (4.13.2.1) l'est (cf. 4.12.3).

\dgrblock{dgrb1069}{The ambient rank cannot be omitted from the threshold}
\begin{dcomment}{development}{p285}{The ambient rank cannot be omitted from the threshold}
The purpose of 4.13.2(c) is to make the eventual comparison $\mathcal S_\pi=\mathcal S\cap\mathcal M_\pi$ effective. Its proposed threshold omits data which the preceding argument explicitly used. We give a small counterexample, so that the correction can be made without weakening the surrounding theory.

\textbf{Counterexample to the threshold in 4.13.2(c).}\ours
Let $X$ be a set of cardinal greater than $2^{\aleph_0}$, let $\mathcal M=\mathcal P(X)$ ordered by inclusion, and let $\mathcal S$ be its full one-object subcategory $\{X\}$. Then $\mathcal M$ is $\aleph_0$-accessible in Grothendieck's convention, $\mathcal S$ is accessible at every infinite cardinal, and its inclusion preserves all nonempty filtered colimits. Nevertheless, for $\pi=2^{\aleph_0}$, the sole object of $\mathcal S_\pi$ does not belong to $\mathcal M_\pi$.

\textit{Proof.} Colimits in $\mathcal P(X)$ are unions. For every infinite $\rho$, a subset $Y\subseteq X$ is $\rho^+$-presentable precisely when $|Y|\leq\rho$. If this inequality holds and $Y\subseteq\bigcup_i Z_i$ for a $\rho^+$-directed family, select an index for each element of $Y$ and then a common upper bound. This proves preservation of the relevant colimits by the representable, whose values are either empty or singleton. If $|Y|>\rho$, the diagram of its subsets of cardinal at most $\rho$ is $\rho^+$-directed and has union $Y$, but none of its members contains $Y$. This proves the converse. Every subset is the union of its countable subsets, with an $\aleph_1$-directed indexing poset; hence $\mathcal M$ has the claimed accessibility. The corresponding assertions for the one-object category and its inclusion are immediate from constant diagrams. Finally $|X|>\pi$ gives the failure of ambient presentability.\hfill$\square$

Here the cardinal prescribed in 4.13.2(c) is $\pi'_0=\aleph_0$ and $\sup(\pi'_0,\operatorname{card}_{\mathrm{red}}\mathcal S_{\pi'_0})=\aleph_0$, whereas $\pi^{\pi'_0}=\pi$. Thus its displayed numerical conditions hold and its proposed conclusion fails. The repair is to retain a cardinal $\theta$ such that
\[
 i(\mathcal S_{\pi'_0})\subseteq\mathcal M_\theta
\]
and to take $\pi'_1$ at least $\theta$, as well as the bounds already listed. To guarantee simultaneous usefulness for $\mathcal M$, also retain the ambient bound $\pi_1$ from p.~281. The existential assertion of 4.13.2(b) is unaffected.

This example is our explicit diagnostic of the printed estimate, not a claim of priority. We checked the uniformization discussion in \citet[\S2.3]{Raptis2020}, which explicitly separates preservation of filtered colimits from preservation of presentable objects, and the standard limits treatment in \citet{Positselski2024Limits}; neither is being cited as a correction of this particular passage.
\end{dcomment}


\textbf{Définition 4.13.3.} Soient $\M$, $S \subset \Ob\M$, $\uS \subset \M$ comme dans 4.12.1, $\pi$ un cardinal infini. On dit que $S$ est une partie $\pi$-\textit{accessible} de $\Ob\M$, ou que $\uS$ est une sous-catégorie $\pi$-accessible dans $\M$, si $S$, $\pi$ satisfont les conditions de 4.13.3 a), i.e. $\uS$ est l'image essentielle % [sic] 4.13.3 a) -> 4.13.2 a), a self-reference of the definition, the conditions a) are those of Proposition 4.13.2
de $\Ind_\pi(S(\pi))$, où $S(\pi)$ est une sous-catégorie pleine de $\M_\pi$.

Notons que, quitte à renforcer $S(\pi)$ par $S'(\pi) = {\natural}$-clôture de $S(\pi)$ dans $\M_\pi$, on peut supposer que $S(\pi)$ est une sous-catégorie \textit{strictement pleine et karoubienne}

\msp{page 287}
de $\M_\pi$. On sait alors que, $S(\pi)$ étant karoubienne, $S(\pi) \longrightarrow (\underbrace{\Ind_\pi(S(\pi))}_{})_\pi$ est une équivalence de catégories. Comme $\Ind_\pi(S(\pi)) \simeq \uS$ et que $S(\pi)$ est strictement pleine dans $\M_\pi$, donc aussi dans $\M$ et a fortiori dans $\uS$, que l'on a
\[
\uS(\pi) = \uS_\pi , \qquad \text{sous-catégorie de } \uS \text{ des objets } \pi\text{-accessibles de } S. \tag*{(4.13.1)}
\]
Je dis qu'on aura aussi
\[
\uS(\pi) = \uS \cap \M_\pi . \tag*{(4.13.2)}
\]
Comme c est clair, il reste à voir que $\uS{\cap}\M_\pi \subset \uS(\pi)$, or il est immédiat que $\uS{\cap}\M_\pi \subset S_\pi$, % [sic] c -> an inclusion sign (\supset), a letter c printed where the reverse inclusion seems meant
et on conclut par (4.13.1). On trouve ainsi :

\textbf{Proposition 4.13.4.} \textit{Soit $\M$ une catégorie accessible, $\pi$ un cardinal infini. Soit $\Sigma_\pi$ l'ensemble des}

\msp{page 288}
\textit{sous-catégories strictement pleines $\uS(\pi)$ de $\M_\pi$ qui sont karoubiennes, ou ce qui revient au même, stables dans $\M_\pi$ par facteurs directs. Soit $\Sigma'_\pi$ l'ensemble des sous-catégories de}
\printedpage{107}
% head: A. Grothendieck : Dérivateurs, Ch. XVIII, Édition des Universités de Montpellier II et Paris VII
% foot: Version 17 décembre 2025
\textit{$\M$ qui sont $\pi$-accessibles dans $\M$. Considérons les applications}
\[
\begin{tikzcd}
\Sigma_\pi \arrow[r, shift left, "\alpha_\pi"] & \Sigma'_\pi \arrow[l, shift left, "\beta_\pi"]
\end{tikzcd}
\tag*{(4.13.4.1)}
\]
\textit{définies ainsi : Si $\uS(\pi) \in \Sigma_\pi$, $\alpha_\pi(\uS(\pi))$ est l'image essentielle du foncteur non pleinement fidèle} $\Ind_\pi(\uS(\pi)) \longrightarrow \M$. \textit{Si $\uS \in \Sigma'_\pi$, $\beta_\pi(\uS) = \uS \cap \M_\pi$. Les applications $\alpha_\pi$, $\beta_\pi$ sont des bijections, inverses l'une de l'autre.}

On n'a pas eu à supposer que $\pi$ soit utile pour $\M$, ni que $\pi$ soit minoré par un cardinal utile pour $\M$, i.e. qu'on ait $\pi \geqslant \pi_0$, où $\pi_0$ est le plus petit cardinal utile pour $\M$.

\msp{page 289}
\textbf{4.13.5.} Soit $\uS$ une sous-catégorie \textit{strictement pleine} de $\M$ (catégorie accessible), $\pi'_0$ un cardinal infini. Dire que $\uS$ est $\pi'_0$-accessible dans $\M$ signifie que les conditions suivantes sont satisfaites.
\begin{enumerate}
\item[1)] $\pi'_0$ est utile pour $\uS$, i.e. $\uS \xleftarrow{\sim} \Ind_{\pi'_0}(\uS_{\pi'_0})$.
\item[2)] $\uS \hookrightarrow \M$ est $\pi'_0$-accessible.
\item[3)] $\uS_{\pi'_0} \hookrightarrow \M_{\pi'_0}$ (qui implique que $S_{\pi'_0}$ est équivalente à une petite catégorie).
\end{enumerate}
Supposons ces conditions satisfaites, et soit $\pi$ un cardinal $> \pi'_0$. Est-ce que $\uS$ est $\pi$-accessible dans $\M$, i.e. les conditions similaires sont-elles satisfaites pour $\pi$ ? C'est évident pour 2°), mais risque d'être faux pour 1°, 2° \texttt{[plutôt 3°\,?]}. Mais faisons appel à 4.13.2 c). Posons
\[
\pi'_1 = \sup(\pi'_0, \card\red S_{\pi'_0}, \underbrace{\pi_0}_{\substack{\text{plus petit cardinal}\\ \text{infini utile}\\ \text{pour } \M}} ) .
\]

\msp{page 290}
Alors pour $\pi \geqslant \pi'_1$ tel que $\pi^{\pi'_0} = \pi$, $\uS$ est $\pi$-accessible dans $\M$ (4.13.2 b et c). Supposons p.ex.
\[
\pi'_0 \geqslant \pi_0
\]
(pratiquement ce sera toujours vérifié) et de plus
\[
\pi'_0 \geqslant \card\red\M_{\pi_0} ,
\]
i.e. supposons
\[
\pi'_0 \geqslant \sup(\pi_0, \card\red\M_{\pi_0}) = \pi_1 .
\]
Alors
\[
\pi'_1 = \sup(\pi'_0, \card\red S_{\pi'_0}) \leqslant \sup(\pi'_0, \card\red\M_{\pi'_0}) \leqslant 2^{({\pi'_0}^{\pi_0})}
\]
(4.7.7.1, p. 177). Donc on trouve ceci :

\textbf{Corollaire 4.13.6.} \textit{Soit $\pi'_0$ un cardinal $\geqslant \pi_1 = \sup(\pi_0, \card\red\M_{\pi_0})$ (\footnote[118]{\textbf{N.B.} $\pi_0$ est le plus petit cardinal utile pour $\M$.}), et soit $\uS$ une sous-catégorie $\pi'_0$-accessible dans $\M$. Alors pour tout cardinal $\pi$ tel que $\pi \geqslant 2^{({\pi'_0}^{\pi_0})}$, et tel que l'on ait $\pi^{\pi'_0} = \pi$ (i.e. $\pi$ de la forme $c^{\pi'_0}$, avec $c$ un cardinal $\geqslant 2^{({\pi'_0}^{\pi_0})}$), $\uS$ est $\pi$-accessible dans $\M$. (De plus, ces cardinaux sont réguliers pour $\uS$} et \textit{pour $\M$, donc $\uS \simeq \Ind_\pi(\uS_\pi)$, $\M \simeq \Ind_\pi(\M_\pi)$).}

(\footnote[119]{Enchaîner ici avec 4.13.23 etc.}, \footnote[120]{Mettre ici le lemme 4.16.7.11 signalé en marge p. 292.}).

\printedpage{108}
% head: A. Grothendieck : Dérivateurs, Ch. XVIII, Édition des Universités de Montpellier II et Paris VII
% foot: Version 17 décembre 2025
\msp{page 291}
\textbf{Proposition 4.13.7} (\footnote[121]{Le rejeter après 4.13.23.}). \textit{Soient $\M$ une catégorie accessible, $(S_\alpha)_{\alpha\in A}$ une famille de parties de} $\Ob\M$, \textit{indexée par un petit ensemble $A$, $(\uS_\alpha)_{\alpha\in A}$ la famille des sous-catégories accessibles dans $\M$ correspondante. Soit $S = \bigcap_{\alpha\in A} S_\alpha$, $\uS = \bigcap_{\alpha\in A} \uS_\alpha$, de sorte que $\uS$ est la sous-catégorie pleine de $\M$ définie par $S \subset$} $\Ob\M$. \textit{Ceci posé, si les $S_\alpha$ sont accessibles, il en est de même de $S$ – ou encore, si les $\uS_\alpha$ sont accessibles dans $\M$, il en est de même de $\uS$.}

\dgrblock{dgrb1070}{A homotopical use of accessible properties of arrows}
\begin{dcomment}{later}{p290}{A homotopical use of accessible properties of arrows}
The passage from accessible categories to accessible full subcategories is useful when the property under consideration belongs to a morphism. Regard a map as an object of the arrow category $\mathcal M^{[1]}$ and a commutative square as a morphism there. The phrase ``the weak equivalences form an accessible subcategory'' then has the same precise meaning as the relative accessibility discussed here: a suitable small family of weak equivalences presents every weak equivalence, and the inclusion preserves the specified filtered colimits.

A concrete later application is the treatment of weak equivalences in combinatorial model categories by \citet[Section~2 and Proposition~2.2]{RaptisRosicky2015}. After a functorial factorization, a map is a weak equivalence exactly when its fibration part is a trivial fibration. Their construction expresses the weak-equivalence category as a pseudopullback and obtains accessibility with an explicit enlarged rank. They then prove the sharper result that weak equivalences of simplicial sets form a finitely accessible category \citep[Corollary~5.1]{RaptisRosicky2015}. The finite conclusion is a further theorem, not a consequence of existence of some accessibility rank.

For the reader of these pages, the structural connection is the relevant one. A property is imposed by pulling back a suitable full subcategory along an accessible functor; several properties are imposed by intersection. The full subcategory retains its morphisms, so the resulting statement concerns coherent squares and their presentations, rather than only a set of admissible arrows. This is why bounds on the inclusion and on the images of presentable objects matter. It also explains why increasing the rank is a substantial part of the method: the homotopical construction may preserve filtered colimits long before it preserves the originally selected small objects. None of this identifies the ordinary cofinality arguments here with a theorem about homotopy colimits; their role is to provide the size control needed before those constructions are applied.
\end{dcomment}


Cela résulte du résultat plus précis :

\textbf{Corollaire 4.13.8} (\footnote[122]{Cf. reformulation plus raisonnable et plus correcte page 302, 302 bis.}). \textit{Avec les notations de 4.13.7 et 4.13.6, soit, pour tout $\alpha \in A$, $\pi'_\alpha$ un cardinal $\geqslant \pi_1$ tel que $S_\alpha$ soit $\pi'_\alpha$-accessible. Soit}
\[
\begin{aligned}
\pi'_0 &= \sup\nolimits_{\alpha\in A} \pi'_\alpha \\
\pi''_0 &= \sup(\card A, \sup\nolimits_\alpha 2^{({\pi'_\alpha}^{\pi_0})})
\end{aligned}
\]
\textit{et soit $\pi \geqslant \pi''_0$ un cardinal tel que $\pi^{\pi'_0} = \pi$, i.e. de la forme $c^{\pi'_0}$, avec $c \geqslant \pi'_0$. Alors $S$ est $\pi$-accessible (i.e. $\uS$ est $\pi$-accessible dans $\M$). En particulier, $S$ est $2^{\pi''_0}$-accessible.} (\footnote[123]{Ça a l'air un peu abracadabrant ! Il faut de plus supposer que $\pi$ n'est pas cardinal limite !}).

\msp{page 292}
\textsc{Démonstration}. Soit $\alpha \in A$. Comme $\pi'_\alpha \leqslant \pi'_0$ et $\pi^{\pi'_0} = \pi$, on a a fortiori
\[
\pi^{\pi'_\alpha} = \pi , \qquad \text{de plus } \pi \geqslant \pi''_0 \geqslant 2^{(\pi_\alpha^{\pi_0})} , % doubt resolved against PDF; see corrections.csv
\]
donc en vertu de 4.13.6, $\uS_\alpha$ est $\pi$-accessible dans $\M$. D'autre part, comme $\uS_\alpha$ est $\pi'_\alpha$-accessible, il est stable dans $\M$ par $\varinjlim$ grandes devant $\pi'_\alpha$, a fortiori par $\varinjlim$ grandes devant $\pi'_0$ ($\geqslant \pi'_\alpha$). D'ailleurs, $\pi^{\pi'_0} = \pi$ implique $\pi \geqslant 2^{\pi_0} > \pi'_0$. Donc 4.13.8 résulte du

\dgrblock{dgrb1071}{Intersections and the later theory of accessible subcategories}
\begin{dcomment}{later}{p291}{Intersections and the later theory of accessible subcategories}
Grothendieck wants a class of admissible properties of objects which is closed under imposing several conditions at once. Proposition 4.13.7 supplies exactly that operation. Its force is that a small intersection still has a small generating family, although the generating families for its separate conditions need not intersect in a useful way at the initial cardinal. The proof therefore raises the cardinal first and then constructs approximations which satisfy all the conditions.

This is now a standard stability theorem. In the higher categorical setting, \citet[Definition 5.4.7.8 and Proposition 5.4.7.10]{Lurie2009HTT} uses the same relative notion and proves closure under small intersections by a homotopy pullback involving the diagonal and a product of inclusions. The ordinary categorical counterpart belongs to the limit theory of accessible categories. A class-sized generalization, with the stronger embedding hypothesis needed there, is \citet[Corollary 3.3]{ChornyRosicky2012}. For explicit preservation ranks, \citet[Corollary 5.1 and Theorem 9.1]{Positselski2024Limits} retains a smaller ordinal-colimit hypothesis while controlling the larger presentability rank.

These formulations illuminate the calculation which follows. The objects of $\mathcal M_\pi/x$ are small approximations equipped with a map to $x$. If $x$ satisfies one of the properties, the approximations satisfying that property form a cofinal full subcategory. The problem is therefore to show that a small family of such cofinal subcategories has cofinal intersection. Cofinality alone does not suffice: in $\mathbb N$ the even and odd subsets are both cofinal and disjoint. Closure under a suitable family of colimits is the additional input which turns repeated approximation into simultaneous approximation.
\end{dcomment}


\textbf{Corollaire 4.13.9.} \textit{Soit $\pi > \pi_0$ un cardinal tel que $\uS_\alpha$ soit $\pi$-accessible dans $\M$ pour tout $\alpha \in A$, et tel que} $\card A \leqslant \pi$. \textit{Supposons de plus qu'il existe un cardinal $\pi'_0$ tel que $\pi_0 \leqslant \pi'_0 < \pi$, tel que les $\uS_\alpha$ soient stables dans $\M$ par $\varinjlim$ grandes devant $\pi'_0$. Alors $\uS = \bigcap_{\alpha\in A} \uS_\alpha$ est $\pi$-accessible dans $\M$.}

(\footnote[124]{Ce serait le moment de formuler ici le Lemme 4.16.7.11 (p. 488), qui devrait même remonter avec les premières généralités sur les parties $\pi$-accessibles dès avant 4.13.7 (immédiat après 4.13.6).}).

Posons, pour tout $\alpha \in A$,
\[
\begin{aligned}
\uS_\alpha(\pi) &\overset{\text{déf}}{=} \uS_\alpha \cap \M_\pi \\
\uS(\pi) &\overset{\text{déf}}{=} \uS \cap \M_\pi = \bigcap \uS_\alpha(\pi) .
\end{aligned}
\]
Chaque $\uS_\alpha(\pi)$ est stable dans $\M$ par $\varinjlim$ grandes devant $\pi'_0$ et de cardinal $\leqslant \pi$ (puisqu'il en est ainsi de $\uS_\alpha$ et de $\M_\pi$) (donc il en est de même de leur intersection

\msp{page 293}
$\uS(\pi)$. A fortiori, $\uS(\pi)$ est une sous-catégorie strictement pleine et karoubienne de $\M_\pi$.)

% ---- from der-ch18_p111-115.tex
\printedpage{109}
% head: A. Grothendieck : Dérivateurs, Ch. XVIII, Édition des Universités de Montpellier II et Paris VII

D'autre part, chaque $\underline{S}_\alpha$ étant stable dans $\mathcal{M}$ par $\varinjlim$ grandes devant $\pi'_0$, il en est de même de leur intersection $\underline{S}$. A fortiori, celle-ci est stable par $\varinjlim$ grandes devant $\pi$, donc le foncteur pleinement fidèle canonique $\Ind_\pi(\underline{S}(\pi))\longrightarrow\mathcal{M}$ se factorise par $S$ en
\[ \Ind_\pi\underline{S}(\pi) \incl \underline{S} \; , \]
il faut prouver que ce foncteur est essentiellement surjectif. Soit donc $x\in S=\Ob\underline{S}$, tout revient à prouver que
\[ \varinjlim_{i\in\underline{S}(\pi)/x} x_i \xrightarrow{\ \sim\ } x \; , \tag*{$(*)$} \]
et que $\underline{S}(\pi)/x$ est grande devant $\pi$.

Notons que, comme $\pi'_0\geqslant\pi_0$, et $\pi^{\pi'_0}=\pi$, on aura $\pi^{\pi_0}=\pi$, donc $\pi$ est cardinal utile pour $\mathcal{M}$, donc $I=\mathcal{M}_\pi/x$ est grande devant $x$, et on a % [sic] grande devant $x$ -> grande devant $\pi$, apparently intended (one speaks of being large before a cardinal)
\[ \varinjlim_{i\in I=\mathcal{M}_\pi/x} x_i \xrightarrow{\ \sim\ } x \; . \tag*{$(**)$} \]
Donc, si on admet que $\underline{S}(\pi)/x$ est grande devant $\pi$,
\msp{page 294}
le fait que $(*)$ soit un isomorphisme équivaut à dire que le foncteur canonique
\[ J \overset{\text{déf}}{=} \underline{S}(\pi)/x \incl \mathcal{M}_\pi \; = \; I \] % [sic] $\mathcal{M}_\pi = I$ -> $\mathcal{M}_\pi/x = I$, since $I=\mathcal{M}_\pi/x$ above and in the display for $J_\alpha$ below
est \textit{cofinal} (cf. 4.4.2, p. 79). Mais on voit tout de suite que ce foncteur est un foncteur d'inclusion d'une sous-catégorie strictement pleine (stable par $\varinjlim$ grandes devant $\pi'_0$, de cardinal $\leqslant\pi$).

Notons maintenant que la flèche $(**)$ similaire à $(*)$, avec $\underline{S}(\pi)$ remplacée par $\underline{S}_\alpha(\pi)$, est bel et bien un isomorphisme, puisque $x\in\Ob\underline{S}\subset\Ob\underline{S}_\alpha$, et que $\underline{S}_\alpha$ est $\pi$-accessible. Cela s'interprète par le fait que le foncteur d'inclusion
\[ J_\alpha \overset{\text{déf}}{=} \underline{S}_\alpha(\pi)/x \incl \mathcal{M}_\pi/x \; = \; I \]
est cofinal. C'est encore une inclusion d'une sous-catégorie $J_\alpha$ strictement pleine de $I$, stable par $\varinjlim$ grandes devant $\pi'_0$ et de cardinal $\leqslant\pi$, et on a
\[ J \; = \; \bigcap J_\alpha \; . \]

Donc la condition voulue résultera du
\msp{page 295}

% marginal hand-drawn mark (manuscript) in the left margin beside the following lemma
\textbf{Lemme 4.13.10} (\footnote[125]{Il faut supposer que $\pi$ ne soit pas un cardinal limite, pour pouvoir conclure que $I$ est stable par $\varinjlim$ de type $I_{\alpha_\pi}$, et les $J_r$ stables dans $I$ par ces limites.},\footnote[126]{Cf. forme plus générale du lemme dans 4.13.10.2, p. 302.}). \textit{Soient $\pi'_0$ et $\pi>\pi'_0$ des cardinaux infinis, $I$ une catégorie équivalente à une petite catégorie, $(J_\alpha)_{\alpha\in A}$ une famille de sous-catégories $J_\alpha$ de $I$. On suppose :}
\begin{itemize}
\item[\cercle{a}] \textit{$I$ est grande devant $\pi$ (}\footnote[127]{Il suffit que $I$ soit grande devant tout $\pi'<\pi$.}\textit{).} % footnote 127 printed at the foot of p. 110
\item[\cercle{b}] \textit{$I$ est stable par $\varinjlim$ de cardinal $\leqslant\pi$, et grandes devant $\pi'_0$ (}\footnote[128]{Il suffit ici la stabilité par $\varinjlim$ de type $I_{\alpha_\pi}$, cf. page 302 [?].}\textit{).} % footnote 128 printed at the foot of p. 110
\end{itemize}
% foot: Version 17 décembre 2025

\printedpage{110}
% head: A. Grothendieck : Dérivateurs, Ch. XVIII, Édition des Universités de Montpellier II et Paris VII
\begin{itemize}
\item[\cercle{c}] \textit{$\forall\ r\in A$, $J_r$ est stable dans $I$ par les limites inductives envisagés dans b) (en particulier, $J_r$ est une sous-catégorie} strictement pleine \textit{de $I$).} % [sic] envisagés -> envisagées, agreement with « limites »
\item[\cercle{d}] \textit{$\forall\, r\in A$, $J_r$ est une sous-catégorie} cofinale \textit{de $I$ (}grande devant $\pi$\textit{).}
\item[\cercle{e}] $\card A\leqslant\pi$.
\end{itemize}
\textit{Soit $J=\bigcap_{r\in A}J_r$. Alors $J$ est encore cofinale dans $I$, et grande devant $\pi$ (et de plus stable dans $I$ par $\varinjlim$ de cardinal $\leqslant\pi$ et grandes devant $\pi'_0$).}

En d'autres termes, si $I$ est grande devant $\pi$, l'intersection $J$ d'une famille $(J_\alpha)_{\alpha\in A}$ de sous-catégories strictement pleines cofinales et grandes devant $\pi$, est encore
\msp{page 296}
cofinale et grande devant $\pi$, \textit{pourvu que} $\card A\leqslant\pi$, et que l'on postule une propriété de stabilité de $I$ et des $J_\alpha$ dans $I$, pour les limites inductives de cardinal $\leqslant\pi$ et grandes devant un cardinal convenable $\pi'_0<\pi$ donné.

\textbf{4.13.10.1.} Notons d'ailleurs que si $f : J\longrightarrow I$ est un foncteur cofinal \textit{pleinement fidèle}, avec $I$ grande devant $\pi$, alors $J$ est automatiquement grande devant $\pi$. Cela résulte p.ex. de 4.9.6.3 2° (p. 225) par le critère $\mathrm{c}_\pi$ (vu que pour $f$ pleinement fidèle, la condition F $3_\pi$ est tautologiquement satisfaite). On peut aussi le voir ainsi : Soit $B$ un ensemble avec $\card B\leqslant\pi$, il faut prouver que $J\longrightarrow J^B$ est cofinal, or on a
\[
\begin{tikzcd}
J \arrow[r, "\text{cof.\,?}"] \arrow[d, "\text{cof.}"'] & J^B \arrow[d, "\text{cof.}"] \\
I \arrow[r, "\text{cof.}"] & I^B \; ,
\end{tikzcd}
\]
où trois des quatre flèches du carré sont cofinales, $J^B\longrightarrow I^B$ l'étant comme produit de foncteurs cofinaux entre catégories \textit{filtrantes} (cf. 4.9.6.5, page 227) -- on sait en effet déjà que $J$ est filtrante.

\msp{page 297}
Donc le composé
\[ J \longrightarrow J^B \xrightarrow{\text{cof.}} I^B \]
est cofinal. Mais comme le deuxième foncteur $J^B\longrightarrow I^B$ est pleinement fidèle, il s'ensuit aussitôt que $J\longrightarrow J^B$ est cofinal (car $\xi\backslash J$ pour $\xi\in\Ob J^B$ est le même, à isomorphisme près, qu'on considère $\xi$ comme objet de $J^B$, ou de $I^B$).

Ainsi dans l'énoncé de la condition d) du lemme, il est inutile de spécifier que $J_\alpha$ est grande devant $\pi$, et de même dans la conclusion sur $J$.

Comme $J\incl I$ est pleinement fidèle, pour voir que ce foncteur est cofinal, il suffit de voir qu'il satisfait F 1, i.e. que pour tout $i\in\Ob I$, existe $j\in J$ avec $i\longrightarrow j$. On va construire $j$ comme limite d'un système inductif ordinal d'objets $i_\alpha$ de $I$, indexé par les ordinaux $\alpha$ tels que $\card\alpha<\pi$. Désignons par $\alpha_\pi$ le plus petite ordinal tel que % [sic] le plus petite ordinal -> le plus petit ordinal, agreement with « ordinal »
% foot: Version 17 décembre 2025
\printedpage{111}
% head: A. Grothendieck : Dérivateurs, Ch. XVIII, Édition des Universités de Montpellier II et Paris VII
\msp{page 298}
% marginal hand-drawn mark (manuscript) in the left margin beside the following two lines
$\card\alpha_\pi=\pi$. L'ensemble ordonné $I_{\alpha_\pi}$ des ordinaux $\alpha<\alpha_\pi$ étant grande devant $\pi'_0$ (\footnote[129]{faux si $\pi$ est un cardinal limite.}) et de cardinal $\pi$, on pourra donc prendre la limite $\varinjlim\limits_{\alpha\in I_{\alpha_0}} i_\alpha$ dans $I$. % [sic] $I_{\alpha_0}$ -> $I_{\alpha_\pi}$, the index set just defined is $I_{\alpha_\pi}$

Construction de $(i_\alpha)_{\alpha\in I_{\alpha_0}}$. Soit % [sic] $I_{\alpha_0}$ -> $I_{\alpha_\pi}$, the index set just defined is $I_{\alpha_\pi}$
\[ \alpha_\pi \; = \; \text{plus petit ordinal } \alpha \text{ tel que } \card\alpha=\pi, \tag*{(1)} \]
donc c'est l'ordinal caractérisé par
\[ \left\{
\begin{aligned}
&\card\alpha_\pi \; = \; \pi \\
&\alpha<\alpha_\pi \;\; \Longrightarrow \;\; \card\alpha \; < \; \pi \; ,
\end{aligned}
\right. \tag*{(2)} \]
% marginal hand-drawn mark (manuscript) in the left margin beside the following paragraph
ou encore par le fait que l'ensemble ordonné $I_{\alpha_\pi}$ des ordinaux $\alpha<\alpha_\pi$ est \textit{parfaitement bien ordonné}. Ainsi, $I_{\alpha_\pi}$ \textit{est grande devant $\pi'$ pour tout cardinal $\pi'<\pi$} (\footnote[130]{Faux si $\pi$ est cardinal limite.}). Dans les conditions b) c) de 4.3.10, il suffira de supposer $I$ stable par $\varinjlim$ de type $I_{\alpha_\pi}$, et les $J_r$ ($r\in A$) stables dans $I$ par ce type de limites. % [sic] 4.3.10 -> 4.13.10, the conditions b) c) are those of Lemme 4.13.10

\textbf{Lemme 4.13.10.1.} \textit{Soit $\pi$ un cardinal infini, $A$ un ensemble non vide tel que} $\card A\leqslant\pi$. \textit{Alors}
\msp{page 299}
\textit{il existe une application}
\[ \varphi \; : \; I_{\alpha_\pi} \longrightarrow A \]
\textit{telle que pour tout $r\in A$, $\varphi^{-1}(\{r\})$ soit une partie cofinale de $I_{\alpha_\pi}$, ou ce que est plus précis,} $\card\varphi^{-1}(\{r\})=\pi$. % [sic] ce que est -> ce qui est, relative pronoun subject

\dgrblock{dgrb1072}{The ordinal calculation needs cofinality, not cardinality alone}
\begin{dcomment}{exposition}{p298}{The ordinal calculation needs cofinality, not cardinality alone}
Grothendieck chooses an ordinal long enough to revisit every condition, and then takes its colimit. Two distinct size arguments enter this construction. An ordinal $\beta$ below the initial ordinal of cardinal $\pi$ has cardinal less than $\pi$, so filteredness supplies a cocone at that stage. By contrast, the whole initial ordinal $\pi$ is large before $\rho$ precisely when
\[
 \operatorname{cf}(\pi)>\rho.
\]
Indeed a subset of size at most $\rho$ has an upper bound exactly when it is not cofinal. The statement on p.~298 that the ordinal is large before every cardinal below $\pi$ is therefore equivalent to regularity of $\pi$. For $\pi=\aleph_\omega$, the countable sequence of initial ordinals $\aleph_n$ is a counterexample. The marginal warning is justified, but ``limit cardinal'' is too broad: a regular limit cardinal would cause no such failure.

There is a useful qualification for 4.13.8. Its equation $\pi^\rho=\pi$ already forces $\operatorname{cf}(\pi)>\rho$. To see this directly, suppose $\delta=\operatorname{cf}(\pi)\leq\rho$ and enumerate all maps $\delta\to\pi$ by indices below $\pi$. Partition those indices into $\delta$ sets $E_\xi$, each of cardinal less than $\pi$. At coordinate $\xi$, choose a value different from the values of all maps indexed by $E_\xi$. The resulting map differs from every map in the enumeration, a contradiction. Thus the particular filtration needed in 4.13.8 follows from its fixed-point equation; the extra successor-cardinal requirement in footnote 123 is unnecessary there.

The ordinal lemma restated in 4.13.10.2 has an independent valid proof under its explicitly stated chain-colimit hypothesis. Partition the initial ordinal $\pi$ into at most $\pi$ subsets, each of cardinal $\pi$, and assign one part to each $J_r$. Every such part is cofinal and, with its induced order, has order type $\pi$. Recursively choose a compatible chain which, at the assigned stages, belongs to the corresponding $J_r$; at a limit stage take a cocone, then extend into the required cofinal subcategory. Start by extending the prescribed object into the first assigned subcategory, rather than requiring an arbitrary $i_0=i$ already to belong to it. The final colimit is the colimit of its restriction to each assigned part, hence belongs to every $J_r$. This proves cofinality. If $I$ is only large before every cardinal below $\pi$, the full subcategory inherits precisely those filtration bounds; inheritance of the stronger bound $\pi$ requires that bound for $I$ itself.
\end{dcomment}


En effet, il suffit de considérer un ensemble $A'$ de cardinal $\pi$ (p.ex. $A'=I_{\alpha_\pi}$), l'application
\[ \psi \; = \; \mathrm{pr}_1 \; : \; A\times A'\longrightarrow A \; , \]
et de se donner une bijection
\[ A\times A' \xrightarrow{\ \sim\ } I_{\alpha_\pi} \]
laquelle existe puisque
\[ \card\,(A\times A') \; = \; \card A\cdot\card A' \; = \; \card A' \; = \; \pi \; = \; \card I_{\alpha_\pi} \; , \]
puisque $0\neq\card A\leqslant\card A'=\pi$ et $\pi$ infini. On suppose choisie $\varphi$ comme dans 4.13.10.1, et on construira le système inductif $(i_\alpha)_{\alpha\in I_{\alpha_\pi}}$ de telle façon qu'on ait
\[ i_\alpha \; \in \; J_{\varphi(\alpha)} \; . \]

Posant
\[ j \; = \; \varinjlim_{\alpha\in I_{\alpha_\pi}} i_\alpha \; , \]
% foot: Version 17 décembre 2025
\printedpage{112}
% head: A. Grothendieck : Dérivateurs, Ch. XVIII, Édition des Universités de Montpellier II et Paris VII
\msp{page 300}
pour tout $r$, comme
\[ I_r \overset{\text{déf}}{=} \varphi^{-1}(r) \; \subset \; I_{\alpha_\pi} \]
est cofinal dans $I_{\alpha_\pi}$, on aura donc
\[ j \; = \; \varinjlim_{\alpha\in I_r} i_\alpha \; , \]
et en vertu de $(*)$, les $i_\alpha$ pour $\alpha\in I_r$ sont dans $J_r$, donc par passage à la limite, utilisant l'hypothèse c), on trouve que $j\in J_r$, donc
\[ j \; \in \; \bigcap_{r\in A} J_r \; = \; J \; , \]
et comme on choisira
\[ i_0 \; = \; i \; , \tag*{$(**)$} \]
on aura bien trouvé $i\longrightarrow j$ avec $i\in J$. Donc il suffit de construire le système inductif $(i_\alpha)$ satisfaisant $(*)$ et $(**)$. % [sic] avec $i\in J$ -> avec $j\in J$, it is $j$ that was shown to lie in $J$

Si le système inductif est construit jusqu'au cran $i_\alpha\in\Ob J_{\varphi(\alpha)}$, on continue par
\[ i_\alpha \longrightarrow i_{\alpha+1} \; \in \; \Ob J_{\varphi(\alpha+1)} \; , \]
ce qui est possible puisque $J_{\varphi(\alpha+1)}$ est
\msp{page 301}
cofinale. Donc il reste à construire $i_\alpha$ pour $\alpha$ ordinal limite. Or l'ensemble des $i_\beta$ pour $\beta<\alpha$ est de cardinal $<\pi$, et comme par a) $I$ est grande devant $\pi$ (il suffit qu'il le soit devant tout cardinal $<\pi$), il existe un $i\in\Ob I$ qui majore tous les $i_\beta$. Choisissons des
\[ i_\beta \xrightarrow{v_\beta} i \; . \]

On n'aura pas a priori commutativité dans
\[
\begin{tikzcd}
i_{\beta'} \arrow[r, "\tau_{\beta,\beta'}"] \arrow[rr, bend right=15, "v_{\beta'}"'] & i_\beta \arrow[r, "v_\beta"] & i \; ,
\end{tikzcd}
\] % doubt resolved against PDF; see corrections.csv
mais comme l'ensemble des double-flèches $i_{\beta'}\mathrel{\overset{v_\beta\tau_{\beta,\beta'}}{\underset{v_{\beta'}}{\rightrightarrows}}} i$ est de cardinal $<\pi$, $I$ étant grande devant $\pi$, il existe $i\longrightarrow i'$ qui les égalise tous. Comme $J_{\varphi(\alpha)}$ est cofinal dans $I$, il existe $i''\in J_{\varphi(\alpha)}$, et $i\longrightarrow i''$. On prend $i_\alpha=i''$, avec les homomorphismes composés
\[
\begin{tikzcd}
i_\beta \arrow[r, "v_\beta"] \arrow[rrr, dotted, bend right=12, "\tau_{\alpha,\beta}"'] & i \arrow[r] & i' \arrow[r] & i''
\end{tikzcd}
\] % doubt resolved against PDF; see corrections.csv
% foot: Version 17 décembre 2025
\printedpage{113}
% head: A. Grothendieck : Dérivateurs, Ch. XVIII, Édition des Universités de Montpellier II et Paris VII
\msp{page 302}
comme morphismes de transition. Cela achève la démonstration de 4.13.10 et par là de 4.13.9, 4.13.8.

\textbf{Remarque 4.13.10.2.} Voici ce qu'on a eu à utiliser vraiment, pour établir la conclusion de 4.13.10, en termes de $I$, $\pi$, $(J_r)_{r\in A}$ (et sans introduire de $\pi'_0$).
\begin{itemize}
\item[\cercle{a$'$}] $I$ est grande devant tout cardinal $\pi'<\pi$.
\item[\cercle{b$'$}] $I$ est stable par $\varinjlim$ de type $I_{\alpha_\pi}$.
\item[\cercle{c$'$}] $\forall\, r\in A$, $J_r$ est stable dans $I$ par les $\varinjlim$ de type $I_{\alpha_\pi}$.
\item[\cercle{d$'$}] $\forall\, r\in A$, $J_r$ est cofinale dans $I$.
\item[\cercle{e$'$}] $\card A\leqslant\pi$.
\end{itemize}
(\textbf{N.B.} d$'$) et e$'$) ne sont autres que d) et e) de 4.3.10, tandis que a$'$), b$'$), c$'$) sont des versions affaiblies de a), b), c) respectivement.) % [sic] 4.3.10 -> 4.13.10, the conditions d) e) are those of Lemme 4.13.10

\dgrblock{dgrb1073}{Simultaneous refinement removes the regularity restriction}
\begin{dcomment}{development}{p302:Voici ce qu'on a eu à utiliser vraiment}{Simultaneous refinement removes the regularity restriction}
The intersection argument can be arranged so that the number of conditions and the length of the construction play separate roles. This gives the conclusion of 4.13.10 for the hypotheses in its main text, without imposing the extra cofinality restriction introduced in its footnotes.

\textbf{Simultaneous cofinal intersection lemma.}\ours
Let $\rho<\pi$ be infinite cardinals. Suppose $I$ is $\pi^+$-filtered and admits all $\rho^+$-filtered colimits indexed by categories with at most $\pi$ arrows. Let $(J_r)_{r\in A}$, where $|A|\leq\pi$, be replete full cofinal subcategories closed under these colimits in $I$. Then $\bigcap_rJ_r$ is cofinal in $I$, is $\pi^+$-filtered, and is closed under the same colimits.

\textit{Proof.} The empty family is immediate. Put $\kappa=\rho^+$, a regular cardinal with $\kappa\leq\pi$, and fix $x_0\in I$. Construct a chain $(x_\beta)_{\beta<\kappa}$. After choosing $x_\beta$, choose, for every $r$, an arrow
\[
 a_{\beta r}:x_\beta\longrightarrow y_{\beta r},\qquad y_{\beta r}\in J_r.
\]
The star consisting of these arrows has at most $\pi$ arrows. Filteredness therefore supplies a cocone with vertex $x_{\beta+1}$ and arrows $b_{\beta r}:y_{\beta r}\to x_{\beta+1}$ satisfying
\[
 b_{\beta r}a_{\beta r}=t_{\beta+1,\beta}
 \quad\text{for every }r.
\]
At a limit $\gamma<\kappa$, choose a cocone on the entire already constructed $x$-chain; its size is less than $\kappa\leq\pi$. The resulting transition arrows $t_{\gamma\beta}$ are compatible by construction. No intermediate colimit is required.

For a fixed $r$ define, when $\beta<\gamma$,
\[
 y_{\beta r}\xrightarrow{b_{\beta r}}x_{\beta+1}
 \xrightarrow{t_{\gamma,\beta+1}}x_\gamma
 \xrightarrow{a_{\gamma r}}y_{\gamma r}.
\]
These maps compose correctly: the equality $b_{\gamma r}a_{\gamma r}=t_{\gamma+1,\gamma}$ reduces composition to the transition identity for the $x$-chain. The maps $a$ and $b$ consequently give mutually inverse maps between the colimits of the two chains; their composites are the maps induced by a shift of the indexing ordinal, hence the identities on those colimits. The ordinal $\kappa$ is $\rho^+$-filtered and has at most $\pi$ arrows. Thus $x=\varinjlim_{\beta<\kappa}x_\beta$ exists, and closure of $J_r$ shows that the isomorphic colimit of the $y_{\beta r}$ lies in $J_r$. Repleteness gives $x\in J_r$ for every $r$. The map $x_0\to x$ proves the first cofinality condition.

For completeness, a full subcategory of a filtered category which receives a map from every ambient object is cofinal: any two objects in a comma category have a common successor by taking a cocone in the ambient category and then extending into the subcategory; parallel arrows are equalized in the same way. The same argument for diagrams with at most $\pi$ arrows shows that the intersection is $\pi^+$-filtered. Closure under the specified colimits is inherited from each $J_r$.\hfill$\square$

Applied to $I=\mathcal M_\pi/x$, this also repairs 4.13.9 without its unstated fixed-point equation, for a nonempty family. Choose one $r$: since $x\in\mathcal S_r$ and $\mathcal S_r$ is $\pi$-accessible in $\mathcal M$, a $\pi^+$-filtered presentation of $x$ inside $\mathcal S_r$ uses objects presentable in $\mathcal M$. Factoring any at-most-$\pi$ diagram of maps from presentable objects through a common stage shows directly that $\mathcal M_\pi/x$ is $\pi^+$-filtered and presents $x$. Global usefulness of $\pi$ for $\mathcal M$ is therefore unnecessary for these particular slices. For an empty family the fixed-rank assertion requires $\pi$ useful for $\mathcal M$, since the intersection is then $\mathcal M$ itself.

We checked the intersection and pseudopullback formulations in \citet{Lurie2009HTT,ChornyRosicky2012,Positselski2024Limits}; we did not find this exact simultaneous-star repair of the printed ordinal argument there. The lemma is recorded as an explicit completion of Grothendieck's method, without a priority claim. It uses the full $\pi^+$-filteredness of $I$ for the stars. The weaker hypothesis suggested in footnote 127 does not support this proof when $|A|=\pi$.
\end{dcomment}


\textbf{4.13.8 (reformulation).} \textit{Soit pour $r\in A$, $\pi''_r\geqslant\pi'_r\geqslant\pi_0$ (= cardinal utile pour $\mathcal{M}$ donné) comme spécifié dans 4.13.2 b) (et c), soit $\pi'=\sup_{r\in A}\pi'_r$, $\pi''=\sup_{r\in A}\pi''_r$, de sorte que pour tout cardinal $\pi$ avec $\pi^{\pi'_0}=\pi\geqslant\pi''$, $\underline{S}_r$ soit $\pi$-accessible dans $\mathcal{M}$ (et $\pi$ utile pour $\underline{S}_r$} et \textit{pour $\mathcal{M}$). Soit $c$ un cardinal} non limite \textit{tel que $c\geqslant$} $\card A$ \textit{et $c>\pi'$ (p.ex. le successeur de} $\sup(\card A,\pi')$\textit{). Prenons $\pi$ tel que}
\[ \pi^{\pi'} \; = \; \pi \; \geqslant \; \sup(\pi'',c) \; . \]
\textit{Alors $\underline{S}=\bigcap_{r\in A}\underline{S}_r$ est $\pi$-accessible.}

En effet, si $\alpha_c$ est le premier ordinal de cardinal $c$, comme $c$ n'est pas limite, $I_{\alpha_c}$ est grande devant tout cardinal $c'<c$, en particulier devant $\pi'$. Or pour tout $x\in\Ob\mathcal{M}$, $I=\mathcal{M}_\pi/x$ et les $J_r=\underline{S}_r/x$ sont stables par $\varinjlim$ grandes devant $\pi'$ et de cardinal $\leqslant\pi$,
\msp{page 302 bis}
et les $J_r$ stables dans $I$ pour ces limites, donc il en est en particulier ainsi pour les $\varinjlim$ de type $I_{\alpha_c}$. On peut donc appliquer le résultat 4.13.10.2 pour conclure que si les $J_r$ sont cofinaux dans $I$ (qui est grande devant $\pi$, $\pi$ étant utile pour $\mathcal{M}$), alors $J=\underline{S}/x$ l'est également. Cela implique alors, comme on a vu, que $\underline{S}$ est $\pi$-accessible dans $\mathcal{M}$.

\textbf{4.13.10.3.} On peut reformuler la reformulation précédente, en utilisant la terminologie introduite plus bas (4.13.23) :

\textit{Soit $c$ un cardinal, $\alpha_c$ le plus petit ordinal $\alpha$ tel que} $\card\alpha=c$, \textit{$\pi$ un cardinal $\geqslant c$. On suppose}
\begin{itemize}
\item[$(*)$] \textit{$\pi$ est utile pour $\mathcal{M}$, et $\mathcal{M}$ est stable par $\varinjlim$ de type $I_{\alpha_c}$. (Cette deuxième condition est le cas p.ex. si $c$ est} non limite \textit{et $c>\pi_0$, $\pi_0$ étant un cardinal tel que $\mathcal{M}$ satisfait \textit{2}$_{\pi_0}$, p.ex. un cardinal} utile \textit{pour $\mathcal{M}$ (}\footnote[131]{Car alors $I_{\alpha_c}$ est grande devant $\pi_0$.}\textit{)).}
\end{itemize}
\textit{Alors l'ensemble des sous-catégories $\underline{S}$ de $\mathcal{M}$ qui sont $\alpha_c$-$\pi$-accessible dans $\mathcal{M}$ (i.e. qui sont $\pi$-accessibles et stables dans $\mathcal{M}$ par $\varinjlim$ de type $I_{\alpha_c}$) est stable par intersection de cardinal $\leqslant c$. Plus généralement, si $\Phi$ est un ensemble de catégories telles que} % [sic] accessible -> accessibles, agreement with « sous-catégories »
\begin{itemize}
\item[\textit{1°)}] \textit{$I\in\Phi\Longrightarrow$} $\card I\leqslant\pi$ \textit{et}
\item[\textit{2°)}] \textit{$I_{\alpha_c}\in\Phi$,}
\end{itemize}
\textit{alors l'ensemble des sous-catégories de $\mathcal{M}$ qui sont $\Phi$-$\pi$-accessibles dans $\mathcal{M}$, est stable par intersection de cardinal $\leqslant c$.}
% foot: Version 17 décembre 2025

% ---- from der-ch18_p116-120.tex
\printedpage{114}
% head: A. Grothendieck : Dérivateurs, Ch. XVIII, Édition des Universités de Montpellier II et Paris VII
% foot: Version 17 décembre 2025

\msp{page 303}\noindent\textbf{Définition 4.13.11} (\footnote[132]{Refondre avec 4.13.23 etc.}). Soient $\mathcal{M}$ une catégorie accessible, $\pi_0$ un cardinal utile pour $\mathcal{M}$, $\pi > \pi_0$ un cardinal. Une catégorie $I$ est dite $\pi_0$-$\pi$-\textit{adaptée} si
\begin{enumerate}
\item[a)] $I$ est équivalente à une petite catégorie,
\item[b)] $I$ est grande devant $\pi$, et
\item[c)] $I$ est stable par $\varinjlim$ de cardinal $\leqslant \pi$, et grandes devant $\pi_0$ (\footnote[133]{\textbf{N.B.} Si $\pi = \pi_0$, la condition c) équivaut à dire que $I$ est stable par facteurs directs.}).
\end{enumerate}
\noindent Un système inductif $(x_i)_{i\in I}$ de type $I$ dans $\mathcal{M}$ est dit $\pi_0$-$\pi$-\textit{adapté}, si $I$ est $\pi_0$-$\pi$-adaptée, i.e. satisfait les conditions a) b) c), et si de plus on a
\begin{enumerate}
\item[d)] le foncteur $i \longmapsto x_i : I \longrightarrow \mathcal{M}$ commute aux $\varinjlim$ (grandes devant $\pi_0$ et de cardinal $\leqslant \pi$) envisagées dans c), et
\item[e)] les $x_i$ sont dans $\mathcal{M}_\pi$.
\end{enumerate}

\noindent\textbf{Proposition 4.13.12.} \textit{Soient $\mathcal{M}$ une catégorie $\pi_0$-accessible, et $\pi > \pi_0$ un cardinal utile pour $\mathcal{M}$. Alors tout $x \in \Ob\mathcal{M}$ admet une représentation comme $x \xleftarrow{\ \sim\ } \ilim_{j\in J} x_j$, où le système inductif $(x_j)_{j\in J}$ dans $\mathcal{M}$ est $\pi_0$-$\pi$-adapté. De façon précise, on peut prendre $J = \mathcal{M}_\pi/x$.}

\msp{page 304}\noindent En effet, la vérification des propriétés a) à d) pour le système inductif canonique, de $\varinjlim$ isomorphe à $x$, est immédiate.

La proposition précédente ne dit pas si on peut prendre pour $J$ un ensemble ordonné. Dans ce cas, la condition c) signifie que pour toute partie $J'$ de $J$ grande devant $\pi_0$ et de cardinal $\leqslant \pi$, $J$ admet une borne supérieure $i_J$ dans $J$, et la condition d) que l'on a alors $x_{i_J} \xleftarrow{\ \sim\ } \ilim_{i\in J'} x_i$.

L'idée naturelle, pour construire un système inductif $\pi_0$-$\pi$-adapté $(x_j)_{j\in J}$, avec $J$ ordonnée, en partant de celui qu'on a canoniquement, avec $J = \mathcal{M}_\pi/x$, est de prendre pour $I$ l'ensemble, ordonné par inclusion, des sous-catégories $J'$ de $J_0$ (sous-catégorie pleine petite de $J$ équivalente à $J$) qui sont de cardinal $\leqslant \pi$ et grandes devant $\pi_0$. (Comparer (4.7.3.3), page 165). Cet ensemble ordonné est bel et bien (petit et) \textit{grand devant} $\pi$, si on suppose que $\pi^{\pi_0} = \pi$ (cf. démonstration 4.7.5, p. 168-171).

\msp{page 305}\noindent\textbf{Lemme 4.13.13.} \textit{Soient $\pi_0$, $\pi > \pi_0$ des cardinaux infinis tels que $\pi^{\pi_0} = \pi$, et $J$ une petite catégorie, grande devant $\pi$. Soit $I$ l'ensemble des sous-catégories $J'$ de $J$ qui sont grandes devant $\pi_0$, et de cardinal $\leqslant \pi$. Alors $I$ est $\pi_0$-$\pi$-adaptée.}

Il reste à voir que pour tout ensemble ordonné $A$ grand devant $\pi_0$ et de cardinal $\leqslant \pi$, si $r \longmapsto J_r$ est une famille croissante d'éléments de $I$ indexée par $A$, alors
\[
J \ \overset{\text{déf}}{=} \ \bigcup_{r\in A} J_r \ = \ \varinjlim_{r\in A} J_r
\]
est dans $I$, i.e. est une sous-catégorie de $I$ de cardinal $\leqslant \pi$ (ce qui est clair) et \textit{grande devant} $\pi_0$. Cela va résulter du

\printedpage{115}
% head: A. Grothendieck : Dérivateurs, Ch. XVIII, Édition des Universités de Montpellier II et Paris VII
% foot: Version 17 décembre 2025

\noindent\textbf{Lemme 4.13.13.1.} \textit{Soit $(J_r)_{r\in A}$ un système inductif de catégories grandes devant $\pi_0$, indexé par un ensemble ordonné $A$ grand devant $\pi_0$. Alors $J = \ilim_{r\in A} J_r$ est grande devant $\pi_0$.}

\msp{page 306}\noindent Il faut voir que si $B$ est un ensemble de cardinal $\leqslant \pi_0$, alors le foncteur diagonal
\[
\delta \ : \ J \longrightarrow J^B \tag*{$(*)$}
\]
est cofinal. Or comme $A$ est grande devant $\pi_0$, $\operatorname{card} B \leqslant \pi_0$, les produits indexés par $B$ commutent aux $\varinjlim$ de type $A$ (SGA 4 I 9.8 appliqué à Cat et à $\pi_0$). Donc la flèche $(*)$ s'interprète comme limite inductive des flèches similaires
\[
\delta_r \ : \ J_r \longrightarrow J_r^B \ , \tag*{$(**)$}
\]
lesquelles sont cofinales, donc $\delta$ est cofinal comme limite inductive de foncteurs cofinaux.

\dgrblock{dgrb1074}{Replacing a presentation by a continuous poset of presentations}
\begin{dcomment}{exposition}{p305}{Replacing a presentation by a continuous poset of presentations}
The adapted presentations are designed to retain more than a filtered colimit expression. They remember colimits of smaller filtered diagrams inside the indexing category, and require the diagram itself to preserve those colimits. This continuity is what will allow a property of $x$ to be detected in any sufficiently well-behaved presentation of $x$.

For the canonical presentation $\mathcal M_\pi/x\to\mathcal M$, the verification is concrete. A diagram with at most $\pi$ arrows has a colimit in $\mathcal M$ when it is large before $\pi_0$. Its colimit carries the induced map to $x$. Moreover it remains $\pi^+$-presentable: its representable functor is an at-most-$\pi$ limit of representables, and those limits of sets commute with $\pi^+$-filtered colimits. Thus the colimit is still an object of $\mathcal M_\pi/x$ and is computed there as claimed.

To replace this indexing category by a poset, one must close small sets of arrows under enough cocone choices. Here is the size argument behind 4.13.13. Start with at most $\pi$ arrows of a small $\pi^+$-filtered category $J$, adjoin identities and composites, and, for every diagram of size at most $\pi_0$ in the resulting subcategory, adjoin a chosen cocone in $J$. There are at most $\pi^{\pi_0}=\pi$ such diagrams. Repeat the operation along $\pi_0^+$ stages, which is permissible because $\pi_0^+\leq\pi$. The union still has at most $\pi$ arrows. Every diagram of size at most $\pi_0$ in the union occurs at one stage, since $\pi_0^+$ is regular, so a subsequent stage contains its cocone. This gives a subcategory large before $\pi_0$.

Applying the construction to the union of at most $\pi$ prescribed subcategories proves that the poset $I$ in 4.13.13 is large before $\pi$. Its smaller directed suprema are literal unions. If $A$ is $\pi_0^+$-directed, $|A|\leq\pi$, and $J_r\subseteq J_s$ for $r\leq s$, then $\bigcup_rJ_r$ has at most $\pi$ arrows. A diagram of size at most $\pi_0$ in this union is contained in one $J_r$, where it has a cocone. The union is therefore again an object of $I$.

This also proves 4.13.13.1 without relying on an unstated stability theorem for final functors. For general transition functors, rather than inclusions, lift the objects and arrows of a small diagram in $\varinjlim_rJ_r$ to a common stage. There are at most $\pi_0$ composition and identity equations to impose, so a further common stage realizes them all. Take a cocone there. Hence the filtered colimit category is again $\pi_0^+$-filtered.

Finally the transition maps between $x_{J'}=\varinjlim_{j\in J'}x_j$ are induced by inclusions of subcategories. For a directed union $J'=\bigcup_rJ_r$, a cocone on $x|_{J'}$ is exactly a compatible family of cocones on the restrictions $x|_{J_r}$. This proves
\[
 \varinjlim_r x_{J_r}\cong x_{\bigcup_rJ_r}
\]
by the universal property, and gives the continuity asserted in 4.13.13.2. Every object and arrow of $J$ is contained in one of these subcategories; the same cocone argument gives $\varinjlim_{J'\in I}x_{J'}\cong\varinjlim_{j\in J}x_j$. This is the precise reason the replacement preserves the object being presented.
\end{dcomment}


\noindent\textbf{Lemme 4.13.13.2.} \textit{Sous les hypothèses de 4.13.13, soient $\mathcal{M}$ une catégorie satisfaisant $\mathrm{L}_{\pi_0}$, et $(x_j)_{j\in J}$ (où $J \longrightarrow \mathcal{M}$) un système inductif dans $\mathcal{M}$, indexé par $J$. Pour toute $J' \in I$, posons}
\[
x_{J'} \ = \ \varinjlim_{j\in J'} x_j
\]
\msp{page 307}\textit{(qui existe dans $\mathcal{M}$, puisque $\mathcal{M}$ satisfait à $\mathrm{L}_{\pi_0}$). Alors les $(x_{J'})_{J'\in I}$ forment un système inductif dans $\mathcal{M}$, et $I \longrightarrow \mathcal{M}$ commute aux limites inductives grandes devant $\pi_0$ et de cardinal $\leqslant \pi$. Si $\mathcal{M}$ est $\pi_0$-accessible, et les $x_j$ ($j \in \Ob J$) sont dans $\mathcal{M}_\pi$, alors $(x_{J'})_{J'\in I}$ est un système inductif $\pi_0$-$\pi$-adapté dans $\mathcal{M}$.}

Tout ceci est essentiellement trivial, et donne :

\noindent\textbf{Corollaire 4.13.14} (de 4.13.12). \textit{Sous les conditions de 4.13.12, si $\pi^{\pi_0} = \pi$, on peut pour tout $x \in \Ob\mathcal{M}$ trouver une représentation $x \simeq \ilim_{i\in I} x_i$, où $(x_i)_{i\in I}$ est un système inductif $\pi_0$-$\pi$-adapté avec $I$} ensemble ordonné. \textit{De plus, si $\pi' > \pi$ et si $x \in \mathcal{M}_{\pi'}$, on peut prendre $I$ tel qu'on ait}
\[
\operatorname{card} I \ \leqslant \ {\pi'}^{\pi} \ .
\]

Ce dernier complément est immédiat,\msp{page 308} via les majorations habituelles.

\dgrblock{dgrb1075}{Adapted systems in a powerset category}
\begin{dcomment}{examples}{p307}{Adapted systems in a powerset category}
The ordered version becomes especially transparent in a powerset category. Let $\mathcal M=\mathcal P(X)$ and let $Y\subseteq X$. For any infinite $\pi$, take
\[
 I=[Y]^{\leq\pi},\qquad D(T)=T.
\]
The poset $I$ is $\pi^+$-directed and its colimit is $Y$. A directed family of at most $\pi$ members has supremum their union, again of cardinal at most $\pi$, and $D$ preserves that supremum. Every value $T$ is $\pi^+$-presentable. Thus this is an adapted presentation, with no cardinal exponent calculation needed in this example.

Fix $Z\subseteq X$ with $|Z|\leq\pi$ and let $\mathcal S$ consist of the subsets containing $Z$. Then
\[
 I_{\mathcal S}=\{T\subseteq Y:|T|\leq\pi,\ Z\subseteq T\}
\]
is cofinal in $I$ exactly when $Z\subseteq Y$. If $Z\subseteq Y$, the extension $T\subseteq T\cup Z$ proves cofinality. Conversely a single object of $I_{\mathcal S}$ implies $Z\subseteq Y$. This calculates the recognition condition of 4.13.15 in full: an eventual supply of small approximations satisfying the property is equivalent to the colimit satisfying it. The word ``eventual'' means cofinality, not that every sufficiently large index must work; in more general diagrams the good indices need not form a final segment.
\end{dcomment}


\noindent\textbf{Théorème 4.13.15} (\footnote[134]{Cf. version plus générale et plus symétrique du théorème 4.13.15 dans 4.13.22.}). \textit{Soient $\mathcal{M}$ une catégorie $\pi_0$-accessible, et $\pi > \pi'_0 \geqslant \pi_0$ tels que $\pi^{\pi_0} = \pi$} (\footnote[135]{On peut se borner à supposer $\pi$ utile pour $\mathcal{M}$ (au lieu de $\pi^{\pi_0} = \pi$), alors on trouve que a, a$'$, b, c sont équivalentes et impliquent c$'$, mais ne sont peut-être pas équivalentes à c$'$.})\textit{. Soit $\uS$ une sous-catégorie strictement pleine de $\mathcal{M}$, et soit $\uS(\pi) = \uS \cap \mathcal{M}_\pi$. Les conditions suivantes sont équivalentes.}
\begin{enumerate}
\item[\circlabel{a}] \textit{$\uS$ est $\pi$-accessible dans $\mathcal{M}$ (4.13.3) et $\uS$ est stable dans $\mathcal{M}$ par limites inductives grandes devant $\pi'_0$.}
\end{enumerate}

\printedpage{116}
% head: A. Grothendieck : Dérivateurs, Ch. XVIII, Édition des Universités de Montpellier II et Paris VII
% foot: Version 17 décembre 2025

\begin{enumerate}
\item[a$'$)] \textit{$\uS$ est $\pi$-accessible dans $\mathcal{M}$, et $\uS(\pi)$ est stable dans $\mathcal{M}_\pi$ par limites inductives grandes devant $\pi'_0$ et de cardinal $\leqslant \pi$} (\footnote[136]{donner un nom à cette dernière condition, soit $\mathrm{L}(\pi'_0, \pi)$ (pour $\uS(\pi)$ dans $\mathcal{M}_\pi$).})\textit{.}
\item[b)] \textit{Pour tout objet $x$ de $\mathcal{M}$, on a ceci : la sous-catégorie strictement pleine $\uS(\pi)/x$ de $J = \mathcal{M}_\pi/x$ est cofinale dans $\mathcal{M}_\pi/x$ si et seulement si $x \in \Ob\uS$. De plus, $\uS(\pi)$ est stable dans $\mathcal{M}_\pi$ pour $\varinjlim$ grandes devant $\pi'_0$ et de cardinal $\leqslant \pi$.}
\item[\circlabel{c}] \textit{Pour tout objet $x$ de $\mathcal{M}$ et pour toute représentation $x \simeq \ilim_{i\in I} x_i$, où $(x_i)_{i\in I}$ est un système inductif $\pi'_0$-$\pi$-adapté, (4.13.11), on a ceci : la sous-catégorie strictement pleine $I_{\uS}$ de $I$ formée des $i \in \Ob I$ tels que $x_i \in \Ob\uS$, est cofinale dans $I$ si et seulement si $x \in \Ob\uS$. De plus, $S(\pi)$ est stable dans $\mathcal{M}_\pi$ par $\varinjlim$ grandes devant $\pi'_0$ et de cardinal $\leqslant \pi$.} % [sic] S(\pi) -> \underline{S}(\pi), printed without the underline used everywhere else for this subcategory
\end{enumerate}
\msp{page 309}
\begin{enumerate}
\item[c$'$)] \textit{Comme d)} \texttt{[plutôt c)]}\textit{, avec $I$ ensemble ordonné.}
\end{enumerate}
\noindent\textbf{Corollaire 4.13.16.} \textit{Les hypothèses sur $\mathcal{M}$, $\pi_0$, $\uS$ étant celles de 4.13.15, pour qu'il existe des cardinaux $\pi > \pi'_0 \geqslant \pi_0$ satisfaisant les conditions $\pi^{\pi'_0} = \pi$ (a fortiori $\pi^{\pi_0} = \pi$) et les conditions équivalentes a) à d$'$)} \texttt{[plutôt c$'$)]} \textit{de 4.13.15, il (faut et il) suffit que $\uS$ soit accessible dans $\mathcal{M}$. De façon plus précise, si $\uS$ est accessible dans $\mathcal{M}$, soit $\pi'_0 \geqslant \pi_0$ un cardinal tel que $\uS$ soit $\pi'_0$-accessible dans $\mathcal{M}$, et soit $\pi$ un cardinal $> \pi'_0$ tel que $\pi^{\pi'_0} = \pi$. Alors $(\uS, \pi'_0, \pi)$ satisfait aux conditions équivalentes de 4.13.15.}

En effet, le “il faut” est évident sur la condition a) de 4.13.15. Pour le “il suffit” on note que $\uS$ $\pi'_0$-accessible dans $\mathcal{M}$ implique que $\uS$ est stable dans $\mathcal{M}$ par $\varinjlim$ grandes devant $\pi'_0$. Donc pour vérifier a), il suffit de vérifier que $\uS$ est $\pi$-accessible dans $\mathcal{M}$. Or comme $\uS$ est\msp{page 310} $\pi'_0$-accessible et que $\pi > \pi'_0$ satisfait $\pi^{\pi'_0} = \pi$, on en conclut que $\uS$ est $\uS_\pi$-accessible, i.e. $\uS \xleftarrow{\ \sim\ } \mathrm{Ind}_\pi(\uS_\pi)$. Donc en vertu de 4.13.3, il reste à voir que $\uS_\pi \subset \mathcal{M}_\pi$, ou ce qui revient au même, que $\uS_\pi \subset \uS(\pi) \overset{\text{déf}}{=} \uS \cap \mathcal{M}_\pi$. (En fait, on aura même égalité, car il est clair que $\uS(\pi) \subset \uS_\pi \ldots$). Or $\uS(\pi)$ est une sous-catégorie de $\uS$ stable par $\varinjlim$ grandes devant $\pi'_0$ et de cardinal $\leqslant \pi$ (puisque $\uS$ et $\mathcal{M}_\pi$ sont stables par ces limites \textit{dans} $\mathcal{M}$, et que ces limites dans $\uS$ et dans $\mathcal{M}$ sont itou), et $\uS(\pi)$ contient $S_{\pi'_0}$. Donc $\uS(\pi)$ contient $\uS_\pi$, en vertu de 4.7.1 (page 161), implication c) $\Longrightarrow$ b$'$), appliquée à $\uS$, $\pi'_0$, $\pi$ au lieu de $\mathcal{M}$, $\pi_0$, $\pi$. Cela prouve donc le corollaire 4.13.16. % [sic] S_{\pi'_0} -> \underline{S}_{\pi'_0}, printed without the underline

\dgrblock{dgrb1076}{The small-colimit condition does not imply the unbounded one}
\begin{dcomment}{development}{p309}{The small-colimit condition does not imply the unbounded one}
Theorem 4.13.15 attempts to recognize stability under all $\pi'_0{}^+$-filtered colimits from stability of the presentable objects under those of size at most $\pi$. This passage requires a relation between $\pi$ and $\pi'_0$, whereas the printed hypothesis controls only $\pi_0$. The distinction has mathematical consequences: the implication a$'$ $\Rightarrow$ a is false under the displayed hypotheses.

\textbf{Counterexample to a$'$ $\Rightarrow$ a.}\ours
Put $\rho=\aleph_1$, $\pi=\beth_{\omega_1}$, $\kappa=\rho^+$, and $\lambda=\pi^+$. The cardinal $\pi$ is a strong limit of cofinality $\rho$. In particular $\pi^{\aleph_0}=\pi$: a countable sequence of ordinals below $\pi$ is bounded, and for every $\theta<\pi$ one has $\theta^{\aleph_0}\leq 2^{\theta\aleph_0}<\pi$. Let
\[
 P=[\pi]^{\leq\rho},\qquad \mathcal M=\mathcal P(P),
\]
with both categories ordered by inclusion. Let $A$ be the full subcategory of $\mathcal M$ consisting of subsets $Y\subseteq P$ with $|Y|\leq\pi$ which are nonempty and $\kappa$-directed in their induced order. Let $\mathcal S$ be the essential image of $\Ind_\pi(A)$ in $\mathcal M$.

The powerset calculation above shows that $\mathcal M$ is $\aleph_0$-accessible in the source convention and that its $\pi$-accessible objects are exactly its subsets of cardinal at most $\pi$. The inclusion of $A$ consequently extends fully faithfully to $\Ind_\pi(A)$: for $Y\in A$, a map $Y\subseteq\bigcup_iZ_i$ into a $\lambda$-directed union factors through one $Z_i$. Thus $\mathcal S$ is $\pi$-accessible in $\mathcal M$. Its $\pi$-accessible objects are exactly $A$, since a $\lambda$-presentable object of an Ind-category is a retract of a generator, and a retract in a poset is the same object.

A $\kappa$-directed union of at most $\pi$ members of $A$ is again in $A$. Its cardinal is at most $\pi$. Any at-most-$\rho$ family of its elements lies in one member of the diagram, where it has an upper bound. Hence $A$ is closed under every $\kappa$-filtered colimit of cardinal at most $\pi$ in $\mathcal M$. This proves condition a$'$ with $\pi_0=\aleph_0$ and $\pi'_0=\rho$.

For each $z\in P$, the principal ideal $\mathord\downarrow z\subseteq P$ already belongs to $A$: it has largest element $z$ and cardinal $2^{|z|}\leq 2^\rho<\pi$, since $\pi$ is a strong limit. But the diagram
\[
 P\longrightarrow\mathcal S,\qquad z\longmapsto\mathord\downarrow z
\]
is $\kappa$-directed and has union $P$, which does not lie in $\mathcal S$. To prove the last assertion, suppose $P$ were a $\lambda$-directed union of members of $A$. The $\pi$-element subset $\{\{\alpha\}:\alpha<\pi\}\subseteq P$ would then lie in a single $Y\in A$. Its $\kappa$-directedness would make $Y$ cofinal in $P$: bound the singletons belonging to any given at-most-$\rho$ subset of $\pi$.

No cofinal subset of $P$ has cardinal at most $\pi$. If $(y_\alpha)_{\alpha<\pi}$ were one, choose increasing cardinals $\pi_\xi<\pi$, $\xi<\rho$, cofinal in $\pi$. For each $\xi$ choose
\[
 a_\xi\in\pi\setminus\bigcup_{\alpha<\pi_\xi}y_\alpha,
\]
which is possible since the union has cardinal less than $\pi$. Then $\{a_\xi:\xi<\rho\}\in P$ is contained in none of the $y_\alpha$, a contradiction. This proves failure of a.\hfill$\square$

A sufficient repair is to strengthen the cardinal equation to $\pi^{\pi'_0}=\pi$, exactly the equation imposed in Corollary 4.13.16. Here is the missing unbounded-colimit argument under that repair. Write $A=\mathcal S\cap\mathcal M_\pi$. The bounded colimits assumed in a$'$ include chains of length $(\pi'_0)^+\leq\pi$. The simultaneous presentation theorem 4.9.3 therefore represents every diagram of cardinal at most $\pi$ in $\Ind_\pi(A)$ by a $\pi^+$-filtered system of diagrams in $A$. For a $\pi'_0{}^+$-filtered such shape, take its colimit in $A$ and then the filtered colimit; interchange of colimits proves that the inclusion $\mathcal S\hookrightarrow\mathcal M$ preserves the result.

For an arbitrarily large $\pi'_0{}^+$-filtered indexing category $K$, the closure construction of 4.13.13, now with $\pi'_0$ in place of $\pi_0$, expresses $K$ as the union of a $\pi^+$-directed system of $\pi'_0{}^+$-filtered subcategories of cardinal at most $\pi$. The equation $\pi^{\pi'_0}=\pi$ is precisely what bounds that construction. Compute the diagram's colimit first on these subcategories and then along their directed system. Both operations preserve membership in $\mathcal S$, proving a.

We checked the source's continuation: p.~337 calls a$'$ $\Rightarrow$ a ``sorital'', and 4.13.19--21, pp.~340--341, states the bounded-to-unbounded extension without resolving this cardinal issue. Our counterexample concerns the printed theorem and those uses of that extension, not the valid existential recognition statement at sufficiently enlarged cardinals. We did not find this particular counterexample in the primary references checked; it is presented as a repair of the argument, without a priority claim.
\end{dcomment}


\textsc{Démonstration de} 4.13.15. Voici les implications $\pm$ tautologiques,
\[
\begin{tikzcd}[row sep=small]
 & & & \mathrm{c}' \\
\mathrm{a} \arrow[r, Rightarrow] & \mathrm{a}' \arrow[r, Rightarrow, "\bullet"] & \mathrm{c} \arrow[ur, Rightarrow] \arrow[dr, Rightarrow] & \\
 & & & \mathrm{b}\ ,
\end{tikzcd}
\tag*{$(*)$}
\]
\msp{page 311}que je vais expliciter un peu. On a a $\Longrightarrow$ a$'$ puisque si a est satisfait, $\uS$ \textit{et} $\mathcal{M}_\pi$ sont stables \textit{dans} $\mathcal{M}$ par $\varinjlim$ grandes devant $\pi'_0$ et de cardinal $\leqslant \pi$, donc aussi leur intersection $\uS(\pi)$,
\printedpage{117}
% head: A. Grothendieck : Dérivateurs, Ch. XVIII, Édition des Universités de Montpellier II et Paris VII
% foot: Version 17 décembre 2025
qui est donc aussi stable dans $\mathcal{M}_\pi$ par lesdites limites. L'implication c $\Longrightarrow$ c$'$ est une tautologie, et c $\Longrightarrow$ b aussi, compte tenu de 4.13.12, qui nous dit que le système inductif canonique $(x_i)_{i\in I}$, avec $I = \mathcal{M}_\pi/x$, de limite $x$ est $\pi_0$-$\pi$-adapté. (\textbf{N.B.} Comme $\pi^{\pi'_0} = \pi$ et $\pi'_0 \geqslant \pi_0$, on a aussi $\pi^{\pi_0} = \pi$, donc $\pi$ est utile pour $\mathcal{M}$, de sorte que 4.13.12 s'applique.)

Reste à prouver la “$\pm$ tautologie” a$'$ $\Longrightarrow$ c. Soit donc une représentation $\pi_0$-$\pi$-adaptée
\[
x \ \simeq \ \varinjlim_{i\in I} x_i \ ,
\]
prouvons que
\[
x \ \in \ \Ob\uS \qquad\quad \Longleftrightarrow \qquad\quad I_{\uS} \text{ cofinal dans } I.
\]
L'implication $\Longleftarrow$ provient du fait que, $\uS$ étant $\pi$-accessible dans $\mathcal{M}$, est stable dans $\mathcal{M}$ par $\varinjlim$ grandes devant $\pi$ (\footnote[137]{On n'utilise ici, pour $\Longleftarrow$, de l'hypothèse $(x_i)_{i\in I}$ $\pi_0$-$\pi$-adapté, que le seul fait que $I$ soit grande devant $\pi$.}). Reste à prouver\msp{page 312} que si $x \in \Ob\uS$, alors $I_{\uS}$ est cofinal dans $I$. Considérons donc le diagramme de catégories
\[
\begin{tikzcd}
I \arrow[r, "\varphi"] & J = \mathcal{M}_\pi/x \\
I' = I_{\uS} \arrow[u, hook] \arrow[r, "\varphi'"] & J' = \uS(\pi)/x \ , \arrow[u, hook, "\alpha"']
\end{tikzcd}
\]
c'est un carré cartésien ($I' = I_{\uS}$ est l'image inverse par $\varphi : I \longrightarrow J$ de la sous-catégorie strictement pleine $J'$ de $J$). On a les propriétés suivantes
\begin{enumerate}
\item[1)] $I$ est grande devant $\pi$, et $\varphi$ cofinal ($\Longrightarrow$ $J$ grande devant $\pi$).
\item[2)] $I$ stable par $\varinjlim$ grandes devant $\pi'_0$ et de cardinal $\leqslant \pi$, et $\varphi$ commute à ces $\varinjlim$ .
\item[3)] $J'$ est stable par $\varinjlim$ du type précédent et $\alpha$ y commute (\footnote[138]{C'est ici qu'intervient l'hypothèse de stabilité faite dans b), car d'après b), $\uS$ est $\pi$-accessible dans $\mathcal{M}$.}).
\item[4)] $J'$ cofinal dans $J$.
\end{enumerate}
\noindent On vient de conclure que $I' = \varphi^{-1}(J')$\msp{page 313} est cofinale dans $I$. Comme c'est une sous-catégorie pleine (et même strictement pleine), il suffit de prouver la propriété F 1, i.e. que pour tout $i \in \Ob I$, existe $i' \in \Ob I'$ et une flèche $i \longrightarrow i'$.

\dgrblock{dgrb1077}{What must be pulled back is a closed cofinal subcategory}
\begin{dcomment}{exposition}{p312}{What must be pulled back is a closed cofinal subcategory}
The difficult implication in the recognition theorem is a$'$ $\Rightarrow$ c. The canonical map from a presentation of $x$ to $\mathcal M_\pi/x$ is cofinal because its values are presentable: maps from at most $\pi$ presentable objects and their compatibility equations factor through a common stage. For $x\in\mathcal S$, the subcategory $\mathcal S(\pi)/x$ is cofinal for the same reason. One must then show that its inverse image in the chosen presentation is cofinal.

A strict inverse image of a cofinal subcategory need not be cofinal without further hypotheses. The inclusions of the even and odd natural numbers into $\mathbb N$ again give a counterexample. Grothendieck's extra continuity conditions are therefore doing real work: after alternating between the presentation and the good approximations, a colimit makes their two chains agree.

The elementary lifting fact used on pp.~321--325 is this. If $F:I\to J$ is final and $I$ is filtered, an arrow $F(i)\to F(k)$ becomes the image of an arrow of $I$ after postcomposition by some $F(k\to k')$. To prove it, put $i$ and $k$ over a common object of $I$. The two resulting arrows out of $F(i)$ become equal after a further transition by the filtered finality criterion. That equality is the required lift. The same criterion equalizes any prescribed small family of parallel arrows after first equalizing them in $J$ and then lifting the resulting transition. This is how the construction corrects commutativity while retaining actual arrows in $I$.

The four compatibility conditions on pp.~316--318 should be read as two factorization identities and naturality for two transformations, one of them shifted by one index. Once the two factorization identities hold at a successor stage, naturality of $u$ with the immediately preceding stage is automatic:
\[
 (u_{\alpha+1}v_\alpha)u_\alpha
   =u_{\alpha+1}(v_\alpha u_\alpha).
\]
Thus the additional equalization proposed on p.~323 is harmless but redundant. Naturality of $v$ with all earlier stages still has to be imposed, especially after a limit stage; this is the substantive correction performed on pp.~324--326. The final block gives the complete construction, including the step taken up on manuscript p.~329.
\end{dcomment}


On va construire encore $i'$ par construction transfinie, sur l'intervalle ordinal $I_{\alpha_\pi}$, en construisant un système inductif
\[
(i_\alpha)_{\alpha\in I_{\alpha_\pi}} \ , \qquad\quad i_\alpha \in \Ob I \ ,
\]
d'objets de $I$, et un système inductif $(j_\alpha)_{\alpha\in I_{\alpha_\pi}}$, $j_\alpha \in \Ob J'$, flèches de transition notées $\tau_{\beta,\alpha} : i_\alpha \longrightarrow i_\beta$, $\tau'_{\beta,\alpha} : j_\alpha \longrightarrow j_\beta$, enfin pour tout $\alpha$ dans $I_{\alpha_\pi}$ des flèches
\[
\begin{tikzcd}
\varphi(i_\alpha) \arrow[r, "u_\alpha"] \arrow[rr, dotted, bend right, "\varphi(\tau_{\alpha+1,\alpha})"'] & j_\alpha \arrow[r, "v_\alpha"] \arrow[rr, dotted, bend left, "\tau'_{\alpha+1,\alpha}"] & \varphi(i_{\alpha+1})\ ( \arrow[r, dotted, "u_{\alpha+1}"] & j_{\alpha+1}\ )
\end{tikzcd}
\tag*{$(*)$}
\]
\printedpage{118}
% head: A. Grothendieck : Dérivateurs, Ch. XVIII, Édition des Universités de Montpellier II et Paris VII
% foot: Version 17 décembre 2025
dans $J$, telles que
\[
\left\{
\begin{aligned}
v_\alpha u_\alpha \ &= \ \varphi(\tau_{\alpha+1,\alpha}) \\
u_{\alpha+1} v_\alpha \ &= \ \tau'_{\alpha+1,\alpha}
\end{aligned}
\right.
\qquad\qquad \text{(\footnotemark[139])} \ .
\tag*{$(**)$}
\]
\footnotetext[139]{cf. aussi $(****)$, page suivante.}

\msp{page 314}\noindent Montrons comment la construction d'un tel système d'objets et de flèches avec
\[
i_0 \ = \ i \tag*{$(***)$}
\]
va impliquer l'existence de $i' \in \Ob I'$ et de $i \longrightarrow i'$. En effet, comme $I_{\alpha_\pi}$ est grande devant $\pi'_0$ (il est grand devant tout cardinal $< \pi$), il existe
\[
i' \ \overset{\text{déf}}{=} \ \varinjlim_{i\in I_{\alpha_\pi}} i_\alpha % [sic] i\in I_{\alpha_\pi} -> \alpha\in I_{\alpha_\pi}, the index of the limit is \alpha
\]
dans $I$, et on a une flèche de transition
\[
i \ = \ i_0 \longrightarrow i' \ ,
\]
donc il reste à prouver que $i' \in \Ob I'$, i.e. $\varphi(i') \in \Ob J'$. Or $\varphi$ commutant aux $\varinjlim$ de type $I_{\alpha_\pi}$, on aura
\[
\varphi(i) \ = \ \varinjlim \varphi(i_\alpha) \ .
\]
Il me faut encore des compatibilités
\[
\begin{tikzcd}
\varphi(i_\alpha) \arrow[r, "u_\alpha"] \arrow[d, "\varphi(\tau_{\beta,\alpha})"'] & j_\alpha \arrow[d, "\tau'_{\beta,\alpha}"] & {} \arrow[d, phantom, "\text{et}" description] & j_\alpha \arrow[r, "v_\alpha"] \arrow[d, "\tau'_{\beta,\alpha}"'] & \varphi(i_{\alpha+1}) \arrow[d, "\varphi(\tau_{\beta+1,\alpha+1})"] \\
\varphi(i_\beta) \arrow[r, "u_\beta"] & j_\alpha & {} & j_\beta \arrow[r, "v_\beta"] & \varphi(i_{\beta+1})
\end{tikzcd}
\tag*{$(****)$} % [sic] j_\alpha (lower right corner of the first square) -> j_\beta, target of \tau'_{\beta,\alpha} and of u_\beta
\]
oubliées tantôt.

\msp{page 315}\noindent Les compatibilités $(****)$ impliquent qu'on a des homomorphismes
\[
\begin{tikzcd}
\varphi(i) \arrow[r, shift left, "u"] & j \arrow[l, shift left, "v"]
\end{tikzcd}
\qquad\qquad \text{où } j \overset{\text{déf}}{=} \varinjlim_{\alpha\in I_{\alpha_\pi}} j_\alpha \quad \text{(limite \textit{dans} $J'$, donc dans $J$)} \ ,
\]
et les compatibilités $(**)$ impliquent que
\[
vu \ = \ \mathrm{id}_{\varphi(i)} \ , \qquad\quad uv = \mathrm{id}_J \ , % [sic] \mathrm{id}_J -> \mathrm{id}_j, uv is an endomorphism of j
\]
donc que $u$, $v$ sont des isomorphismes inverses l'un de l'autre. Comme $J'$ est une sous-catégorie \textit{strictement} pleine de $J$, il s'ensuit bien que $\varphi(i) \in J'$, i.e. $i \in I'$, q.e.d.

Donc tout revient à construire les
\[
\underbrace{i_\alpha}_{\in \Ob I} \ , \ \tau_{\beta,\alpha} \ , \qquad\quad \underbrace{j_\alpha}_{\in \Ob J'} \ , \ \tau'_{\beta,\alpha} \ , \qquad\quad u_\alpha \ , \ v_\alpha \ ,
\]

% ---- from der-ch18_p121-125.tex
\printedpage{119}
% head: A. Grothendieck : Dérivateurs, Ch. XVIII, Édition des Universités de Montpellier II et Paris VII
satisfaisant aux contraintes $(**)$, $(***)$, $(****)$. On va faire la construction en trois étapes.

\circled{1} Départ :
\[
i_0 \;=\; i\ .
\]
Comme $J'$ est cofinal dans $J$, existe $j_0\in\Ob J'$ et
\[
\varphi(i_0)\xrightarrow{u_0} j_0\ .
\]

\msp{page 316}

\circled{2} Supposons que les $i_\beta$, $j_\beta$ soient construits pour $\beta<\alpha$, ainsi que les $\tau_{\beta',\beta}:i_\beta\longrightarrow i_{\beta'}$, $\tau'_{\beta',\beta}:j_\beta\longrightarrow j_{\beta'}$ pour $\beta,\beta'<\alpha$, enfin les $u_\beta:\varphi(i_\beta)\longrightarrow j_\beta$, $v_\beta:j_\beta\longrightarrow\varphi(i_{\beta+1})$ pour $\beta<\alpha$ resp. pour $\beta+1<\alpha$, de façon à satisfaire aux contraintes $(**)$, $(****)$, i.e. commutativités dans
\[
\left\{
\begin{array}{l}
\left.
\begin{array}{ll}
\circled{1} &
\begin{tikzcd}[ampersand replacement=\&]
\varphi(i_\beta) \arrow[r, "u_\beta"] \arrow[rr, bend right=15, "\varphi(\tau_{\beta+1,\beta})"'] \& j_\beta \arrow[r, "v_\beta"] \& \varphi(i_{\beta+1})
\end{tikzcd}
\\[5ex]
\circled{2} &
\begin{tikzcd}[ampersand replacement=\&]
j_\beta \arrow[r, "v_\beta"] \arrow[rr, dotted, bend right=15, "\tau'_{\beta+1,\beta}"'] \& \varphi(i_{\beta+1}) \arrow[r, "u_{\beta+1}"] \& j_{\beta+1}
\end{tikzcd}
\end{array}
\right\}\quad \textit{si}\ \beta+1<\alpha
\\[8ex]
\begin{array}{ll}
\circled{3} &
\begin{tikzcd}[ampersand replacement=\&]
\varphi(i_\beta) \arrow[r, "u_\beta"] \arrow[d, "\varphi(\tau_{\beta',\beta})"'] \& j_\beta \arrow[d, "\tau'_{\beta',\beta}"] \\
\varphi(i_{\beta'}) \arrow[r, "u_{\beta'}"] \& j_{\beta'}
\end{tikzcd}
\qquad\qquad \text{si } \beta<\beta'<\alpha
\end{array}
\\[10ex]
\begin{array}{ll}
\circled{4} &
\begin{tikzcd}[ampersand replacement=\&]
j_\beta \arrow[r, "v_\beta"] \arrow[d, "\tau'_{\beta',\beta}"'] \& \varphi(i_{\beta+1}) \arrow[d, "\varphi(\tau_{\beta'+1,\beta+1})"] \\
j_{\beta'} \arrow[r, "v_{\beta'}"] \& \varphi(i_{\beta'+1})
\end{tikzcd}
\qquad\qquad \text{si } \beta<\beta' \text{ et } \beta'+1<\alpha\ ,
\end{array}
\end{array}
\right.
\]

\msp{page 317}

il faut construire
% foot: Version 17 décembre 2025

\printedpage{120}
% head: A. Grothendieck : Dérivateurs, Ch. XVIII, Édition des Universités de Montpellier II et Paris VII
\begin{itemize}[leftmargin=*]
\item[1)] $i_\alpha\in\Ob I$,\quad et un homomorphisme $(i_\beta)_{\beta\in I_\alpha}\longrightarrow i_\alpha$, i.e. des $\tau_{\alpha,\beta}$ : $i_\beta\longrightarrow i_\alpha$ $(\beta\in I_\alpha)$, compatibles aux translations $\tau_{\beta',\beta}$ pour $\beta<\beta'<\alpha$.
\item[2)] $j_\alpha\in\Ob J'$,\quad et un homomorphisme $(j_\beta)_{\beta\in I_\alpha}\longrightarrow j_\alpha$, i.e. des $\tau'_{\alpha,\beta}$ : $j_\beta\longrightarrow j_\alpha$ $(\beta\in I_\alpha)$, compatibles aux translations $\tau'_{\beta',\beta}$ pour $\beta<\beta'<\alpha$.
\item[3)] Une flèche
\[
u_\alpha:\varphi(i_\alpha)\longrightarrow j_\alpha\ .
\]
\item[4)] Si $\alpha$ n'est pas ordinal limite, i.e. $\alpha=\alpha_0+1$, une flèche
\[
v_{\alpha_0}:j_{\alpha_0}\longrightarrow\varphi(i_{\alpha_0+1})=\varphi(i_\alpha)\ ;
\]
\end{itemize}

et ceci de façon à satisfaire aux contraintes \circled{1}\,\circled{2}\,\circled{3}\,\circled{4} quand on y suppose $\beta+1\leqslant\alpha$ (au lieu de $\beta'<\alpha$) dans \circled{1} et \circled{2}, et $\beta'\leqslant\alpha$ (au lieu de $\beta'<\alpha$) dans \circled{3}, et $\beta'+1\leqslant\alpha$ (au lieu de $\beta'+1<\alpha$) dans \circled{4}. Donc pour les contraintes $1^\circ$, $2^\circ$, $4^\circ$, il n'y a en tenir compte \textit{que} si $\alpha$ n'est pas ordinal limite, donc $\alpha=\alpha_0+1$, et en y faisant % [sic] il n'y a en tenir -> il n'y a à en tenir, the preposition « à » appears to be missing
\[
\beta' \;=\; \alpha_0\ .
\]
Quant à la contrainte $3^\circ$, il faut en
\msp{page 318}
tenir compte pour $\beta'=\alpha$. Donc au total, il faut exiger
\[
(3_\alpha)(\beta)\qquad\qquad
\begin{tikzcd}
\varphi(i_\beta) \arrow[r, "u_\beta"] \arrow[d, "\varphi(\usebox{\dctauAB})"'] & j_\beta \arrow[d, "\tau'_{\alpha,\beta}" {draw, circle, inner sep=1pt}] \\
\varphi(i_\alpha) \arrow[r, "u_\alpha" {draw, circle, inner sep=1pt}] & j_\beta % [sic] j_\beta -> j_\alpha, the target of \tau'_{\alpha,\beta} and u_\alpha is j_\alpha
\end{tikzcd}
\qquad \text{commutatif pour } \beta<\alpha\ (\footnotemark[140]),
\]
\footnotetext[140]{\textbf{N.B.} Il suffit de l'exiger pour $\beta=\alpha_0$, dans le cas $\alpha=\alpha_0+1$, cf. ci-dessous, p. 323.}
(par restriction sur $\alpha$), et si $\alpha=\alpha_0+1$ est ordinal non limite, il faut exiger de plus
\[
(4_\alpha)(\beta)\qquad\qquad
\begin{tikzcd}
j_\beta \arrow[r, "v_\beta"] \arrow[d, "\tau'_{\alpha_0,\beta}"'] & \varphi(i_{\beta+1}) \arrow[d, "\varphi(\usebox{\dctauABone})"] \\
j_{\alpha_0} \arrow[r, "v_{\alpha_0}" {draw, circle, inner sep=1pt}] & \varphi(i_{\alpha_0+1})=\varphi(i_\alpha)
\end{tikzcd}
\qquad \beta<\alpha_0
\]
\[
(1_\alpha)\qquad\qquad
\begin{tikzcd}
\varphi(i_{\alpha_0}) \arrow[r, "u_{\alpha_0}"] \arrow[rr, bend right=15, "\varphi(\usebox{\dctauAAzero})"'] & j_{\alpha_0} \arrow[r, "v_{\alpha_0}" {draw, circle, inner sep=1pt}] & \varphi(i_{\alpha_0+1})=\varphi(i_\alpha)
\end{tikzcd}
\]
% foot: Version 17 décembre 2025

\printedpage{121}
% head: A. Grothendieck : Dérivateurs, Ch. XVIII, Édition des Universités de Montpellier II et Paris VII
\[
(2_\alpha)\qquad\qquad
\begin{tikzcd}
j_{\alpha_0} \arrow[r, "v_{\alpha_0}" {draw, circle, inner sep=1pt}] \arrow[rr, bend right=15, "\tau'_{\alpha,\alpha_0}"' {draw, circle, inner sep=1pt}] & \varphi(i_{\alpha_0+1})=\varphi(i_\alpha) \arrow[r, "u_\alpha" {draw, circle, inner sep=1pt}] & j_\alpha\ .
\end{tikzcd}
\]

Dans ces compatibilités, j'ai cerclé les flèches $\tau_{\alpha,\beta}$, $\tau'_{\alpha,\beta}$ $(\beta<\alpha)$, $\Bigl(\ \tau_{\alpha,\beta+1}\ (\beta+1<\alpha)\ \Bigr)$, $v_{\alpha_0}$, $\Bigl(\ \tau_{\alpha,\alpha_0}\ \Bigr)$, $u_\alpha$, $\Bigl(\ \tau'_{\alpha,\alpha_0}\ \Bigr)$, qui ont été introduites dans les donnés 1) 2) 3) 4), page 317. % [sic] donnés -> données, the noun « données » is meant

\msp{page 319}

On va distinguer donc les deux cas.

\textit{2 A)} $\alpha=\alpha_0+1$.

Pour construire $i_\alpha$ et les $\tau_{\alpha,\beta}:i_\beta\longrightarrow i_\alpha$, il suffit de construire $i_\alpha$ et
\[
\circled{$\tau_{\alpha,\alpha_0}$}\ :\ i_{\alpha_0}\ \longrightarrow\ i_\alpha=i_{\alpha_0+1}\ .
\]
De même pour $j_\alpha$ et les $\tau'_{\alpha,\beta}$, il suffit de construire
\[
\circled{$\tau'_{\alpha,\alpha_0}$}\ :\ j_{\alpha_0}\ \longrightarrow\ j_\alpha\ .
\]
Il faut de plus se donner
\[
j_{\alpha_0}\ \xrightarrow{\circled{$v_{\alpha_0}$}}\ \varphi(\underbrace{\,i_\alpha\,}_{=i_{\alpha_0+1}})\ \xrightarrow{\circled{$u_\alpha$}} j_\alpha\ ,
\]
donc en tout deux objets $i_\alpha$, $j_\alpha$ (dans $I$ et dans $J'$ respectivement) et quatre flèches $\tau_{\alpha,\alpha_0}$, $\tau'_{\alpha,\alpha_0}$, $v_{\alpha_0}$, $u_\alpha$. Il faut quatre compatibilités (page 318), où les deux dernières $3_\alpha(\beta)$, $4_\alpha(\beta)$ dépendent d'un paramètre $\beta<\alpha$, i.e. $\beta\leqslant\alpha_0$, ou de $\beta<\alpha_0$.

\msp{page 320 biffé}

\msp{page 321}

On va d'abord s'occuper de deux premières compatibilités $1_\alpha$, $2_\alpha$, en termes de la donnée de
\[
i_{\alpha_0}\ \in\ \Ob I\ ,\qquad j_{\alpha_0}\ \in\ \Ob J'\ ,\qquad u_{\alpha_0}\ :\ \varphi(i_{\alpha_0})\ \longrightarrow\ j_{\alpha_0}\ .
\]
Les deux compatibilités voulues pour les deux objets $i_\alpha$, $j_\alpha$ et les quatre flèches $\tau_{\alpha,\alpha_0}$, $\tau'_{\alpha,\alpha_0}$, $v_{\alpha_0}$, $u_\alpha$ sont indiquées dans le diagramme
\[
(*)\qquad\qquad
\begin{tikzcd}
\varphi(i_{\alpha_0}) \arrow[r, "u_{\alpha_0}"] \arrow[rr, bend right=15, "\varphi(\tau_{\alpha,\alpha_0})"'] & j_{\alpha_0} \arrow[r, "v_{\alpha_0}\ ?"] \arrow[rr, bend left=20, "\tau'_{\alpha,\alpha_0}"] & \varphi(\underbrace{\,i_\alpha\,}_{?}) \arrow[r, "u_\alpha\ ?"] & j_\alpha\ .
\end{tikzcd}
\]

On note que comme $\varphi$ est cofinal, $\exists\ i_\alpha=i\in\Ob I$ et $j_{\alpha_0}\xrightarrow{v_{\alpha_0}}\varphi(i)$. Mais on veut que le composé
\[
w\ :\ \varphi(i_{\alpha_0})\ \xrightarrow{u_{\alpha_0}}\ j_{\alpha_0}\ \xrightarrow{v_{\alpha_0}}\ \varphi(i)
\]
soit de la forme $\varphi(\tau)$, où $\tau=\tau_{\alpha,\alpha_0}:i_{\alpha_0}\longrightarrow i_\alpha$ est une flèche dans $I$. Comme $\varphi$ est un foncteur cofinal de catégories filtrantes, on sait
% foot: Version 17 décembre 2025
\printedpage{122}
% head: A. Grothendieck : Dérivateurs, Ch. XVIII, Édition des Universités de Montpellier II et Paris VII
\msp{page 322}
qu'on peut trouver $i\xrightarrow{\rho}i'$ dans $I$, tel que la flèche composée $\varphi(i_{\alpha_0})\xrightarrow{w}\varphi(i)\xrightarrow{\varphi(\rho)}\varphi(i')$ soit de la forme $\varphi(\tau)$. Donc quitte à remplacer $i$ par $i'$, et $v_{\alpha_0}$ par le composé $v'_{\alpha_0}$ : $j_{\alpha_0}\xrightarrow{v_{\alpha_0}}\varphi(i)\xrightarrow{\varphi(\rho)}\varphi(i')$, on peut supposer que $v_{\alpha_0}u_{\alpha_0}$ est de la forme $\varphi(\tau)$. Cela nous donne donc en principe l'objet $i=i_\alpha$, et les flèches $v_{\alpha_0}:j_{\alpha_0}\longrightarrow\varphi(i)$ et $\tau_{\alpha,\alpha_0}:i_{\alpha_0}\longrightarrow i_\alpha$, de façon à satisfaire à la première compatibilité dans $(*)$. Il reste à construire $j_\alpha$, $\tau'_{\alpha,\alpha_0}$, $u_\alpha$ de façon à satisfaire à la deuxième compatibilité $(*)$. Or $J'$ étant cofinal dans $J$, il existe $j_\alpha$ dans $J'$ et une flèche $u_\alpha:\varphi(i_\alpha)\longrightarrow j_\alpha$. D'autre part, comme $J'\longrightarrow J$ est pleinement fidèle, la flèche composée $j_{\alpha_0}\xrightarrow{v_{\alpha_0}}\varphi(i_\alpha)\xrightarrow{u_\alpha}j_\alpha$ dans $J$, entre objets de $J'$, provient d'une unique flèche $\tau'_{\alpha,\alpha_0}:j_{\alpha_0}\longrightarrow j_\alpha$ dans $J'$.

\msp{page 323}

Ainsi, on a construit les deux objets et les quatre flèches voulues, de façon à satisfaire aux contraintes $1_\alpha$, $2_\alpha$ de la page 318. Les contraintes $3_\alpha(\beta)$, $4_\alpha(\beta)$ sont-elles alors satisfaites automatiquement ?
\[
\begin{tikzcd}
\varphi(i_\beta) \arrow[r, "u_\beta"] \arrow[d] & j_\beta \arrow[d] \\
\varphi(i_{\alpha_0}) \arrow[r, "u_{\alpha_0}"] \arrow[d, "\varphi(\tau_{\alpha,\alpha_0})"'] & j_{\alpha_0} \arrow[d, "\tau'_{\alpha,\alpha_0}"] \\
\varphi(i_\alpha) \arrow[r, "u_\alpha"] & j_\alpha
\end{tikzcd}
\qquad\qquad \beta<\alpha, \text{ i.e. } \beta\leqslant\alpha_0\ .
\]
Le carré supérieur est commutatif par hypothèse de récurrence, pour que le carré total (composé) soit commutatif, il suffit que le carré inférieur le soit. On a donc une double-flèche $(\tau'_{\alpha,\alpha_0}u_{\alpha_0},u_\alpha\varphi(\tau_{\alpha,\alpha_0})):\varphi(i_\alpha)\rightrightarrows j_\alpha$, elle n'est peut-être pas triviale, mais on sait, $J'$ étant filtrante, qu'on peut l'égaliser par une flèche $j_\alpha\longrightarrow\bar{j}_\alpha$ dans $J'$. Donc quitte à remplacer $j_\alpha$ par % [sic] \varphi(i_\alpha) -> \varphi(i_{\alpha_0}), both composites of the double arrow start at \varphi(i_{\alpha_0})
\msp{page 324}
$\bar{j}_\alpha$, $u_\alpha$ par le composé $\bar{u}_\alpha:\varphi(i_0)\longrightarrow j_\alpha\longrightarrow\bar{j}_\alpha$ et $\tau'_{\alpha,\alpha_0}$ par le composé $\bar{\tau}'_{\alpha,\alpha_0}:j_{\alpha_0}\longrightarrow j_\alpha\longrightarrow\bar{j}_\alpha$ (sans toucher à $i_\alpha$, $v_{\alpha_0}$, $\tau_{\alpha_0,\alpha}$), on peut obtenir la commutativité, d'où $3_\alpha(\beta)$. % [sic] \varphi(i_0) -> \varphi(i_\alpha), u_\alpha has source \varphi(i_\alpha); [sic] \tau_{\alpha_0,\alpha} -> \tau_{\alpha,\alpha_0}, the arrow is written \tau_{\alpha,\alpha_0} everywhere else

Reste la contrainte $4_\alpha(\beta)$. Pour ceci, considérons le diagramme composite
\[
\begin{tikzcd}
j_\beta \arrow[r, "v_\beta"] \arrow[d, "\tau'_{\alpha_0,\beta}"'] & \varphi(i_{\beta+1}) \arrow[d, "\varphi(\tau_{\alpha,\beta+1})"] \\
j_{\alpha_0} \arrow[r, "v_{\alpha_0}"] & \varphi(i_{\alpha_0+1})=\varphi(i_\alpha)\ .
\end{tikzcd}
\]

Si $\alpha_0$ n'est pas un ordinal limite, mais de la forme $\alpha_1+1$, on peut encore se ramener au cas où $\beta=\alpha_1$, en procédant comme ci-dessus. Mais on peut procéder dans tous les cas, sans un tel artifice (\,!), en considérant les double-flèches
\[
(*)\qquad\qquad
\begin{tikzcd}
j_\beta \arrow[r, shift left, "\varphi(\tau_{\alpha,\beta+1})v_\beta"] \arrow[r, shift right, "v_{\alpha_0}\tau'_{\alpha_0,\beta}"'] & \varphi(i_\alpha)
\end{tikzcd}
\]
% foot: Version 17 décembre 2025
\printedpage{123}
% head: A. Grothendieck : Dérivateurs, Ch. XVIII, Édition des Universités de Montpellier II et Paris VII
déduites du carré précédent. Comme
\msp{page 325}
l'ensemble de ces double-flèches est de cardinal $\card\alpha<\pi$, et que $J$ est grande devant $\pi$ (il suffirait qu'il le soit pour les cardinaux $\pi'<\pi$), on peut trouver une flèche $\varphi(i_\alpha)\xrightarrow{u}j$ dans $J$ qui égalise toutes ces double-flèches. Comme $\varphi$ est cofinal, on peut supposer que $j$ est de la forme $\varphi(i)$. Comme $\varphi$ est un foncteur cofinal de catégories filtrantes, il existe alors $i\xrightarrow{\rho}i'$ dans $I$, tel que le composé $\varphi(i_\alpha)\xrightarrow{u}\varphi(i)\xrightarrow{\varphi(\rho)}\varphi(i')$ soit de la forme $\varphi(\sigma)$, pour $\sigma:i_\alpha\longrightarrow i'$. Au total, on a trouvé $\sigma:i_\alpha\longrightarrow i'$, tel que $\varphi(\sigma)$ $(=\varphi(\rho)u):\varphi(i_\alpha)\longrightarrow\varphi(i')$ égalise toutes les double-flèches envisagées dans $(*)$, p. 324. Donc il s'impose de remplacer $i_\alpha$ par $\bar{i}_\alpha=i'$, $v_{\alpha_0}$ par le composé $j_{\alpha_0}\longrightarrow\varphi(i_\alpha)\xrightarrow{\varphi(\sigma)}\varphi(\bar{i}_\alpha)$, $\tau_{\alpha,\alpha_0}$
\msp{page 326}
par le composé $\bar{\tau}_{\alpha,\alpha_0}:i_{\alpha_0}\xrightarrow{\tau_{\alpha,\alpha_0}}i_\alpha\xrightarrow{\sigma}\bar{i}_\alpha$.

Pour résumer, voici la recette pour construire, d'abord \ovalbox{$i_\alpha$, $v_{\alpha_0}$, $\tau_{\alpha,\alpha_0}$,}et \textit{ensuite} \ovalbox{$j_\alpha$, $u_\alpha$, $\tau'_{\alpha,\alpha_0}$}.

On procède d'abord comme page 321, 322 pour construire le premier triple, de façon à satisfaire à la première compatibilité $(*)$ (p. 321), i.e. à $1_\alpha$. On “corrige” ensuite ce choix de $i_\alpha$, $v_{\alpha_0}$, $\tau_{\alpha,\alpha_0}$, en remplaçant $i_\alpha$ par $\bar{i}_\alpha$ pour $\sigma:i_\alpha\longrightarrow\bar{i}_\alpha$, $v_{\alpha_0}$ par $\bar{v}_{\alpha_0}=\varphi(\sigma)v_{\alpha_0}$, et $\tau_{\alpha,\alpha_0}$ par $\bar{\tau}_{\alpha,\alpha_0}=\sigma\tau_{\alpha,\alpha_0}$, de sorte à obtenir les compatibilités $4_\alpha(\beta)$ pour $\beta<\alpha_0$, comme on vient de l'indiquer. On ne touchera plus à ces choix de $i_\alpha$, $v_{\alpha_0}$, $\tau_{\alpha,\alpha_0}$, dans la suite de la construction. On construit $(j_\alpha,u_\alpha,\tau'_{\alpha,\alpha_0})$ de façon à satisfaire à la deuxième compatibilité $(*)$ (p. 321), comme
\msp{page 327}
indiqué page 322. Ensuite, on corrige ce choix, en remplaçant $j_\alpha$ par $\bar{j}_\alpha$, avec $\rho:j_\alpha\longrightarrow\bar{j}_\alpha$, et $u_\alpha$ par $\bar{u}_\alpha=\rho u_\alpha$, $\tau'_{\alpha,\alpha_0}$ par $\bar{\tau}'_{\alpha,\alpha_0}=\rho\tau'_{\alpha,\alpha_0}$, de façon à satisfaire aux contraintes $3_\alpha(\beta)$ de la page 318.

Cela achève la construction dans le cas où $\alpha=\alpha_0+1$.

\textit{Cas 2 B : $\alpha$ est ordinal limite.}

Il faut donc construire les données 1) 2) 3) de la page 317 (à l'exclusion de la donnée 4), de façon à satisfaire à la seule contrainte $3_\alpha(\beta)$.

On commence à construire $i_\alpha$ et les $\tau_{\beta,\alpha}$ dans $I$. C'est possible, car $I_\alpha=\{\beta\mid\beta<\alpha\}$ est de cardinal $\pi_\alpha<\pi$ – il suffit donc de supposer $I$ grande devant les cardinaux $\pi'<\pi$.

On construit ensuite de même $J_\alpha$ dans $J'$ et les $\tau'_{\beta,\alpha}$ $(\beta<\alpha)$, en utilisant le % [sic] J_\alpha -> j_\alpha, the object of J' is meant
\msp{page 328}
fait que $J'$ est grande devant $\card\alpha=\card I_\alpha$. On doit de plus construire une flèche
\[
u_\alpha\ :\ \varphi(i_\alpha)\ \longrightarrow\ j_\alpha\ .
\]
C'est possible, car $J'$ étant cofinale dans $J$, il existe $j'$ dans $J'$ et $\varphi(i_\alpha)\longrightarrow J'$, et $J'$ étant filtrante, il existe $\bar{j}_\alpha$ et des flèches % [sic] \varphi(i_\alpha)\longrightarrow J' -> \varphi(i_\alpha)\longrightarrow j', the arrow goes to the object j'
$\begin{tikzcd}[ampersand replacement=\&, baseline=(\tikzcdmatrixname-1-1.base)]
j_\alpha \arrow[rr] \& \& j_\alpha \\
\& j' \arrow[ur] \&
\end{tikzcd}$, donc quitte à remplacer $j_\alpha$ par $\bar{j}_\alpha$, les $\tau'_{\alpha,\beta}$ par des $\bar{\tau}'_{\alpha,\beta}$, on trouve bien une flèche $u_\alpha:\varphi(i_\alpha)\longrightarrow j_\alpha$. Il faut ensuite s'occuper des % [sic] target j_\alpha in the small diagram -> \bar{j}_\alpha, the two arrows go to the object \bar{j}_\alpha just introduced
% foot: Version 17 décembre 2025


\dgrblock{dgrb1078}{Completing the alternating construction at the page boundary}
\begin{dcomment}{exposition}{p328}{Completing the alternating construction at the page boundary}
Grothendieck is finishing a construction which makes a closed cofinal class visible in every continuous presentation. These pages end on manuscript p.~328, before the last naturality equations have been equalized. The complete chapter continues on manuscript p.~329 and states the result as Lemma 4.13.17 on pp.~329--330. We give a proof here with a suitable regular ordinal in place of the possibly singular ordinal $\pi$. This is a completion of his proof, not an additional assertion of originality.

Let $F:I\to J$ be final, with $I$ $\pi^+$-filtered, and let $J'\subseteq J$ be replete, full, and final. Let $\rho<\pi$ be infinite. Assume that $I$ and $J'$ admit $\kappa$-indexed chain colimits for $\kappa=\rho^+$, and that both $F$ and the inclusion $J'\to J$ preserve them. Then $I'=F^{-1}(J')$ is final in $I$. These assumptions hold in the application on p.~312, taking $\rho=\pi'_0$; no equation $\pi^{\pi'_0}=\pi$ is needed for this implication.

Fix $i_0\in I$. We construct chains $i_\alpha\in I$, $j_\alpha\in J'$ for $\alpha<\kappa$, with arrows
\[
 F(i_\alpha)\xrightarrow{u_\alpha}j_\alpha
 \xrightarrow{v_\alpha}F(i_{\alpha+1})
\]
and transition maps $t_{\beta\alpha}$, $s_{\beta\alpha}$. Require
\[
 v_\alpha u_\alpha=F(t_{\alpha+1,\alpha}),\qquad
 u_{\alpha+1}v_\alpha=s_{\alpha+1,\alpha},
\]
and naturality of $u_\alpha:F(i_\alpha)\to j_\alpha$ and $v_\alpha:j_\alpha\to F(i_{\alpha+1})$. Choose $j_0,u_0$ using finality of $J'\to J$.

At a successor stage, finality of $F$ gives $j_\alpha\to F(k)$. The lifting fact explained above allows us to enlarge $k$ and arrange that its composite with $u_\alpha$ is $F(t)$ for an actual $t:i_\alpha\to k$. For each $\beta<\alpha$, compare the two arrows
\[
 j_\beta\rightrightarrows F(k):\quad
 F(t\,t_{\alpha,\beta+1})v_\beta,\qquad
 v_\alpha s_{\alpha\beta}.
\]
There are at most $\pi$ of them. The category $J$ is $\pi^+$-filtered, as follows by applying finality to diagrams in $I$ or the criterion 4.9.6.3. Equalize this family in $J$, then use finality and the lifting fact to replace the equalizing transition by $F(k\to\bar k)$. Set $i_{\alpha+1}=\bar k$ and postcompose the already chosen arrows. This imposes naturality of $v$ without disturbing its factorization identity. Choose $u_{\alpha+1}:F(i_{\alpha+1})\to j_{\alpha+1}$ with $j_{\alpha+1}\in J'$, and define $s_{\alpha+1,\alpha}=u_{\alpha+1}v_\alpha$. Naturality of $u$ follows from the first identity and earlier naturality; the remaining transition maps are composites.

At a nonzero limit $\gamma<\kappa$, choose compatible cocones $i_\beta\to i_\gamma$ and $j_\beta\to j_\gamma$ on the two preceding chains. This is possible by filteredness, since $|\gamma|\leq\rho<\pi$. Extend $j_\gamma$ inside $J'$ if necessary to obtain $u_\gamma:F(i_\gamma)\to j_\gamma$. For each $\beta<\gamma$, the two composites
\[
 F(i_\beta)\rightrightarrows j_\gamma:\quad
 s_{\gamma\beta}u_\beta,\qquad u_\gamma F(t_{\gamma\beta})
\]
need not yet coincide. Equalize them simultaneously in $J$ and then extend the target into $J'$. Postcomposition corrects $u_\gamma$ and the whole $j$-cocone, preserving all its transition identities. This is exactly the operation begun on p.~328 and completed on p.~329. There is no $v$-naturality condition with target index $\gamma$ at this step: its target is the shifted index $\alpha+1$, so the new conditions are imposed at the next successor stage.

Take $i_*=\varinjlim_{\alpha<\kappa}i_\alpha$ and $j_*=\varinjlim_{\alpha<\kappa}j_\alpha$. Preservation and naturality give arrows $F(i_*)\rightleftarrows j_*$. Their composites are identities because they are induced by the transition maps of the two chains. Hence $F(i_*)\cong j_*\in J'$, and repleteness gives $i_*\in I'$. The map $i_0\to i_*$ proves the first finality condition; fullness and filteredness imply full finality as above.

The source's $i$ in the colimit identities on pp.~314--315 must mean this final $i_*$, not the initial object $i_0$. That distinction is essential: otherwise the proof would assert that the starting object already lies in $I'$. With the corrected notation, the construction proves a$'$ $\Rightarrow$ c. The continuation of the complete chapter also derives b and c$'$ $\Rightarrow$ a$'$ using the same lemma. The separate implication a$'$ $\Rightarrow$ a has the cardinal defect described in the preceding development block; completing the alternating construction does not repair that different step.
\end{dcomment}
\begin{otherlanguage}{english}
\dgroverview{XVIII-c}{Chapter XVIII \textperiodcentered\ \textup{XVIII-c} \textperiodcentered\ Overview}
% Commentary and audit by GPT-6 Astra (OpenAI, ChatGPT Pro mode).
\begin{otherlanguage*}{english}
\noindent\textbf{Presentations which retain structure.}
These pages turn the existence of small approximations into a method for handling objects with structure. An accessible category supplies filtered presentations of its objects; a diagram requires compatible presentations of its arrows and their equations. A section of a fibred category adds transport between different fibres. A full subcategory defined by a property requires approximations which already possess that property. The successive arguments address these problems with the same operations: choose a common stage, impose a bounded family of equations, and take a suitably long colimit. The cardinal estimates specify when this process can still be carried out among the chosen small objects.

These pages begin with Theorem 4.9.3 and end during the proof of Theorem 4.13.15, on manuscript page 328 \citep{Grothendieck1990XVIII}. They are printed pages 79--123 of the edition; the manuscript labels run from an opening continuation before 220 through 328, including 302 bis and the cancelled page 320. These systems should not be identified with one another. The complete chapter was consulted where it supplies the opening hypotheses or the continuation of an unfinished proof. The end of these pages is not itself a mathematical gap.

\medskip\noindent\textbf{The cardinal convention.}
Grothendieck's phrase ``grande devant $\pi$'' means modern $\pi^+$-filteredness: every diagram with at most $\pi$ arrows has a cocone. His $\pi$-accessible objects are therefore $\pi^+$-presentable objects, and his $\Ind_\pi C$ is the completion under $\pi^+$-filtered colimits. This translation matters even for $\pi=\aleph_0$, where countable constraints, rather than only finite constraints, are admitted. It also prevents an unnoticed assumption that $\pi$ is regular. The adjective ``regular'' applied to a category's filtration elsewhere in the chapter denotes a separate categorical condition.

A second distinction concerns a full subcategory $\mathcal S\subseteq\mathcal M$. Intrinsic accessibility, accessibility of the inclusion, and preservation of presentable objects by that inclusion are different requirements. The source's fixed-rank notion of accessibility \emph{in} $\mathcal M$ requires the last comparison as well. We keep the intrinsic notation $\mathcal S_\pi$ separate from $\mathcal S\cap\mathcal M_\pi$ until the hypotheses identify them.

\medskip\noindent\textbf{Cofinality and simultaneous presentation.}
The first task is to show that a small diagram in $\Ind_\pi C$ comes from a filtered system of diagrams in $C$. Full faithfulness follows from expressing natural transformations as an equalizer of products of hom sets, provided the indexing category has at most $\pi$ arrows. Essential surjectivity is the substantial part. The source supplies it when $C$ has colimits of one suitable limit-ordinal shape. In the ensuing section theorem, that colimit turns approximately compatible structure maps into actual structure maps.

The cofinality discussion 4.9.6 isolates the mechanism. We use ``final'' for right-cofinal functors. A faithful functor can reflect an equality once the requisite arrows have been lifted after enlargement. Without faithfulness, eventual equalization of arrows having the same image suffices; the commentary proves this relative variant explicitly. It also completes the cardinal descent and the discrete-opfibration lemma. Common upper bounds, connectedness, and equalization are kept distinct, with worked examples showing why none can silently replace another.

One clause is missing from the printed Lemma 4.9.6.3(b$'$): filteredness of the comma categories. The faithful functor from a two-object discrete category to the terminal category disproves the formulation as printed. The proof on page 236 specifies the missing condition. The possible empty-category issue in b$''$ is instead a convention to clarify: its cardinal version already includes the empty family among the families to be bounded.

The finite-diagram question now has a definitive answer. The comparison works for every $C$ exactly for essentially small enough well-founded shapes, in the sense of \citet[Theorem 1.3]{Henry2025Ind}. The finite rigid case in 4.11.6 is included. A nontrivial finite group supplies an obstruction, and the commentary calculates the signed-ordinal involution example at every infinite completion bound. This is a known construction, with its antecedent in SGA~4, I, 8.8.8, not a result claimed for the commentators \citep{SGA4I1972}.

\medskip\noindent\textbf{Sections and the transfer of colimits.}
In 4.10 ``cofibrant'' describes a Grothendieck opfibration. Its sections are families $X_i$ with maps $a_*X_i\to X_j$ for $a:i\to j$; those maps need not be cocartesian. The approximation category of sections over a given section maps finally to the product of the fibrewise approximation categories. Evaluating its canonical colimit therefore recovers the original section. We write out this passage, which the source asks to have made explicit, and prove full faithfulness separately.

The transfinite construction has a coherence defect as displayed. Compatibility of consecutive stages does not compare choices across an internal limit ordinal. A constant two-point diagram whose involution changes at $\omega$ satisfies the displayed conditions but gives no induced map on its constant colimit. The repair imposes compatibility with every earlier stage when constructing a new stage; there are still at most $\pi$ equations. The completed induction proves the intended section theorem. A separate endpoint adjustment handles the parenthesis allowing an auxiliary cardinal equal to $\pi$.

The application to paratoposes concerns locally presentable categories, without an additional assertion of topos exactness. Ordinary sections of an opfibration have pointwise limits whenever the fibres have them: their structure maps are induced by the limit projections. Preservation by transport is needed for pointwise colimits. This distinction vindicates the printed limit argument in 4.10.9 and prevents the editors' tentative alternative from becoming an unwarranted correction.

Question 4.10.11 about ordinary and cartesian sections is settled by the later theory of accessible limits. Products, inserters, and equifiers encode arbitrary sections; requiring invertible structure maps gives the pseudolimit describing cartesian sections. After choosing common ranks, \citet[Theorems 8.1, 8.3 and Corollary 9.2]{Positselski2024Limits} apply. Their fixed-rank conclusions retain hypotheses about a smaller chain-colimit cardinal, close to the distinction Grothendieck is already exploiting.

Section 4.11 transfers bounded colimits from $C$ to $\Ind_\pi C$ by first presenting the entire diagram. The universal-property argument interchanges two limits of hom sets. The separate assertion that $C\to\Ind_\pi C$ preserves existing bounded colimits uses filtered-colimit/limit commutation in sets. Idempotent completeness enters in the converse, through the description of presentable ind-objects as retracts of objects of $C$. We also give the finite rigid induction and the reduction of finite limits and colimits to products and equalizers, or coproducts and coequalizers.

\medskip\noindent\textbf{Accessible functors and accessible properties.}
Uniformization in 4.12 bounds the ranks of the images of a small generating family. At a sufficiently large fixed point $\pi^{\pi_0}=\pi$, an accessible functor carries $\pi$-accessible objects to $\pi$-accessible objects. Fixing a preservation rank reduces the functor category to one with small domain. The proof of 4.12.4 invokes the ind universal property without explicitly securing the required colimits in the target. Choosing the rank useful for both categories repairs the proof, as in 4.12.5. We distinguish this verified form from the broader first assertion rather than manufacture a counterexample from an omitted hypothesis.

There is no common rank for all accessible endofunctors of sets: the representables $X\mapsto X^B$ prove this as $|B|$ grows. Small colimits of accessible set-valued functors are pointwise, so the same example rules out a small accessible generating family for their entire category. This supplies the counterexample sought in Remark 4.12.6.

The existential theorems about accessible full subcategories and their small intersections remain valid. A later higher-categorical counterpart is \citet[Proposition 5.4.7.10]{Lurie2009HTT}. The effective threshold in 4.13.2(c), however, omits the ambient ranks of the generators. The singleton subcategory $\{X\}\subseteq\mathcal P(X)$ gives a complete counterexample to that estimate. The repair retains the missing bound already used on page 282.

For intersections, an ordinal of cardinality $\pi$ need not have cofinality greater than the smaller bound. We prove a replacement construction which attends to all the subcategories at each stage of a chain of regular length $\rho^+\leq\pi$. Its colimit belongs to every subcategory. This proves the main cofinal-intersection lemma without an unnecessary regularity restriction on $\pi$. A related argument repairs the cofinal-pullback construction in the final proof.

Theorem 4.13.15 needs a more substantial correction. Under its printed hypothesis $\pi^{\pi_0}=\pi$, condition a$'$ need not imply a): closure of the presentable part under $\pi_0'^+$-filtered colimits of size at most $\pi$ need not give unrestricted closure. We construct a counterexample with $\pi_0=\aleph_0$, $\pi_0'=\aleph_1$, and $\pi=\beth_{\omega_1}$, proving both the bounded closure and the failure of unrestricted closure. Strengthening the hypothesis to $\pi^{\pi_0'}=\pi$ gives a sufficient repair: the indexing categories can then be exhausted by suitably filtered subcategories of size at most $\pi$. The bounded membership criteria and the accessible-intersection theorem survive this failed implication.

\end{otherlanguage*}
\end{otherlanguage}
\dgrchunkend
